\documentclass{article}
\RequirePackage{etex}
\makeatletter
\def\input@path{{./../}}
\makeatother



\usepackage{titletoc}
% Custom formatting for appendix TOC entries
\titlecontents{section}
  [0pt]
  {\addvspace{1ex}}
  {\contentspush{\thecontentslabel\enspace}}
  {}
  {\titlerule*[0.7em]{.}\contentspage}
  



%\usepackage[toc,page,header]{appendix}
%\usepackage{minitoc}

% if you need to pass options to natbib, use, e.g.:
\PassOptionsToPackage{numbers, compress}{natbib}
% before loading neurips_2023
% ready for submission
%\usepackage{neurips_2023}
%\usepackage[nonatlib]{neurips_2023}

\usepackage{subcaption}
\usepackage{booktabs}
\usepackage{lscape}
% to compile a preprint version, e.g., for submission to arXiv, add add the
% [preprint] option:
%\usepackage[preprint]{neurips_2023}


% to compile a camera-ready version, add the [final] option, e.g.:
     \usepackage[final]{neurips_2023}


% to avoid loading the natbib package, add option nonatbib:
%    \usepackage[nonatbib]{neurips_2023}


\usepackage[utf8]{inputenc} % allow utf-8 input
\usepackage[T1]{fontenc}    % use 8-bit T1 fonts
\usepackage[hidelinks]{hyperref}       % hyperlinks
\usepackage{url}            % simple URL typesetting
\usepackage{booktabs}       % professional-quality tables
\usepackage{amsfonts}       % blackboard math symbols
\usepackage{amsmath}
\usepackage{amsthm}
\newtheorem{theorem}{Theorem}
\usepackage{nicefrac}       % compact symbols for 1/2, etc.
\usepackage{microtype}      % microtypography
\usepackage{xcolor}         % colors
\usepackage{graphicx}
\usepackage{mathtools}


%\usepackage{algorithm}
%\usepackage[noend]{algpseudocode}
%\algnewcommand{\LineComment}[1]{\State \(\triangleright\) #1}
%\usepackage[ruled,vlined]{algorithm2e}


%%%%%%%%%%%%%%%%%%%%%%%%%%%%%%%%
% Luhuan's custom comments stuff
%%%%%%%%%%%%%%%%%%%%%%%%%%%%%%%%
\newif\ifcomments
% \commentsfalse
\commentstrue
\ifcomments\newcommand{\comments}[1]{#1}\else\newcommand{\comments}[1]{}\fi

\definecolor{clrgp}{rgb}{.9,0,.9}
\definecolor{red}{rgb}{.8,0,0}
\definecolor{blue}{rgb}{0,0, 0.8}

%\newcommand{\lw}[1]{}
%\newcommand{\bt}[1]{}
%\newcommand{\can}[1]{}
%\newcommand{\BLT}[1]{}
\newcommand{\lw}[1]{\comments{\textcolor{blue}{[LW: #1]}}}
\newcommand{\bt}[1]{\comments{\textcolor{clrgp}{[BT: #1]}}}
\newcommand{\can}[1]{\comments{\textcolor{red}{[CAN: #1]}}}
\newcommand{\BLT}[1]{\textcolor{blue}{ {\small \ (brian)} #1}}
% !TEX root = main.tex



\usepackage{centernot}
\usepackage{amsthm}
\usepackage{amsmath}
\usepackage{amssymb}
\usepackage{dsfont}
\usepackage{appendix}




\DeclareRobustCommand{\mb}[1]{\ensuremath{\mathbf{\boldsymbol{#1}}}}
% \DeclareRobustCommand{\mb}[1]{\mathbold{#1}}

\DeclareRobustCommand{\KL}[2]{\ensuremath{\textrm{KL}\left[#1\;\|\;#2\right]}}

\DeclareMathOperator*{\argmax}{arg\,max}
\DeclareMathOperator*{\argmin}{arg\,min}




\newcommand\independent{\protect\mathpalette{\protect\independenT}{\perp}}
\def\independenT#1#2{\mathrel{\rlap{$#1#2$}\mkern2mu{#1#2}}}

\renewcommand{\mid}{~\vert~}
\newcommand{\prm}{\:;\:}

\newcommand{\mbw}{\mb{w}}
\newcommand{\mbW}{\mb{W}}

\newcommand{\mbx}{\mb{x}}
\newcommand{\mbX}{\mb{X}}

\newcommand{\mby}{\mb{y}}
\newcommand{\mbY}{\mb{Y}}

\newcommand{\mbz}{\mb{z}}
\newcommand{\mbZ}{\mb{Z}}
\newcommand{\mbT}{\mb{T}}
\newcommand{\mbA}{\mb{A}}
\newcommand{\mba}{\mb{a}}

\newcommand{\mbI}{\mb{I}}
\newcommand{\mbone}{\mb{1}}

\newcommand{\mbL}{\mb{L}}

\newcommand{\mbtheta}{\mb{\theta}}
\newcommand{\mbTheta}{\mb{\Theta}}
\newcommand{\mbomega}{\mb{\omega}}
\newcommand{\mbOmega}{\mb{\Omega}}
\newcommand{\mbsigma}{\mb{\sigma}}
\newcommand{\mbSigma}{\mb{\Sigma}}
\newcommand{\mbphi}{\mb{\phi}}
\newcommand{\mbPhi}{\mb{\Phi}}

\newcommand{\mbalpha}{\mb{\alpha}}
\newcommand{\mbbeta}{\mb{\beta}}
\newcommand{\mbgamma}{\mb{\gamma}}
\newcommand{\mbeta}{\mb{\eta}}
\newcommand{\mbmu}{\mb{\mu}}
\newcommand{\mbrho}{\mb{\rho}}
\newcommand{\mblambda}{\mb{\lambda}}
\newcommand{\mbzeta}{\mb{\zeta}}

\newcommand\dif{\mathop{}\!\mathrm{d}}
\newcommand{\diag}{\textrm{diag}}
\newcommand{\supp}{\textrm{supp}}

\newcommand{\V}{\mathbb{V}}
\newcommand{\bbH}{\mathbb{H}}

\newcommand{\bbN}{\mathbb{N}}
\newcommand{\bbZ}{\mathbb{Z}}
\newcommand{\bbR}{\mathbb{R}}
\newcommand{\bbS}{\mathbb{S}}

\newcommand{\cL}{\mathcal{L}}
\newcommand{\cD}{\mathcal{D}}

\newcommand{\cN}{\mathcal{N}}
\newcommand{\cT}{\mathcal{T}}
\newcommand{\Gam}{\textrm{Gam}}
\newcommand{\InvGam}{\textrm{InvGam}}

% \newcommand{\qedsymbol}{\rule{0.7em}{0.7em}}

\newcommand{\g}{\, | \,}
\newcommand{\s}{\, ; \,}

\newcommand{\indpt}{\protect\mathpalette{\protect\independenT}{\perp}}
%\newcommand{\E}[2]{\mathbb{E}_{#1}\left[#2\right]}

\def\checkmark{\tikz\fill[scale=0.4](0,.35) -- (.25,0) -- (1,.7) -- (.25,.15) -- cycle;} 


\usepackage{booktabs,arydshln}
\makeatletter
\def\adl@drawiv#1#2#3{%
        \hskip.5\tabcolsep
        \xleaders#3{#2.5\@tempdimb #1{1}#2.5\@tempdimb}%
                #2\z@ plus1fil minus1fil\relax
        \hskip.5\tabcolsep}
\newcommand{\cdashlinelr}[1]{%
  \noalign{\vskip\aboverulesep
           \global\let\@dashdrawstore\adl@draw
           \global\let\adl@draw\adl@drawiv}
  \cdashline{#1}
  \noalign{\global\let\adl@draw\@dashdrawstore
           \vskip\belowrulesep}}
\makeatother

\newenvironment{proofsk}{%
  \renewcommand{\proofname}{Proof sketch}\proof}{\endproof}



% ===== useful

\newcommand{\inner}[2]{#1 \cdot #2}
\newcommand{\RR}{\mathbb{R}}
\newcommand{\half}{\frac{1}{2}}
\newcommand{\new}{\mathrm{new}}
\newcommand{\const}{\mathrm{const.}}
\newcommand{\dd}{\mathrm{d}}
\newcommand{\ddx}[1]{\frac{\partial}{\partial #1}}
\newcommand{\dfdx}[2]{\textstyle \frac{\partial #1}{\partial #2}}


\newcommand{\bas}[1]{\begin{align*}#1\end{align*}}
\newcommand{\bac}[1]{\begin{equation}\begin{aligned}#1\end{aligned}\end{equation}}
\newcommand{\ba}[1]{\begin{align}#1\end{align}}
\newcommand{\beq}[1]{\begin{equation}#1\end{equation}}
\newcommand{\bsplt}[1]{\begin{split}#1\end{split}}

\newcommand{\norm}[1]{\left\lVert#1\right\rVert}
\newcommand{\ip}[2]{\langle #1,#2 \rangle}
\newcommand{\idx}[1]{\{1,2,\cdots,#1\}}
\newcommand{\pde}[1]{\frac{\partial}{\partial #1}}
\newcommand{\distas}[1]{\mathbin{\overset{#1}{\kern\z@\sim}}}

\newcommand{\one}{\mathbf{1}}
\newcommand{\zero}{\mathbf{0}}
\newcommand{\indep}{\rotatebox[origin=c]{90}{$\models$}}
\newcommand{\abs}[1]{\lvert #1\rvert}

% ===== probability


\def\independenT#1#2{\mathrel{\rlap{$#1#2$}\mkern2mu{#1#2}}}

\newcommand{\EE}[2]{\mathbb{E}_{#1}\left[#2\right]}
\newcommand{\var}[1]{\textrm{Var}\left[#1\right]}
\newcommand{\h}[1]{\mathbb{H}\left( #1 \right)}
\newcommand{\kl}[1]{\mathrm{KL}\left(#1\right)}
\newcommand{\ind}{\mathbb{1}}
\newcommand{\rmp}{\mathrm{p}}
\newcommand{\rmq}{\mathrm{q}}
\newcommand{\prob}{\mathrm{P}}


% ===== distributions

\newcommand{\mult}{\mathrm{Mult}}
\newcommand{\dir}{\mathrm{Dir}}
\newcommand{\discrete}{\mathrm{Discrete}}
\newcommand{\Bern}{\mathrm{Bernoulli}}
\newcommand{\DP}{\mathrm{DP}}
\newcommand{\GP}{\mathrm{GP}}
\newcommand{\Bet}{\mathrm{Beta}}
\newcommand{\pois}{\mathrm{Pois}}

% ===== calligraphy



% ===== bold mathematics


\newcommand{\Amat}[0]{{{\bf A}}}
\newcommand{\Bmat}[0]{{{\bf B}}}
\newcommand{\Cmat}[0]{{{\bf C}} }
\newcommand{\Dmat}{{\bf D}}
\newcommand{\Emat}[0]{{{\bf E}}}
\newcommand{\Fmat}[0]{{{\bf F}} }
\newcommand{\Gmat}[0]{{{\bf G}} }
\newcommand{\Hmat}[0]{{{\bf H}} }
\newcommand{\Imat}{{\bf I}}
\newcommand{\Jmat}[0]{{{\bf J}} }
\newcommand{\Kmat}[0]{{{\bf K}} }
\newcommand{\Lmat}[0]{{{\bf L}}}
\newcommand{\Mmat}{{\bf M}}
\newcommand{\Nmat}[0]{{{\bf N}}}
\newcommand{\Omat}[0]{{{\bf O}}}
\newcommand{\Pmat}[0]{{{\bf P}}}
\newcommand{\Qmat}[0]{{{\bf Q}} }
\newcommand{\Rmat}[0]{{{\bf R}}}
\newcommand{\Smat}[0]{{{\bf j}}}
\newcommand{\Tmat}[0]{{{\bf T}}}
\newcommand{\Umat}[0]{{{\bf U}}}
\newcommand{\Vmat}[0]{{{\bf V}}}
\newcommand{\Wmat}[0]{{{\bf W}}}
\newcommand{\Xmat}[0]{{{\bf X}}}
\newcommand{\Ymat}[0]{{{\bf Y}}}
\newcommand{\Zmat}{{\bf Z}}


\newcommand{\av}[0]{{\boldsymbol{a}}}
\newcommand{\bv}[0]{{\boldsymbol{b}}}
\newcommand{\cv}[0]{{\boldsymbol{c}}}
\newcommand{\dv}{\boldsymbol{d}}
\newcommand{\ev}[0]{{\boldsymbol{e}} }
\newcommand{\fv}[0]{{\boldsymbol{f}} }
\newcommand{\gv}[0]{{\boldsymbol{g}} }
\newcommand{\hv}[0]{{\boldsymbol{h}} }
\newcommand{\iv}[0]{{\boldsymbol{i}} }
\newcommand{\jv}[0]{{\boldsymbol{j}} }
\newcommand{\kv}[0]{{\boldsymbol{k}} }
\newcommand{\lv}[0]{{\boldsymbol{l}}}
\newcommand{\mv}[0]{{\boldsymbol{m}}}
\newcommand{\nv}[0]{{\boldsymbol{n}}}
\newcommand{\ov}[0]{{\boldsymbol{o}} }
\newcommand{\pv}[0]{{\boldsymbol{p}}}
\newcommand{\qv}[0]{{\boldsymbol{q}} }
\newcommand{\rv}{\boldsymbol{r}}
\newcommand{\sv}{\boldsymbol{s}}
%\newcommand{\sv}[0]{{\boldsymbol{j}}}
\newcommand{\tv}[0]{{\boldsymbol{t}} }
\newcommand{\uv}{\boldsymbol{u}}
\newcommand{\vv}{\boldsymbol{v}}
\newcommand{\wv}{\boldsymbol{w}}
\newcommand{\xv}{\boldsymbol{x}}
\newcommand{\yv}{\boldsymbol{y}}
\newcommand{\zv}{\boldsymbol{z}}
\newcommand{\cdotv}{\boldsymbol{\cdot}}

\newcommand{\Gammamat}[0]{{\boldsymbol{\Gamma}} }
\newcommand{\Deltamat}[0]{{\boldsymbol{\Delta}} }
\newcommand{\Thetamat}{\boldsymbol{\Theta}}
\newcommand{\Betamat}{\boldsymbol{\Beta}}
\newcommand{\Lambdamat}{\boldsymbol{\Lambda}}
\newcommand{\Ximat}[0]{{\boldsymbol{\Xi}} }
\newcommand{\Pimat}[0]{{\boldsymbol{\Pi}} }
\newcommand{\Sigmamat}[0]{{\boldsymbol{\Sigma}}}
\newcommand{\Upsilonmat}[0]{{\boldsymbol{\Upsilon}} }
\newcommand{\Phimat}{\boldsymbol{\Phi}}
\newcommand{\Psimat}{\boldsymbol{\Psi}}
\newcommand{\Omegamat}[0]{{\boldsymbol{\Omega}}}

\newcommand{\alphav}{\boldsymbol{\alpha}}
\newcommand{\betav}[0]{{\boldsymbol{\beta}}}
\newcommand{\gammav}[0]{{\boldsymbol{\gamma}} }
\newcommand{\deltav}[0]{{\boldsymbol{\delta}} }
\newcommand{\epsilonv}{\boldsymbol{\epsilon}}
\newcommand{\zetav}[0]{{\boldsymbol{\zeta}} }
\newcommand{\etav}[0]{{\boldsymbol{\eta}}}
\newcommand{\ellv}[0]{{\boldsymbol{\ell}}}
\newcommand{\thetav}{\boldsymbol{\theta}}
\newcommand{\iotav}[0]{{\boldsymbol{\iota}}}
\newcommand{\kappav}[0]{{\boldsymbol{\kappa}} }
\newcommand{\lambdav}[0]{{\boldsymbol{\lambda}}}
\newcommand{\muv}[0]{{\boldsymbol{\mu}}}
\newcommand{\nuv}[0]{{\boldsymbol{\nu}}}
\newcommand{\xiv}[0]{{\boldsymbol{\xi}} }
\newcommand{\omicronv}[0]{{\boldsymbol{\omicron}} }
\newcommand{\piv}{\boldsymbol{\pi}}
\newcommand{\rhov}[0]{{\boldsymbol{\rho}} }
\newcommand{\sigmav}[0]{{\boldsymbol{\sigma}} }
\newcommand{\tauv}[0]{{\boldsymbol{\tau}} }
\newcommand{\upsilonv}[0]{{\boldsymbol{\upsilon}} }
\newcommand{\phiv}{\boldsymbol{\phi}}
\newcommand{\chiv}[0]{{\boldsymbol{\chi}} }
\newcommand{\psiv}{\boldsymbol{\psi}}
\newcommand{\varthetav}{\boldsymbol{\vartheta}}
\newcommand{\omegav}[0]{{\boldsymbol{\omega}}}

\newcommand{\bR}{\mathbb{R}}
\newcommand{\bE}{\mathbb{E}}
\newcommand{\bN}{\mathbb{N}}
\newcommand{\cE}{\mathcal{E}}
\newcommand{\E}{\mathbb{E}}
\newcommand{\G}{\mathcal{G}}

\newcommand{\grad}{\nabla}



% Brian's macros
\newcommand{\R}{\mathbb{R}}
\newcommand{\Var}{\mathrm{Var}}
\usepackage{cancel}
%\def\Q{\mathbb{Q}}
\def\Q{Q}
\def\P{\mathbb{P}}
\def\inv{^{-1}}
\def\X{\mathcal{X}}
\def\M{\mathcal{M}}
\def\I{\mathbb{I}}
\def\pinv{p_{\mathrm{inv}}}
\def\pref{p_{\mathrm{ref}}}
\def\pmodel{p_{\theta}}
\def\pmodelapprox{\tilde{p}_{\theta}}
%\def\denoiser{\mu_\theta}
%\def\pmodel{p}
%\def\pmodelapprox{\tilde{p}}
%\def\denoiser{\mu}
\def\denoiser{{\hat x}_\theta }
\def\scoremodel{s_\theta}
\newcommand{\dt}[2]{#1^{#2}}
\newcommand{\xt}{\dt{x}{t}}
\newcommand{\xtpo}{\dt{x}{t+1}}
\newcommand{\xprev}{\dt{x}{t-1}}
\newcommand{\xtmo}{\xprev}
\newcommand{\xT}{\dt{x}{T}}
\newcommand{\xzero}{\dt{x}{0}}
%\newcommand{\plikelihood}[2]{f_{#1}(#2)}
\newcommand{\plikelihood}[2]{p_{y|\xzero}({#1}|{#2})}
\newcommand{\proposalVar}{\tilde \sigma^2}
%\newcommand{\plikelihood}[2]{p({#1};{#2})}

\usepackage[ruled,vlined,linesnumbered]{algorithm2e}


\newcommand{\methodname}{Twisted Diffusion Sampler}
\newcommand{\methodshort}{TDS}

\def\psivar{\bar \sigma^2}

\def\truelabel{z^{(0)}}

% SMC
\newcommand{\proposal}{r}
\newcommand{\weight}{w}
\newcommand{\target}{\nu}
\newcommand{\mask}{\mathbf{M}}
\newcommand{\umask}{\mathbf{\bar M}}

% SMC twisted
\newcommand{\twist}{\tilde p_\theta} % please use \pTwist instead
\newcommand{\twistfcn}{\tilde p}
\newcommand{\pTwist}{\tilde p_\theta} % please use this one for consistency
\newcommand{\wTwist}{w}
\newcommand{\tTwist}{\nu}
\newtheorem{definition}{Definition}
\newtheorem{lemma}{Lemma}
\newtheorem{remark}{Remark}
\newtheorem{cor}{Corollary}
\newtheorem{proposition}{Proposition}
\newtheorem{assumption}{Assumption}
\newtheorem{example}{Example}
\usepackage{cleveref}
\newtheorem{fkModel}{SMC Sampler}[section]
\crefname{fkModel}{sampler}{samplers}
\Crefname{fkModel}{Sampler}{Samplers}

\usepackage{enumitem}
\newlist{requirements}{enumerate}{1}
\setlist[requirements,1]{label=\textbf{R\arabic*.},
                   ref  =\textbf{R\arabic*.}}
\crefname{requirementsi}{requirement}{requirements}
\Crefname{requirementsi}{Requirement}{Requirements}

%\title{Twisting Towards Perfection: Asymptotically Accurate Conditional Sampling in Diffusion Models}
\title{Practical and Asymptotically Exact Conditional Sampling in Diffusion Models}


% The \author macro works with any number of authors. There are two commands
% used to separate the names and addresses of multiple authors: \And and \AND.
%
% Using \And between authors leaves it to LaTeX to determine where to break the
% lines. Using \AND forces a line break at that point. So, if LaTeX puts 3 of 4
% authors names on the first line, and the last on the second line, try using
% \AND instead of \And before the third author name.


\author{%
  Luhuan Wu\thanks{Equal contribution, order by coin flip.} \\
  Columbia University\\
  \texttt{lw2827@columbia.edu} \\
  % examples of more authors
  \And
  Brian L. Trippe\textsuperscript{*} \\
  Columbia University\\
  \texttt{blt2114@columbia.edu} \\
  \And
  Christian A. Naesseth \\
  University of Amsterdam \\
  \texttt{c.a.naesseth@uva.nl} \\
  \And
  David M. Blei \\
  Columbia University\\
  \texttt{david.blei@columbia.edu} \\
  \And
  John P. Cunningham \\
  Columbia University\\
  \texttt{jpc2181@columbia.edu} \\
  % Address \\
  % Address \\
  % \texttt{email} \\
  % \AND
  % Coauthor \\
  % Affiliation \\
  % Address \\
  % \texttt{email} \\
  % \And
  % Coauthor \\
  % Affiliation \\
  % Address \\
  % \texttt{email} \\
  % \And
  % Coauthor \\
  % Affiliation \\
  % Address \\
  % \texttt{email} \\
}

% Included in NeurIPS submission for space, deleting for arxiv
%\setlength{\abovedisplayskip}{1pt}
%\setlength{\belowdisplayskip}{1pt}
%\setlength{\abovedisplayshortskip}{1pt}
%\setlength{\belowdisplayshortskip}{1pt}

\usepackage{autonum}

\newcommand{\workshopsubmission}{0} % 1 for true, 0 for false

\makeatletter
\renewcommand\section{\@startsection {section}{1}{\z@}
                                   {-0.5ex \@plus -1ex \@minus -.2ex}                               {0.3ex \@plus.2ex}
                                   {\large\bf\raggedright}}
\renewcommand\subsection{\@startsection{subsection}{2}{\z@}                                     {-0.25ex\@plus -1ex \@minus -.2ex}
{0.5ex \@plus .2ex}               {\normalsize\bf\raggedright}}
\makeatother


\begin{document}



\maketitle

%\doparttoc % Tell to minitoc to generate a toc for the parts
%\faketableofcontents % Run a fake tableofcontents command for the partocs



\begin{abstract}
Diffusion models have been successful on a range of conditional generation tasks including molecular design and text-to-image generation.
However, these achievements have primarily depended on task-specific conditional training or error-prone heuristic approximations.
Ideally, a conditional generation method should provide exact samples for a broad range of conditional distributions without requiring task-specific training.
To this end, we introduce the \emph{Twisted Diffusion Sampler}, or \emph{TDS}.
TDS is a sequential Monte Carlo (SMC) algorithm that targets the conditional distributions of diffusion models through simulating a set of weighted particles. 
The main idea is to use \emph{twisting}, an SMC technique that enjoys good computational efficiency, to incorporate heuristic approximations without compromising asymptotic exactness. 
We first find in simulation and in conditional image generation tasks that TDS provides a computational statistical trade-off, yielding more accurate approximations with many particles but with empirical improvements over heuristics with as few as two particles.
We then turn to motif-scaffolding, a core task in protein design, using a TDS extension to Riemannian diffusion models; on benchmark tasks, TDS allows flexible conditioning criteria and often outperforms the state-of-the-art, conditionally trained model.\footnote{
 Code: \url{https://github.com/blt2114/twisted_diffusion_sampler}
 }

%Diffusion models have been successful in areas like molecular design and text-to-image generation.
%However, these achievements have primarily depended on expensive, task-specific conditional models or error-prone heuristic approximations. 
%By contrast, an ideal method should provide exact samples from a broad range of conditional distributions without requiring conditional neural network training.
%While the exact conditional distributions 
%of diffusion models are analytically intractable, the sequential structure of diffusion models permits the application of an established toolkit of sequential Monte Carlo (SMC) methods.
%To this end, we introduce the Twisted Diffusion Sampler (TDS), an SMC algorithm that targets conditional distributions of trained diffusion generative models. 
%The main idea behind TDS is to incorporate heuristic approximations as \emph{twisting} functions -- an SMC technique that enjoys good computational efficiency, but without compromising asymptotic exactness.
%We study the performance of TDS on image inpainting, class-conditioned image generation tasks, and diffusion models on Riemannian manifolds, which is crucial for molecular modeling applications.
%In an application to protein design, and across various functional motif-scaffolding benchmark problems, TDS enables more flexible conditioning criteria and outperforms the state-of-the-art.

%We introduce the Twisted Diffusion Sampler (TDS), a sequential Monte Carlo (SMC)  algorithm that accurately samples from the conditional distributions of trained diffusion generative models. 
%While diffusion models have been successful in areas like molecular design and text-to-image generation, these achievements have primarily depended on expensive task-specific conditional models or error-prone heuristic approximations. 
%The work presented here, by contrast, accommodates a broader range of conditioning information and samples exact conditional distributions in a large-compute limit. 
%The main idea behind TDS is to 
%incorporate heuristic approximations as twisting functions and proposals in an exact SMC sampler. The resulting algorithm enjoys good computational efficiency, but without compromising asymptotic accuracy.
%We study the performance of TDS on image inpainting, class-conditioned image generation tasks, and diffusion models on Riemannian manifolds, which is crucial for molecular modeling applications. In an application to protein design, and across various functional motif-scaffolding benchmark problems, TDS enables more flexible conditioning criteria and outperforms the state-of-the-art.
%
%We introduce the Twisted Diffusion Sampler (TDS), a sequential Monte Carlo algorithm that accurately samples from the conditional distributions of trained diffusion generative models. While diffusion models have been successful in areas like molecular design and text-to-image generation, these achievements have primarily depended on task-specific conditional models. These models, however, only accommodate a limited range of conditioning information and are computationally burdensome to train. Alternatively, heuristic approximations are used but can introduce significant approximation errors.  TDS, by contrast, accommodates a broader range of conditioning information and samples exact conditional distributions in a large-compute limit. It incorporates previous heuristic approximations as twisting functions and proposals, which enhance computational efficiency without compromising asymptotic accuracy. We demonstrate TDS's effectiveness in image inpainting and class-conditioned image generation tasks, and its applicability to diffusion models on Riemannian manifolds, crucial for molecular modeling applications. In an application to protein design, TDS enables more flexible conditioning criteria and outperforms the state-of-the-art, conditionally-trained diffusion model across various functional motif-scaffolding benchmark problems.

\end{abstract}

\section{Introduction}\label{sec:intro} 

Conditional sampling is an essential primitive in the machine learning toolkit.
%When combined with a likelihood $p(y\mid x)$, generative models imply conditional distributions $\pmodel(x \mid y) \propto \pmodel(x) p(y\mid x)$ given conditioning information $y$.
One begins with a generative model that parameterizes a distribution $\pmodel(x)$ on data $x$, and then augments the model to include information $y$ in a joint distribution $\pmodel(x, y)=\pmodel(x)p(y\mid x).$
This joint distribution then implies a conditional distribution $\pmodel (x\mid y),$ from which desired outputs are sampled.

For example, in protein design, $\pmodel(x)$ can represent a distribution of physically realizable protein structures, 
$y$ a substructure that imparts a desired biochemical function, 
and samples from $\pmodel(x \mid y)$ are then physically realizable structures that contain the substructure of interest \citep[e.g.][]{trippe2022diffusion,watson2022broadly}.
%In computer vision, $\pmodel(x)$ can be a distribution of images, $y$ a class label,and samples from $\pmodel (x\mid y)$ are then images likely to be classified as label $y$ \citep{ho2022classifier}.  

Diffusion models are a class of generative models that have demonstrated success in conditional generation tasks \citep{ho2020denoising,ramesh2022hierarchical,song2020score,watson2022broadly}.
%Diffusion models approximate sampling from complicated data distributions $\pmodel (x)$ through an iterative refinement process that builds up data from noise \citep{ho2020denoising, song2020score}.
They parameterize distributions $\pmodel (x)$ through an iterative refinement process that builds up data from noise.
When a diffusion model is used for conditional generation, this refinement process is modified to account for conditioning at each step \citep{dhariwal2021diffusion,hovideo,ramesh2022hierarchical,song2020score}.
%Moreover, diffusion models apply not only to Euclidean data, but also to data better described as living on a Riemannian manifold \citep{deriemannian}.
%This type of controllable generation can be formalized as a conditional sampling problem, i.e.\ sample from $\pmodel(x \mid y)$ where $y$ is some conditioning information.

One approach is to incorporate the conditioning information into training \citep[e.g.][]{dhariwal2021diffusion,ramesh2022hierarchical};
one either modifies the unconditional model to take $y$ as input or trains a separate conditional model to predict $y$ from partially noised inputs.
%This approach has provided most of the major successes of diffusion models to date.
However, conditional training requires 
(i) assembling a large set of paired examples of the data and conditioning information $(x, y)$, and (ii) designing and training a task-specific model when adapting to new conditioning tasks.
For example, image inpainting and class-conditional image generation can be both formalized as conditional sampling problems based on the same (unconditional) image distribution $\pmodel(x)$; however, the conditional training approach requires training separate models on two curated sets of conditioning inputs. 

To avoid conditional training, a separate line of work uses heuristic approximations that directly operate on unconditional diffusion models: once an unconditional model $\pmodel(x)$ is trained, it can be flexibly combined with various conditioning criteria to generate customized outputs. These approaches have been applied to \emph{inpainting} problems \citep{meng2021sdedit,song2020score,trippe2022diffusion}, and other inverse problems \citep{bansal2023universal,chung2022diffusion,hovideo, song2023pseudoinverse}.
But it is unclear how well these heuristics approximate the exact conditional distributions they are designed to mimic; for example on inpainting tasks they often fail to return outputs consistent with both the conditioning information and unconditional model \citep{zhang2023towards}.
These concerns are critical in domains that require accurate conditionals. 
In molecular design, for example, even a small approximation error could result in atomic structures with chemically implausible bond distances.

% SMC is revelant. cite smcdiff.  Naive SMC does not work 
This paper develops a practical and asymptotically exact method for conditional sampling from an unconditional diffusion model.
We use sequential Monte Carlo (SMC), a general tool for asymptotically exact inference in sequential probabilistic models  \citep{chopin2020introduction,doucet2001sequential,naesseth2019elements}.
SMC simulates an ensemble of weighted trajectories, or \emph{particles}, through a sequence of proposals and weighting mechanisms. 
These weighted particles then form an asymptotically exact approximation to a desired target distribution. 

The premise of this work is to recognize that the sequential structure of diffusion models permits the application of SMC for sampling from conditional distributions $\pmodel(x \mid y)$.
We design an SMC algorithm that leverages \emph{twisting},
an SMC technique that modifies proposals and weighting schemes to approach the optimal choices \citep{ guarniero2017iterated,whiteley2014twisted}.
While optimal twisting is intractable, we effectively approximate it with recent heuristic approaches to conditional sampling \citep[eg.][]{hovideo},
and correct the errors by the weighting mechanisms. 
The resulting algorithm maintains asymptotic exactness to $\pmodel (x \mid y)$, and empirically it can outperform heuristics alone even with just two particles.

We summarize our contributions: 
(i) We propose a practical SMC algorithm, \emph{Twisted Diffusion Sampler} or TDS,  for asymptotically exact conditional sampling from diffusion models; 
(ii) We show that TDS applies to a range of conditional generation problems, and extends to Riemannian manifold diffusion models; 
(iii) On MNIST inpainting and class-conditional generation tasks we demonstrate TDS's empirical improvements beyond heuristic approaches; and  
(iv) On protein motif-scaffolding problems with short scaffolds TDS provides greater flexibility and achieves higher success rates than the state-of-the-art conditionally trained model. 

%The paper is organized as follows. \Cref{sec:background} gives background materials on diffusion models and sequential Monte Carlo. \Cref{sec:method}  introduces the main methodology, Twisted Diffusion Sampler, and its applications to various conditioning criteria and Riemannian diffusion models. \Cref{sec:related-works} describes the related works. We evaluate our method on synthetic settings and conditional MNIST image generation in \Cref{sec:results}, and benchmark protein motif-scaffolding problems in \Cref{sec:exp-protein}. Finally, \Cref{sec:discussion} summarizes and discusses the limitation of the proposed method. 

\section{Background: Diffusion models and sequential Monte Carlo}\label{sec:background}


%Our goal in this work is to sample from the conditional distribution of a fitted diffusion model $\pmodel (\xzero)$ given some conditioning information $y$, i.e. to sample from \begin{align}    \pmodel (\xzero |y) &\propto \pmodel (\xzero ) \plikelihood{y}{\xzero}\end{align}where $\plikelihood{y}{\xzero}$ is a pre-specified conditional likelihood model. 

\textbf{Diffusion models.} 
A diffusion model generates a data point $\xzero$ by iteratively refining a sequence of noisy data $\xt$, starting from pure noise $\xT$.
This procedure parameterizes a distribution of $\xzero$ as the marginal of a length $T$ Markov chain
\vspace{-0.5em}
\begin{equation}\label{eqn:diffusion_marginal}
        \pmodel(\xzero)=\int p (\xT) \prod_{t=1}^T \pmodel(\xtmo  \mid \xt)d \dt{x}{1:T},
\end{equation}
where $p(\xT)$ is an easy-to-sample noise distribution,
and each $\pmodel(\xtmo \mid \xt)$ is the transition distribution defined by the $(T-t)^\mathrm{th}$ refinement step. 

Diffusion models $\pmodel$ are fitted to match a data distribution $q (\xzero)$ from which we have samples.
To achieve this goal, a \emph{forward process} $q(\xzero) \prod_{t=1}^{T} q(\xt\mid \xtmo)$ is set to gradually add noise to the data, where $q(\xt \mid \xtmo) = \mathcal{N} \left ( \xt; \xtmo ,  \sigma^2\right)$, and $\sigma^2$ is a positive variance.
To fit a diffusion model, one finds $\theta$ such that $p_\theta(\xtmo \mid \xt) \approx q(\xtmo\mid \xt)$, which is the reverse conditional of the forward process.
If this approximation is accomplished for all $t$, and if $T\sigma^2$ is big enough that $q(\xT)$ may be approximated as $q(\xT)\approx \mathcal{N}(0, T\sigma^2)=: p(\xT)$, then we will have $\pmodel(\xzero) \approx q(\xzero)$.
% Note the marginal conditional is given by $q(\xt \mid \xzero) \sim \mathcal{N}(\xt ; \xzero, \bar \sigma_t^2),$ for $\bar \sigma_t^2:= \sum_{t^\prime=1}^T \sigma_{t^\prime}^2.$
% Note that if (1) the $\bar \sigma_T$ is large enough such that $q(\xT \mid \xzero) \approx \mathcal{N}(\xT ; 0, \bar \sigma_T^2)$ we can choose $p(\xT)=\mathcal{N}(\xT ; 0, \bar \sigma_T^2),$ and 
% (2) $\pmodel(\xt \mid \xtpo) \approx q(\xt \mid \xtpo)$ for every $t=0, 1, \dots, T-1,$ then $\pmodel(\xzero)\approx q(\xzero).$
%The key idea behind fitting a diffusion model is to note that  if (1) the $\sigma_t$'s are large enough, $$q(\xT \mid \xzero) \sim \mathcal{N}(\xT ; \xzero, \bar \sigma_T^2)\approx \mathcal{N}(\xT ; 0, \bar \sigma_T^2), \footnote{Note that the dependence on $\xzero$ is lost for large enough $t.$ E.g. $\lim_{\bar\sigma^2 \rightarrow \infty} d_{TV}\left(\mathcal{N}(\xzero, \bar \sigma^2), \mathcal{N}(0, \bar \sigma^2)\right)= 0$.} $$ so we can choose $p(\xT)=\mathcal{N}(\xT ; 0, \bar \sigma_T^2),$ and  (2) $\pmodel(\xt \mid \xtpo) \approx q(\xt \mid \xtpo)$ for every $t=0, 1, \dots, T-1,$ then $\pmodel(\xzero)\approx q(\xzero).$

In particular, when $\sigma^2$ is small enough 
%if $\sigma_t^2$ is much smaller than $\bar \sigma_t^2 := \sum_{t^\prime=1}^t \sigma_{t^\prime}^2$ 
then the reverse conditionals of $q$ are approximately Gaussian,
\begin{align}\label{eq:q_reverse_conditional}
\begin{split}
q(\xtmo \mid \xt) 
    &\approx \mathcal{N}\left(\xtmo; \xt + \sigma^2\nabla_{\xt} \log q(\xt), \sigma^2 \right),
\end{split}
\end{align}
where $q(\xt)=\int q(\xzero) q(\xt \mid \xzero) d\xzero$ and $\nabla_{\xt} \log q(\xt)$ is known as the (Stein) score \citep{song2020score}.
To mirror \cref{eq:q_reverse_conditional}, diffusion models parameterize $\pmodel(\xtmo \mid \xt)$ via a \emph{score network} $\scoremodel(\xt, t)$
\begin{align}\label{eq:unconditional-diffusion-transition}
    \pmodel(\xtmo \mid \xt) := \mathcal{N}\left(
    \xtmo;\xt   +  \sigma^2 \scoremodel(\xt, t),\sigma^2 \right). 
\end{align}
When $\scoremodel (\xt, t)$ is trained to approximate $ \nabla_{\xt} \log q(\xt )$, we have $\pmodel(\xtmo \mid \xt) \approx q(\xtmo \mid \xt)$. 

Notably, approximating the score is equivalent to learning a \emph{denoising} neural network $\denoiser (\xt, t)$ to approximate $\E_q[\xzero \mid \xt]$. 
The reason is that by Tweedie's formula \citep{efron2011tweedie,robbins1992empirical} $\nabla_{\xt}\log q(\xt) = (\E_q[\xzero \mid \xt] - \xt)/t\sigma^2$ 
and one can set $\scoremodel(\xt,t):=  (\denoiser (\xt,t)- \xt)/t\sigma^2$. 
The neural network $\denoiser (\xt,t)$ may be learned by denoising score matching (see \citep[e.g.][]{ho2020denoising,vincent2011connection}). 
For the remainder of paper we drop the argument $t$ in $\denoiser$ and  $\scoremodel$ when it is clear from context. 

\Cref{sec:supp_method} generalizes the formulation above to diffusion process formulations that are commonly used in practice.



\textbf{Sequential Monte Carlo.}
%\label{sec:SMC_background}
Sequential Monte Carlo (SMC) is a general tool to approximately sample from a sequence of distributions on variables $\dt{x}{0:T}$, terminating at a final target of interest
\citep{chopin2020introduction,doucet2001sequential,gordon1993novel,naesseth2019elements}. 
SMC approximates these targets by generating a collection of $K$ \emph{particles} $\{\dt{x}{t}_k\}_{k=1}^K$ across $T$ steps of an iterative procedure.
The key ingredients are proposals $r_T(\xT)$, $\{r_t(\xt \mid \xtpo)\}_{t=0}^{T-1}$ 
and weighting functions $\weight_T(\xT)$, $\{\weight_t(\xt, \xtpo)\}_{t=0}^{T-1}$. 
At the initial step $T$, one draws $K$ particles of $\xT_k \sim \proposal_T(\xT)$ and sets $\dt{w}{T}_k:= \weight_T(\xT_k)$, 
%setting $\dt{w}{T}_k:= \target_T(\xT_k) / \proposal_T (\xT_k)$
and sequentially repeats the following for $t=T-1, \dots, 0$: 
\vspace{-0.2em}
\begin{itemize}
    \item \emph{resample} \qquad $\{ \dt{x}{t+1:T}_k \}_{k=1}^K \sim \textrm{Multinomial}(\{ \dt{x}{t+1:T}_k \}_{k=1}^K; \{\dt{w}{t+1}_k \}_{k=1}^K)$ %\footnote{Other resampling strategies can be used as well. See Appendix for discussion.}
    \item  \emph{propose} \qquad $\xt_k \sim \proposal_t(\xt \mid \xtpo_k), \quad k=1, \cdots, K$ 
    \item \emph{weight}  \qquad \quad $\dt{w}{t}_k := \weight_t(\xt_k, \xtpo_k), \quad k=1,\cdots, K$
\end{itemize}
\vspace{-0.2em}

The proposals and weighting functions together define a sequence of \emph{intermediate target} distributions,
%$\nu_t,$ each defined on $\dt{x}{0:T}$ by the proposals, and the weighting functions at the first $T-t$ steps as
%    The weighted set of particles form an approximation to the target distribution $\target_t$ defined by $\proposal_t$ and $\weight_t$ 
    %for $t=T,\cdots, 0$, that is,
%\begin{align}\label{eq:final-target}
\begin{equation}\label{eq:final-target}
    %\target_t(\dt{x}{0:T}) &:= \frac{1}{L_t}
    \target_t(\dt{x}{0:T}) 
    :=\frac{1}{\mathcal{L}_t} 
        \Big[ \proposal_T(\dt{x}{T})  \prod_{t^\prime=0}^{T-1} \proposal_{t^\prime}(\dt{x}{t^\prime}\mid  \dt{x}{t^\prime+1})  \Big] 
        \Big[ \weight_T(\dt{x}{T})
        \prod_{t^\prime =t}^{T-1} \weight_{t^\prime}(\dt{x}{t^\prime}, \dt{x}{t^\prime+1}) \Big]  
\end{equation}
where $\mathcal{L}_t$ is a normalization constant.
A classic example that SMC applies is the state space model 
\citep[Chapter 1]{naesseth2019elements} 
that describes a distribution over a sequence of latent states $\dt{x}{0:T}$ and observations  $\dt{y}{0:T}$. 
Each intermediate target $\target_t$ is constructed to be the posterior $p(\dt{x}{0:T}\mid \dt{y}{t:T})$ given the first $T-t+1$ observations, and the final target $p(\dt{x}{0:T} \mid \dt{y}{0:T})$ is the  posterior given all observations. 


%Then $\target_t$ is a target posterior $p(\dt{x}{0:T}\mid \dt{y}{t:T})$ after the first $T-t$ observations in a filtering equation where $p(\xT)$ and each $p(\xt \mid \xtpo)$ and $p(\dt{y}{t}\mid \dt{x}{t})$ are known but the posterior is intractable.
%This example also points to the connection to diffusion and to our conditional sampling problem of interest.\BLT{\footnote{IIUC, John suggested a line like this last sentence. But I think we can drop it, since making it sufficiently precise is essentially the goal of the beginning of the next section.}}

The defining property of SMC is that the weighted particles at each $t$ form discrete approximations
$(\sum_{k^\prime=1}^K \dt{w}{t}_{k^\prime})^{-1}\sum_{k=1}^K w_t^k \delta_{\dt{x}{t}_k} (\dt{x}{t})$  (where $\delta$ is a Dirac measure)
to $\target_t (\dt{x}{t})$ that become arbitrarily accurate in the limit that many particles are used \citep[Proposition 11.4]{chopin2020introduction}.
So, by choosing $\proposal_t$ and $\weight_t$ so that $\target_0(\xzero)$ matches the desired distribution, one can guarantee arbitrarily low approximation error in the large compute limit.% wherein many particles are simulated.


%Each step of SMC involves importance sampling of a target $\target_{t}$ with $\target_{t+1}$ as the proposal, and so the procedure inherits the operating characteristics of importance sampling.
%This includes an exponential dependence of the number of samples needed for accurate estimation of $\target_{t}$ on the Kullback-Leibler (KL) divergence of each $\target_{t+1}$ from $\target_{t}$ \citep{chatterjee2018sample}.
%In the context of  \Cref{eq:final-target}, the KL divergence reduces to
%$$\mathrm{KL}(\target_{t} \| \target_{t+1})=\ln \E_{\target_{t}}[\weight_{t}(\xt, \xtpo)\inv] - \E_{\target_{t}}[\ln \weight_{t}(\xt, \xtpo)\inv].$$
%\BLT{\footnote{If we choose to keep this equation, we should probably add a derivation of it to the appendix.}}This divergence will be large when a small fraction of particles receive a disproportionately large weight.
%In that case only a small fraction of the particles receive nontrivial weights, the rest will be discarded, and 
%the procedure provides a poor approximation of the target. 

%See \Cref{sec:SMC_details} for details.  

%the number of particles needed for accurate estimation of each $\target_t$ is exponential in the KL divergence between successive targets \citep{chatterjee2018sample}: \begin{align}\label{eq:kl-intermediate-targets}\mathrm{KL}(\target_{t-1} \| \target_t)=\ln \E_{\target_{t-1}}[\weight_{t-1}\inv(\xt, \xtpo)] - \E_{\target_{t-1}}[\ln \weight_{t-1}\inv(\xt, \xtpo)].\end{align}
%This divergence will be large when a small fraction of particles receive a disproportionately large weight. 
%Otherwise only a small fraction of the particles will receive nontrivial weights, and the rest will be discarded. 
%Then exponential amount of particles in $T$ are needed for accurate approximation of $\target_0$ (see appendix for a formal discussion). 
%In that case most particles are discarded, the procedure provides only small effective number of samples and therefore a poor approximation of the target distribution. 


%More formally, each step of SMC involves importance sampling of a target $\target_{t-1}$ with $\target_t$ as the proposal, and so the procedure inherits the operating characteristics of importance sampling. Notably this includes an exponential dependence of the number of samples needed for accurate estimation on the Kullback-Leibler divergence of each $\target_t$ from $\target_{t-1}$ \citep{chatterjee2018sample}.In the context of \Cref{eqn:final-target} we see that $$\mathrm{KL}(\target_{t-1} \| \target_t)=\ln \E_{\target_{t-1}}[\weight_{t-1}\inv(\xt, \xtpo)] - \E_{\target_{t-1}}[\ln \weight_{t-1}\inv(\xt, \xtpo)].$$ This divergence will be large when a small fraction of particles receive a disproportionately large weight.





\section{\methodname{}: SMC sampling for diffusion model conditionals }\label{sec:method}
%We further assume the conditional independence of $\dt{x}{1:T}$ and $y$ given $\dt{x}{0}$.


%\textbf{Problem statement.}
%Our goal in this work is to approximate conditional distributions of diffusion generative models.
%We introduce conditioning information $y$ through a smooth likelihood function $\plikelihood{y}{\xzero}.$ 
%We extend our joint model to include $y$ and $\plikelihood{y}{\xzero}$ as $\pmodel(\dt{x}{0:T}, y)=\pmodel(\dt{x}{0:T})\pmodel(y\mid \xzero)$ where we define $\pmodel(y \mid \xzero)=\plikelihood{y}{\xzero}$. The goal is to sample from the conditional $\pmodel(\xzero \mid y)$. 
%%and write the conditional of interest as $\pmodel(\xzero \mid y).$


%To this end, we

%In this section, we develop \methodname{} (\methodshort{}), a practical SMC algorithm targeting $\pmodel(\xzero \mid y)\propto \pmodel(\xzero)\plikelihood{y}{\xzero},$ the conditional distribution of a diffusion model $\pmodel(\xzero)$ defined though a likelihood denoted $\plikelihood{y}{\xzero}$. 
Consider conditioning information $y$ associated with a given likelihood function $\plikelihood{y}{\xzero}.$
We embed $y$ in a joint model over $\dt{x}{0:T}$ and $y$ as
$\pmodel(\dt{x}{0:T},y) = \pmodel(\dt{x}{0:T}) \plikelihood{y}{\xzero},$ 
where $y$ and $\dt{x}{1:T}$ are conditionally independent given $\xzero.$
Our goal is to sample from the conditional $\pmodel(\xzero\mid y)$.% implied by this joint model.

In this section, we develop \methodname{} (\methodshort{}), a practical SMC algorithm targeting $\pmodel(\xzero \mid y).$ 
%a joint model 
First, we describe how the Markov structure of diffusion models permits a factorization of an extended conditional distribution to which SMC applies.
Then, we show how a diffusion model's denoising predictions support the application of \emph{twisting}, an SMC technique in which one uses proposals and weighting functions that approximate the ``optimal'' ones.
%We show that \methodshort{} provides asymptotically exact estimates of conditional distributions and applies to a broad range of conditioning criteria.
%After introducing a first approximation, 
Lastly, we extend TDS to certain ``inpainting'' problems where $\plikelihood{y}{\xzero}$ is not smooth,
and to Riemannian diffusion models.


\ifnum\workshopsubmission=1
TDS builds on diffusion generative models and SMC;
we provide background and specify the conventions we adopt for diffusion models and SMC in \Cref{sec:background}.
Empirical results are left to \Cref{sec:results}.
\fi 

\begin{algorithm}[t]
\SetAlgoLined
\SetKwFor{For}{for}{}{}
\SetKwBlock{Begin}{function}{end}
\BlankLine 

\For{$k=1:K$}{
\vspace{-0.1em}
\BlankLine
$\xT_k \sim p(\xT), \quad w_k \leftarrow \dt{\twistfcn_k}{T} =\twist(y\mid \xT_k)$ \quad \tcp{initial proposal and weight} \label{alg_line:initial_twisting_function}
}
\vspace{-0.1em}
\BlankLine
\For{$t = T-1,\cdots, 0$}{
     \vspace{-0.1em}
    \BlankLine
     $ \{\dt{x}{t+1}_k, \dt{\twistfcn}{t+1}_k\}_{k=1}^K \sim \mathrm{Multinomial}\left(\{ \dt{x}{t+1}_{k}, \dt{\twistfcn}{t+1}_{k}\}_{k=1}^K; \{ w_{k}\}_{k=1}^K \right)$ \quad \tcp{resample}
     \label{alg_line:resampling}
     \vspace{-0.1em}
     \BlankLine
     \For{$k=1:K$}{
          \vspace{-0.1em}
    \BlankLine
    $\tilde s_k =  (\denoiser(\xtpo_k) - \xtpo_k)/t\sigma^2 + \nabla_{\xtpo_k} \log \dt{\twistfcn}{t+1}_k$ \quad \tcp{conditional score approx.} \label{alg_line:conditional_score}
         \vspace{-0.1em}
    \BlankLine
    $\xt_k \sim \pTwist(\cdot {\mid} \xtpo_k, y):=\mathcal{N}\left(\xtpo_k + \sigma^2 \tilde s_k, \proposalVar \right) 
    $ \quad \tcp{proposal}\label{alg_line:proposal}
         \vspace{-0.1em}
    \BlankLine
     $\dt{\twistfcn_k}{t} \leftarrow \twist(y\mid \xt_k)$ \quad \tcp{twisting function (in eqs.\ \eqref{eq:diffusion-twisting-functions}, \eqref{eqn:inpaint_twist}, or \eqref{eqn:mask_dof_twist})}
     \vspace{-0.1em}
    \BlankLine
    $ w_k \leftarrow p_\theta(\xt_k {|} \xtpo_k)\dt{\twistfcn}{t}_k/ [\pTwist(\xt_k {|} \xtpo_k, y)\dt{\twistfcn}{t+1}_k]$  \quad \tcp{weight} 
    } 
}
\caption{Twisted Diffusion Sampler (TDS)}
\label{alg:TDS}
 \vspace{0.1em}
\end{algorithm}

\subsection{Conditional diffusion sampling as an SMC procedure}\label{subsec:SMC_framing}
The Markov structure of the diffusion model permits a factorization that is recognizable as the final target of an SMC algorithm.
We  write the conditional distribution, extended to  include $\dt{x}{1:T},$ as
\begin{align}\label{eq:diffusion-final-target}
    \pmodel(\dt{x}{0:T} \mid y) &=  
    \ifnum\workshopsubmission=0
        \frac{\pmodel (\dt{x}{0:T} ,y)}{\pmodel(y)} = 
    \fi
    \frac{1}{\pmodel(y)}\Big[p(\xT) \prod_{t=0}^{T-1} \pmodel (\xt \mid \xtpo )  \Big] \plikelihood{y}{\xzero}
\end{align}
with the desired marginal, $\pmodel(\xzero \mid y).$
Comparing the diffusion conditional of \cref{eq:diffusion-final-target} to the SMC target of \cref{eq:final-target} suggests SMC can be used.

For example, consider a first attempt at an SMC algorithm.
Set the proposals as
%\begin{equation}
$\proposal_T(\xT)= p(\xT)\quad \mathrm{and} \quad
\proposal_t(\xt \mid \xtpo) = \pmodel (\xt \mid \xtpo)\quad \mathrm{for}\quad 1\le t \le T,$
%\end{equation}
and weighting functions as  
%\begin{equation}
$\weight_T(\xT)=\weight_t(\xt, \xtpo)=1\quad \mathrm{for} \quad 1\le t \le T \quad \mathrm{and} \quad \weight_0 (\xzero, \dt{x}{1}) = \plikelihood{y}{\xzero}.$
%\end

Substituting these choices into \cref{eq:final-target} results in the desired final target  $\target_0 = \pmodel(\dt{x}{0:T}\mid y)$ with normalizing constant $\mathcal{L}_0 = \pmodel(y)$. 
%and so defines an SMC sampler whose final target is $\pmodel(\dt{x}{0:T} \mid y),$ and with normalization constant $L_0=\pmodel(y),$ as desired.
As a result, the associated SMC algorithm produces a final set of $K$ samples and weights $\{ \xzero_k; \dt{w}{0}_k\}_{k=1}^K$ that provides an asymptotically accurate approximation $\mathbb{P}_K (\xzero \mid y):=(\sum_{k}^K w^0_k)\inv \sum_{k=1}^K \dt{w}{0}_k \delta_{\xzero_k}(\xzero)$ to the desired $\pmodel(\xzero \mid y).$ 
%Specifically, if we let $\mathbb{P}_K=(\sum_{k}^K w^0_k)\inv \sum_{k=1}^K \dt{w}{0}_k \delta_{\xzero_k}$ be the $K$ sample approximation, then $\mathbb{P}_K$ converges weakly to $\pmodel(\xzero \mid y)$ as $K$ goes to infinity. \BLT{\footnote{A previous version equated the discrete measure and density pointwise; I've modified this to say weak convergence instead which can be more correct but which we don't define until the appendix...  So this can be improved.}}
%\begin{align}\label{eq:naive-is}
%\lim_{K\rightarrow \infty} \int_A (\sum_{k}^K w^0_k)\inv \sum_{k=1}^K \dt{w}{0}_k \delta_{\xzero_k} (\xzero)d\xzero 
%= \int_A \target_0 (\xzero) d\xzero = \int_A \pmodel(\xzero \mid y) d \xzero
%\ifnum\workshopsubmission=0
%    \qquad \mathrm{with}\ \ \dt{x}{0}_k \overset{i.i.d.}{\sim} \pmodel(\dt{x}{0}).
%\fi
%\end{align}
%\ifnum\workshopsubmission=1
%with $\dt{x}{0}_k \overset{i.i.d.}{\sim} \pmodel(\dt{x}{0})$.
%\fi

%And when $K \to \infty$, this approximation becomes arbitrarily accurate.
The approach above is simply importance sampling with proposal $\pmodel (\dt{x}{0:T})$; with all intermediate weights set to 1, one can skip resampling steps to reduce the variance of the procedure.
Consequently, this approach will be impractical if $\pmodel(\xzero\mid y)$ is too dissimilar from $\pmodel(\xzero)$
as only a small fraction of unconditional samples will have high likelihood: the number of particles required for accurate estimation of $\pmodel(\dt{x}{0:T} \mid y)$ is exponential in $\KL{\pmodel(\dt{x}{0:T} \mid y)}{\pmodel (\dt{x}{0:T})}$ \citep{chatterjee2018sample}. 

\subsection{Twisted diffusion sampler}\label{subsec:PF_twisted}

%In this part, we show how we can \emph{twist} the choice of the above proposals $\proposal_t$ and weighting functions $\weight_t$ to circumvent the exponential complexity of particles.

\emph{Twisting} is a technique in the SMC literature intended to reduce the number of particles required for good approximation accuracy \citep{guarniero2017iterated,heng2020controlled}. %\citep[Chapter 3] {naesseth2019elements}.
%The key idea is to relieve the discrepancy between successive intermediate targets by bringing them closer to the final target. 
%, so only a practical amount of particles are need for accurate estimation  (recall this amount scales exponentially with $\mathrm{KL}(\target_{t-1} \| \target_{t})$ for each step $t$)
Loosely, it introduces a sequence of \emph{twisting functions} that modify the naive proposals and weighting functions, so that the resulting intermediate targets are closer to the final target of interests. 
% This idea is achieved by introducing a sequence of \emph{twisting functions} that define new \emph{twisted} proposals and \emph{twisted} weighting functions. 
We refer the reader to \citet[Chapter 3.2]{naesseth2019elements} for background on twisting in SMC.

\textbf{Optimal twisting.}
%For example, consider twisting the  proposals $\proposal_t$ in the naive procedure  described in \Cref{subsec:SMC_framing} to define 
Consider defining the twisted proposals $\proposal_t^*$ by multiplying the naive proposals $\proposal_t$ described in \Cref{subsec:SMC_framing} by $\pmodel (y \mid \xt)$ as
\ifnum\workshopsubmission=0
\begin{align}\label{eq:optimal-twisted-proposals}
\proposal_T^* (\xT) &\propto p(\xT) \pmodel(y {|} \xT) \ \mathrm{and} \ 
     \proposal_t^* (\xt {|} \xtpo) \propto \pmodel(\xt {|} \xtpo) \pmodel(y{|} \xt) \quad \mathrm{for}  \qquad 0\leq t < T.
\end{align}
\else
\begin{align}\label{eq:optimal-twisted-proposals}
\proposal_T^* (\xT) &\propto \proposal_T (\xT) \twistfcn_T^* (\xT) \ \mathrm{and} \ 
\\
     \proposal_t^* (\xt \mid \xtpo) &\propto \proposal_t (\xt \mid \xtpo) \twistfcn_t^* (\xt) \ \mathrm{for}  \qquad 0\leq t < T.
\end{align}
\fi
The factors $\pmodel(y \mid \xt)$ are the \emph{optimal} twisting functions
because they permit an SMC sampler that draws exact samples from $\pmodel(\dt{x}{0:T} \mid y)$ even when run with a single particle.

To see that a single particle is an exact sample, by Bayes rule, the proposals in \cref{eq:optimal-twisted-proposals} reduce to 
\ifnum\workshopsubmission=0
\begin{align}\label{eq:optimal-twisted-proposals-reduced}
     \proposal_T^{*} (\xT) &= \pmodel(\xT \mid y)\ \mathrm{and}  \ 
     \proposal_t^{*} (\xt \mid \xtpo) = \pmodel ( \xt \mid \xtpo, y)\quad \mathrm{for}  \qquad 0\leq t < T.
\end{align}
 \else
\begin{align}\label{eq:optimal-twisted-proposals-reduced}
     \proposal_T^{*} (\xT) &= \pmodel(\xT \mid y)\ \mathrm{and}  \\
     \proposal_t^{*} (\xt \mid \xtpo) &= \pmodel ( \xt \mid \xtpo, y)\quad \mathrm{for}  \qquad 0\leq t < T.
\end{align}
 \fi
As a result, 
if one samples $\xT\sim \proposal_T^{*}(\xT)$ and $\xt\sim \proposal_t^{*} (\xt \mid \xtpo)$
for  $t=T-1,\dots, 0$,
by (i) the law of total probability and (ii) the chain rule of probability, 
 one obtains $\xzero \sim \pmodel(\xzero \mid y)$ as desired.
%If we choose twisted weighting functions as $\weight^*_T(\xT)=\weight^*_t(\xt, \xtpo) = 1$ for  $t=T-1,\dots, 0$,  
%then resulting twisted targets $\target_t^*$ are all equal to $\pmodel(\dt{x}{0:T}\mid y)$. 
%Specifically, substituting  $\proposal_t^{*}$ and $\weight_t^{*}$ into \cref{eq:final-target} gives 
%\ifnum\workshopsubmission=0
%\begin{align}\label{eqn:opt_intermediate}
%    \target_t^{*} (\dt{x}{0:T}) = \frac{1}{1}  \left[\pmodel (\xT \mid y) \prod_{t=0}^{T-1}\pmodel ( \xt \mid \xtpo, y)\right]
%    \left[1 \cdot \prod_{t=0}^{T-1} 1 \right]
%    =\pmodel(\dt{x}{0:T}\mid y), \quad \mathrm{for}\quad 0\le t \le T.
%\end{align}
%\else
%, for $t=0,\cdots, T,$
%\begin{align}\label{eqn:opt_intermediate}
%    \target_t^{*} (\dt{x}{0:T}) = \pmodel (\xT \mid y) \prod_{t=0}^{T-1}\pmodel ( \xt \mid \xtpo, y)
%    =\pmodel(\dt{x}{0:T}\mid y), 
%\end{align}
%\fi

%As a consequence of \cref{eqn:opt_intermediate}, $\mathrm{KL} (\target_{t-1}^{*} \| \target_t^{*})=0$ for each $t,$ and essentially one particle is sufficient for exact conditional sampling.
However, we cannot readily sample from each $r^*_t$
because $\pmodel(y\mid \xt)$ is not analytically tractable. The challenge is that $\pmodel(y\mid \xt) = \int \plikelihood{y}
{\xzero} \pmodel(\xzero \mid \xt) d\xzero $ depends on $x_t$ through $\pmodel(\xzero\mid \xt)$. The latter in turn requires marginalizing out $\dt{x}{1},\dots, \dt{x}{t-1}$ from the joint density $\pmodel(\dt{x}{0:t-1} \mid \xt),$ 
whose form depends on $t$ calls to the neural network $\denoiser$.

%However, we cannot generate samples from  the optimally twisted proposals $\proposal_t^*$ in \cref{eq:optimal-twisted-proposals-reduced}.
%%But we can develop a practical sampler through tractable twisting functions.   
%$\nu_0(\dt{x}{0:T})=$
%\begin{align}\label{eq:optimal-twisting-functions}
%    \psi^*_t (\xt) := \pmodel (y \mid \xt) =\int \plikelihood{y}{\xzero} \pmodel (\xzero \mid \xt)  d\xzero. 
%\end{align}
%which define the optimally twisted proposal as 
%
%Notably, the \emph{optimal} twisting functions are chosen as 
%\begin{align}\label{eq:optimal-twisting-functions}
%    \psi^*_t (\xt) := \pmodel (y \mid \xt) =\int \plikelihood{y}{\xzero} \pmodel (\xzero \mid \xt)  d\xzero. 
%\end{align}
%which define the optimally twisted proposal as 
%\begin{align}\label{eq:optimal-twisted-proposals}
%     \proposal_t^{\psi^*} (\xt \mid \xtpo) &\propto \proposal_t (\xt \mid \xtpo) \psi_t^* (\xt) 
%     %= \pmodel (\xt \mid \xtpo) \pmodel (y \mid \xt)  
%     \propto \pmodel ( \xt \mid \xtpo, y), \qquad 0\leq t < T \\ 
%     \proposal_T^{\psi^*} (\xT) &\propto \proposal_T(\xT) \psi_T^* (\xT)\propto \pmodel(\xT \mid y), 
%\end{align}
%and the optimally twisted weighting functions as 1. 

\textbf{Tractable twisting.}
To avoid this intractability, we approximate the optimal twisting functions by
\begin{align}\label{eq:diffusion-twisting-functions}
     \twist(y\mid \xt) := \plikelihood { y}{\denoiser (\xt)}    \approx  \pmodel (y \mid \xt),
\end{align}
which is the likelihood function evaluated at $\denoiser(\xt)$, the denoising estimate of $\xzero$ at step $t$ from the diffusion model.
This tractable twisting function $ \twist(y\mid \xt)$ is the key ingredient needed to define the Twisted Diffusion Sampler (TDS, \Cref{alg:TDS}).
We motivate and develop its components below.

The approximation  in \cref{eq:diffusion-twisting-functions} offers two favorable properties.
First, $\twist(y \mid \xt)$ is computable because it depends on $\xt$ only though one call to $\denoiser,$ instead of an intractable integral over many calls as in the case of optimal twisting.
Second, $\twist(y\mid \xt)$ becomes an increasingly accurate approximation of $\pmodel(y \mid \xt),$  as $t$ decreases and $\pmodel(\xzero \mid \xt)$ concentrates on $\E_{\pmodel}[\xzero \mid \xt]$, which $\denoiser$ is fit to approximate;
at $t=0$, where we can choose $\denoiser(\xzero)=\xzero$, we obtain $\twist(y\mid \xzero) = \plikelihood{y}{\xzero}.$



We next use \cref{eq:diffusion-twisting-functions}
to develop a sequence of twisted proposals $\pTwist(\xt \mid \xtpo, y),$ to approximate the optimal proposals $\pmodel(\xt \mid \xtpo, y)$ in \cref{eq:optimal-twisted-proposals-reduced}.
Specifically, we define twisted proposals as
\begin{align}\label{eq:diffusion-twisted-proposals}
    \pTwist(\xt \mid \xtpo, y) &:= \mathcal{N} \left(\xt; \xtpo + \sigma^2 \scoremodel(\xtpo, y), \proposalVar \right), \\
\mathrm{where}\quad \scoremodel(\xtpo, y) &:= \scoremodel(\xtpo) + \nabla_{\xtpo} \log  \twist (y\mid \xtpo)\label{eqn:conditional_score_approx}
\end{align}
is an approximation of the conditional score, $\scoremodel(\xtpo, y) \approx \nabla_{\xtpo} \log \pmodel(\xtpo \mid y),$ and
$\proposalVar$ is the proposal variance.
For simplicity one could choose $\proposalVar=\sigma^2$ to match the variance of $\pmodel (\xt \mid \xtpo)$.
%\lw{discuss whether we use $\tilde \sigma_t^2$ or $\sigma_t^2$ in \cref{eq:diffusion-twisted-proposals}} 

\Cref{eq:diffusion-twisted-proposals} builds on previous works (e.g.\ \citep{song2020score,shi2022conditional}) that seek to approximate the reversal of a conditional forward process.
%;appendix XXX shows how this expression obtains from a Taylor approximation of $\log \pmodel(\xt \mid \xtpo, y).$ \lw{should we keep this paragraph now that we have reference above?}
The gradient in \cref{eqn:conditional_score_approx} is computed by back-propagating through  $\denoiser(\xtpo).$
We further discuss this technique, which has been used before \citep[e.g.\ in][]{hovideo,song2023pseudoinverse}, in \Cref{sec:related-works}. 


\paragraph{Twisted targets and weighting functions.} 
Because $\pTwist(\xt \mid \xtpo, y)$ will not in general coincide with the optimal twisted proposal $\pmodel (\xt \mid \xtpo)$, 
we must introduce non-trivial weighting functions to ensure the resulting SMC sampler converges to the desired final target.
In particular, we define \emph{twisted} weighting functions as 
\ifnum\workshopsubmission=0
\begin{align}\label{eq:diffusion-twisted-weigthing-functions}
    &\weight_t (\xt , \xtpo ):= \frac{\pmodel (\xt \mid \xtpo) \twist(y \mid \xt)}{\twist(y \mid \xtpo) \pTwist(\xt \mid \xtpo, y)}, \qquad t=0,\dots, T-1 %\\
    %&\weight_T^{\psi} (\xT) := \pTwist(\xT).
\end{align}
\else
\begin{align}\label{eq:diffusion-twisted-weigthing-functions}
    &\weight_t (\xt , \xtpo ):= \frac{\pmodel (\xt \mid \xtpo) \twist(y \mid \xt)}{
    \pTwist(\xt \mid \xtpo, y)\twist(y \mid \xtpo)}
\end{align}
for $t=0,\dots, T-1,$
\fi
and $\weight_T (\xT) := \pTwist(y \mid \xT)$. 
The weighting functions in \cref{eq:diffusion-twisted-weigthing-functions} recover the optimal, constant weighting functions if all other approximations at play are exact. 

These tractable twisted proposals and weighting functions define intermediate targets that gradually approach the final target $\target_0=\pmodel (\dt{x}{0:T} \mid y)$.
Substituting into \cref{eq:final-target} each $\pTwist(\xt \mid \xtpo, y)$ in \cref{eq:diffusion-twisted-proposals} for $\proposal_t(\xt\mid \xtpo)$ and $\weight_t(\xt, \xtpo)$ in \cref{eq:diffusion-twisted-weigthing-functions} %for each $\weight_t(\xt, \xtpo)$ 
and then simplifying we obtain the intermediate targets  
\begin{align}\label{eq:twisted-intermediates}
   \nu_t(\dt{x}{0:T}) %&\propto\left[p(\xT)\prod_{t^\prime=0}^T \pTwist(\xt \mid \xtpo, y) \right] \left[\twist(y\mid \xT) \prod_{t^\prime=t}^{T-1}  \frac{\pmodel (\xt \mid \xtpo) \twist(y \mid \xt)}{\twist(y \mid \xtpo) \pTwist(\xt \mid \xtpo, y)}\right] \label{eq:twisted-intermediates-1} \nonumber \\
   &\propto  \pmodel(\dt{x}{0:T}\mid y) \Big[\frac{\twist(y\mid \xt)}{\pmodel(y \mid \xt)} \prod_{t^\prime=0}^{t-1} \frac{\pTwist(\dt{x}{t^\prime} \mid \dt{x}{t^\prime +1}, y)}{
   \pmodel(\dt{x}{t^\prime}\mid \dt{x}{t^\prime + 1}, y)} \Big].
\end{align}
The right-hand bracketed term in \cref{eq:twisted-intermediates} 
can be understood as the discrepancy of $\target_t$ from the final target $\target_0$ accumulated from step $t$ to $0$ (see \Cref{sec:twisted_intermediates_derivation} for a derivation).
As $t$ approaches $0,$
$\twist(y\mid \xt)$ improves as an approximation of $\pmodel(y \mid \xt),$ and the $t$-term product inside the bracket consists of fewer terms -- the latter accounts for the discrepancy between practical and optimal proposals. 
Finally, at $t=0,$ because $\twist(y\mid \xzero) = \pmodel(y \mid \xzero)$ by construction, \cref{eq:twisted-intermediates} reduces to $\target_0(\dt{x}{0:T}) = \pmodel(\dt{x}{0:T}\mid y),$
as desired.

\textbf{The \methodshort{} algorithm and  asymptotic exactness.}
Together, the twisted  proposals $\pTwist(\xt\mid \xtpo, y)$  and  weighting functions $\weight_t(\xt, \xtpo)$ lead to \emph{Twisted Diffusion Sampler}, or TDS (\Cref{alg:TDS}).
While Algorithm~\ref{alg:TDS} states multinomial resampling for simplicity, in practice other resampling strategies (e.g.\ systematic \citep[Ch. 9]{chopin2020introduction}) may be used as well.
Under additional conditions, TDS provides arbitrarily accurate estimates of $\pmodel (\xzero \mid y)$.
Crucially, this guarantee does not rely on assumptions on the accuracy of the approximations used to derive the twisted proposals and weights. \Cref{sec:supp_method} provides the formal statement with complete conditions and proof.



\begin{theorem}\label{prop:twisted_SMC} (Informal) 
Let $\mathbb{P}_{K}(\xzero)=(\sum_{k^\prime}^K w_{k^\prime})\inv \sum_{k=1}^K w_k\delta_{\dt{x_k}{0}}(\xzero)$ denote the discrete measure defined by the particles and weights returned by \Cref{alg:TDS} with $K$ particles.
Under regularity conditions on the twisted proposals and weighting functions, $\mathbb{P}_K(\xzero)$
converges setwise to $\pmodel(\xzero\mid y)$ as $K$ approaches infinity.
\end{theorem}

%\Cref{sec:supp_method} provides the formal statement with complete conditions and proof.





\subsection{TDS for inpainting, additional degrees of freedom}\label{subsec:conditioning-problems}
The twisting functions $\twist(y \mid \xt):=\plikelihood{y}{\denoiser(\xt)}$ we introduced above are one convenient option,
but are sensible only when $\plikelihood{y}{ \denoiser(\xt)}$ is differentiable and strictly positive.
We now show how alternative twisting functions lead to proposals and weighting functions that address inpainting problems and  more flexible conditioning specifications.
In these extensions, Algorithm~\ref{alg:TDS} still applies with the new definitions of twisting functions.
\Cref{sec:supp_method} provides additional details, including the adaptation of TDS to variance preserving diffusion models \citep{ho2020denoising}.


%We present three scenarios of conditioning problems where TDS can be applied to. The main methodology is demonstrated on the first scenario. We extend our method to two other scenarios that are commonly seen in real-world applications. 

\textbf{Inpainting.}
Consider the case that $\xzero$ can be segmented into observed dimensions $\mask$ and unobserved dimensions $\umask$ such that we may write $\xzero=[\xzero_\mask, \xzero_\umask]$
and let $y=\xzero_\mask,$
and take $\plikelihood{y}{\xzero}=\delta_y(\xzero_\mask).$
The goal, then, is to sample
$\pmodel(\xzero \mid \xzero_\mask =y)= \pmodel(\xzero_\umask \mid \xzero_\mask)\delta_{y}(\xzero_\mask).$
Here we define the twisting function for each $t>0$ as 
\begin{align}\label{eqn:inpaint_twist}
    \twist(y\mid \xt, \mask):= \mathcal{N}\left(y; \denoiser (\xt)_\mask, t\sigma^2 \right),
\end{align}
and set twisted proposals and weights according to 
\cref{eq:diffusion-twisted-proposals,eq:diffusion-twisted-weigthing-functions}.
The variance in \cref{eqn:inpaint_twist}
is chosen as $t \sigma^2 =\mathrm{Var}_{\pmodel}[\xt \mid \xzero]$ for simplicity;
in general, choosing this variance to more closely match $\mathrm{Var}_{\pmodel}[y\mid \xt]$ may be preferable.
For $t=0,$ we define the twisting function analogously with small positive variance, 
for example as $\twist(y\mid \xzero)= \mathcal{N}(y; \xzero_\mask, \sigma^2).$
This choice of simplicity changes the final target slightly;
alternatively, the final twisting function, proposal, and weights may be chosen to maintain asymptotic exactness (see \Cref{sec:supp_method_inpaint_alt_final_step}).



\textbf{Inpainting with degrees of freedom.} 
We next consider the case when we wish to condition on some observed dimensions,
but have additional degrees of freedom.
%In particular, 
%let $\pmodel(\xzero \mid y) = \pmodel(\xzero \mid \exists \mask \in \mathcal{M} \ \mathrm{s.t.}\ \xzero_\mask =y),$
%where $\mathcal{M}$ is a set of possible observed dimensions.
For example in the context of motif-scaffolding in protein design,
we may wish to condition on a functional motif $y$ appearing \emph{anywhere} in a protein structure, rather than having a pre-specified set of indices $\mask$ in mind. 
To handle this situation, we
(i) let $\mathcal{M}$ be a set of possible observed dimensions,
(ii) express our ambivalence in which dimensions are observed as $y$ using randomness by placing a uniform prior on $p(\mask)= 1/|\mathcal{M}|$ for each $\mask \in \mathcal{M},$ and
(iii) again embed this new variable into our joint model to define the degrees of freedom likelihood by 
$p(y\mid \xzero)=\sum_{\mask \in \mathcal{M}} \pmodel(\mask, y \mid \xzero)=|\mathcal{M}|\inv \sum_{\mask \in \mathcal{M}} p(y\mid \xzero, \mask).$
Accordingly, we approximate $\pmodel(y\mid \xt)$ with the twisting function
\begin{align}\label{eqn:mask_dof_twist}
    \twist(y\mid \xt) := |\mathcal{M}|\inv \sum_{\mask \in \mathcal{M}} \twist(y\mid \xt, \mask),
\end{align}
with each $\twist(y\mid \xt,\mask)$ defined
as in \cref{eqn:inpaint_twist}.
Notably, \cref{eqn:mask_dof_twist,eqn:inpaint_twist} coincide when $|\mathcal{M}|=1.$

The sum in \cref{eqn:mask_dof_twist} may be computed efficiently because each term depends on $\xt$ only through the same denoising estimate $\denoiser(\xt),$
which must be computed only once.
Since computation of $\denoiser(\xt)$ is the expensive step, the overall run-time is not significantly increased by using even large $|\mathcal{M}|$.


\subsection{TDS on Riemannian manifolds}\label{sec:TDS_riemannian_main}
TDS extends to Riemannian diffusion models on with little modification.
Riemannian diffusion models \citep{deriemannian} are structured like in the Euclidean case, but with conditionals defined by tangent normal distributions parameterized with a score approximation followed by a manifold projection step (see e.g.\ \citep{chirikjian2014gaussian,deriemannian}). 
When we assume that (as, e.g., in \citep{yim2023se}) the model is associated with a denoising network $\denoiser$, twisting functions are also constructed analogously.
For conditional tasks defined by likelihoods, $\twist(y\mid \xt)$ in \cref{eq:diffusion-twisting-functions} applies.
 For inpainting (and by extension, degrees of freedom), we propose 
 \begin{align}\label{eqn:tangent_normal}
     \twist(y\mid \xt) = \mathcal{TN}_{\denoiser(\xt)_\mask}(y ; 0, t \sigma^2),
      \end{align}
%\begin{align}\label{eqn:likelihood_approx_riemannian} \twist(y\mid \xt) {=} \mathcal{TN}_{\denoiser^M(\xt)}(y ; 0, \bar \sigma_t^2),\end{align}
where $\mathcal{TN}_{\denoiser(\xt)_\mask}(0, t \sigma^2)$ is a tangent normal distribution centered on $\denoiser(\xt)_\mask.$
As in the Euclidean case, $\pTwist(\xt \mid \xtpo, y)$ is defined with 
conditional score approximation $s_\theta(\dt{x}{t},y)=s_\theta(\dt{x}{t}) + \nabla_{\dt{x}{t}} \log \tilde p(y \mid \dt{x}{t})$, which is computable by automatic differentiation.
\Cref{sec:riemannian} provides details.
%Though $\nabla_{\xt}\log \twist(y; \xt)$ is a $\mathcal{T}_{\xt}$-valued Riemannian gradient, it is nevertheless computable 

%\section{Riemannian \methodname{}}\label{sec:riemannian}

This section provides additional details on the extension of TDS to Riemannian diffusion models introduced in \Cref{sec:method}.
We first introduce the tangent normal distribution.
We then provide with background on Riemannian diffusion models, which we parameterize with the tangent normal.
Then we describe the extension of TDS to these models.
Finally we show how \Cref{alg:TDS} modifies to this setting.

\paragraph{The tangent normal distribution.}
Just as the Gaussian is the parametric family underlying Euclidean diffusion models in \Cref{eq:unconditional-diffusion-transition}, 
the tangent normal (see e.g.\ \citep{chirikjian2014gaussian,deriemannian})
underlies generation in Riemannian diffusion models so we review it here.

We take the tangent normal to be the distribution is implied by a two step procedure.
Given a variable $x$ in the manifold, the first step is to sample a variable $\bar y$ in $\mathcal{T}_{x},$ the tangent space at $x;$
if $\{h_1, \dots, h_D\}$ is an orthonormal basis of $\mathcal{T}_{x}$
one may generate
\begin{align}
    \bar y = \mu + \sum_{d=1}^D \sigma \epsilon_d \cdot h_d,
\end{align}
with $\epsilon_d\overset{i.i.d.}{\sim}\mathcal{N}(0, 1),$
and $\sigma^2>0$ a variance parameter. 
The second step is to project $\bar y$ back onto the manifold to obtain
$y = \exp_{x}\{ \bar y\}$
where $\exp_{x}\{ \cdot \}$ denotes the exponential map at $x.$ The resulting distribution of $y$ is denoted as $\mathcal{TN}_{x}(\mu, \sigma^2)$. 
By construction, $\mathcal{T}_x$ is a Euclidean space with its origin at $x,$
so when $\|\mu\|_2=0, \mathcal{TN}_{x}(\mu, \sigma^2)$ is centered on $x.$
And since the geometry of a Riemannian manifold is locally Euclidean, when $\sigma^2$ is small
the exponential map is close to the identity map and the tangent normal is essentially a narrow Gaussian distribution in the manifold at $x.$
Finally, we use $\mathcal{TN}_{x}(y;\mu, \sigma^2)$ to denote the density of the tangent normal evaluated at $y$.

Because this procedure involves a change of variables from $\bar y$ to $y$ (via the exponential map),
to compute the tangent normal density one computes 
\begin{equation}\label{eq:TN_density}
    \mathcal{TN}_{x}(y;\mu, \sigma^2) = \mathcal{N}(\exp^{-1}_x\{y\}; \mu,  \sigma^2) \left| \frac{\partial}{\partial y} \exp_x^{-1}\{y\}\right|
\end{equation}
where $\exp_x^{-1}\{y\}$ is the inverse of the exponential map (from the manifold into $\mathcal{T}_x$), 
and $\left| \frac{\partial}{\partial y} \exp_x^{-1}\{y\}\right|$ is the determinant of the Jacobian of the exponential map; when the manifold lives in a higher dimensional subset of $R,$ we take 
$\left| \frac{\partial}{\partial y} \exp_x^{-1}\{y\}\right|$ to be the product of the positive singular values of the Jacobian.

\paragraph{Riemannian diffusion models and the tangent normal distribution.}
Riemannian diffusion models proceeds through a geodesic random walk \citep{deriemannian}.
At each step $t,$
one first samples a variable $\dt{\bar x}{t}$ in $\mathcal{T}_{\xtpo};$
if $\{h_1, \dots, h_D\}$ is an orthonormal basis of $\mathcal{T}_{\xtpo}$
one may generate
\begin{align}
    \dt{\bar x}{t}  =\sigma_{t+1}^2 s_\theta(\dt{x}{t+1})+ \sum_{d=1}^D \sigma_{t+1} \epsilon_d \cdot h_d,
\end{align}
with $\epsilon_d\overset{i.i.d.}{\sim}\mathcal{N}(0, 1).$
One then projects $\dt{\bar x}{t}$ back onto the manifold to obtain
$\xt = \exp_{\xtpo}\{ \dt{\bar x}{t}\}$
where $\exp_{x}\{ \cdot \}$ denotes the exponential map at $x.$

This is equivalent to sampling $\xt$ from a tangent normal distribution as 
\begin{equation}\label{eqn:reverse_transition_unconditional_riemannian}
    \dt{x}{t} \sim p(\dt{ x}{t} \mid \xtpo) = \mathcal{TN}_{\xtpo}\left(\dt{ x}{t};\sigma_{t+1}^2 s_\theta(\dt{x}{t+1}),
\sigma_{t+1}^2 \right).
\end{equation}

\paragraph{TDS for Riemannian diffusion models.}
To extend TDS, appropriate analogues of the twisted proposals and weights are all that is needed.
For this extension we require that the diffusion model is also associated with a manifold-valued denoising estimate $\denoiser$ as will be the case when, for example, $\scoremodel(\xt, t) := \nabla_{\xt} \log q(\xt \mid \xzero=\denoiser)$ for $\denoiser = \denoiser(\xt, t).$
In contrast to the Euclidean case, a relationship between a denoising estimate and a computationally tractable score approximation may not always exist for arbitrary Riemannian manifolds;
however for Lie groups when the the forward diffusion is the Brownian motion, tractable score approximations do exist \citep[Proposition 3.2]{yim2023se}.

For the case of positive and differentiable $\plikelihood{y}{\xzero},$ we again choose twisting functions $\twist(y\mid \xt):= \plikelihood{y}{\denoiser(\xt)}.$

Next are the inpainting and inpainting with degrees of freedom cases.
Here, assume that $\xzero$ lives on a multidimensional manifold (e.g.\ $SE(3)^N$)
and the unmasked observation $y=\xzero_\mask$ with $\mask \subset \{1, \dots, N\}$ on a lower-dimensional sub-manifold (e.g.\ $SE(3)^{|\mask|},$ with $|\mask|<N$).
In this case, twisting functions are constructed exactly as in \Cref{subsec:conditioning-problems}, 
except with the normal density in \Cref{eqn:inpaint_twist} replaced with a Tangent normal as
\begin{align}
\twist(y\mid \xt) {=} \mathcal{TN}_{\denoiser(\xt)_\mask}( y ; 0, \bar \sigma_t^2).
\end{align}

For all cases, we propose the twisted proposal as
\begin{equation}\label{eqn:likelihood_approx_riemannian_supp}
    \pTwist(\dt{ x}{t} \mid \xtpo, y) {=} \mathcal{TN}_{\xtpo}\left(\dt{ x}{t};\sigma_{t+1}^2 s_\theta(\dt{x}{t+1}, y), \tilde\sigma_{t+1}^2 \right)
\end{equation}
%where $\bar y := \exp_{\denoiser(\dt{x}{t})_\umask}^{-1}(y)\in \mathcal{T}_{\denoiser(\dt{x}{t})_\umask}$ 
%is $\dt{x}{0}$ mapped into $\mathcal{T}_{\denoiser(\dt{x}{t})}$ by the inverse of the exponential map at $\denoiser(\dt{x}{t})_\umask,$
where as in the Euclidean case  $s_\theta(\dt{x}{t}, y)=s_\theta(\dt{x}{t}) + \nabla_{\dt{x}{t}} \log \tilde p(y \mid \dt{x}{t})$.


Weights at intermediate steps are computed as in the Euclidean case (\Cref{eq:diffusion-twisted-weigthing-functions}): 
%However, 
%since the tangent normal includes a projection step back to the manifold, its density involves not only usual Gaussian density but also the Jacobian of the exponential map to account for the change of variables.
%For example, in the case that the ambient dimension of the manifold and the tangent space are the same dimension,
%$$
%\mathcal{TN}_{x}(x^\prime;\mu, \sigma^2)= \mathcal{N}(\exp^{-1}_x\{x^\prime\}; \mu, \sigma^2)\left|\frac{\partial}{\partial x} \exp_x^{-1}\{x^\prime\}\right|,
%$$
%where $\exp^{-1}_x\{x^\prime\}$ is the inverse of the exponential map at $x$ from the manifold into $\mathcal{T}_x,$
%and 
%$\left|\frac{\partial}{\partial x} \exp_x^{-1}\{x^\prime\}\right|$ is the determinant of its Jacobian.
%The weights are then computed as
\begin{align}\label{eq:diffusion-twisted-weigthing-functions-riemannian}
    \weight_t (\xt , \xtpo ):&= \frac{\pmodel (\xt \mid \xtpo) \twist(y \mid \xt)}{\twist(y \mid \xtpo) \pTwist(\xt \mid \xtpo, y)} \\
    &= \frac{\mathcal{TN}_{\xtpo}(\xt; \sigma_{t+1}^2\scoremodel(\xtpo), \sigma_{t+1}^2) \twist(y \mid \xt)
    }{\twist(y \mid \xtpo) \mathcal{TN}_{\xtpo}(\xt; \sigma_{t+1}^2\scoremodel(\xtpo, y), \tilde\sigma_{t+1}^2)}.
\end{align}
While the proposal and target contribute identical Jacobian determinant terms that cancel out, they remain in the twisting functions.


\paragraph{Adapting the TDS algorithm to the Riemannian setting.}
To translate the TDS algorithm to the Riemannian setting we require only two changes.

The first is on \Cref{alg:TDS} 
\Cref{alg_line:conditional_score}.
Here we assume that the unconditional score is computed through the transition density function: 
\begin{align}
\scoremodel(\xtpo) := \nabla_{\xtpo} \log q_{t+1{|}0}(\xtpo\mid \denoiser) \quad \mathrm{\ for\ } \denoiser = \denoiser(\xtpo).
\end{align}
Note that the gradient above ignores the dependence of $\denoiser$ on $\xtpo$. 

%\BLT{Is it clear that the gradient above ignores the dependence of $\denoiser$ on $\xt$?}
The conditional score approximation in \Cref{alg_line:conditional_score} is then replaced with 
$$
\tilde s_k \leftarrow \pTwist (\cdot \mid \xtpo_k, y):= \nabla_{\xtpo_k} \log q_{t+1{|} 0}(\xtpo_k\mid \denoiser) + \nabla_{\xtpo_k} \log \tilde p_k^{t+1} \quad \mathrm{ \ for\  } \denoiser = \denoiser(\xtpo_k).
$$
Notably, $\tilde s_k$ is a $\mathcal{T}_{\xtpo_k}$-valued Riemannian gradient.

The second change is to make the proposal on \Cref{alg:TDS} \Cref{alg_line:proposal} a tangent normal, as defined in \cref{eqn:likelihood_approx_riemannian_supp}.


\section{Related work}\label{sec:related-works}
%\BLT{This section could be moved to the appendix if needed.}
There has been much recent work on conditional generation using diffusion models.
These prior works demand either task specific conditional training,
or involve unqualified approximations and can suffer from poor performance in practice.
We discuss additional related work in \Cref{sec:additional_related_work}.

\textbf{Gradient guidance.}
We use \emph{gradient guidance} to refer to a general approach that incorporates conditioning information with gradients through the neural network $\denoiser(\xt).$
For example, in an inpainting setting, \citet{hovideo} propose a Gaussian approximation $\mathcal{N}(y\mid \denoiser(\xt)^\mask, \alpha)$ to $\pmodel(y \mid \xt)$. This approximation motivates a modified transition distribution as in \Cref{eq:diffusion-twisted-proposals}, with the corresponding approximation to the conditional score as 
%corresponding to an approximation of the conditional score as
$
s_\theta(\xt,y ) = s_\theta(\xt) - \nabla_{\xt} \| y - \denoiser(\xt)_M\|^2/\alpha.
$
This approach coincides exactly with TDS applied to the inpainting task when a single particle is used.
Similar Gaussian approximations have been used by \citep{chung2022diffusion,song2023pseudoinverse} for other inverse problems.

Gradient guidance can also be used with non-Gaussian approximations;
e.g.\ using $\plikelihood{y}{\denoiser (\xt}) \approx \pmodel (y \mid \xt)$ for a given likelihood $\plikelihood{y}{\xzero}.$
This choice again recovers TDS with one particle. 

Empirically, these heuristics can have unreliable performance, e.g. in image inpainting problems \citep{zhang2023towards}.
By comparison, TDS enjoys the benefits of gradient guidance by using it in proposals, while also providing a mechanism to eliminate approximation error by simulating additional particles.
%A limitation of this class of methods is its sensitivity to required hyperparameters (often referred to as the \emph{guidance scale}) that influence the approximation which they sample. 

\textbf{Replacement method.}
The replacement method \citep{song2020score} is a method introduced for image inpainting using only unconditional diffusion models.
The idea is to replace the observed dimensions of intermediate samples $\xt_\mask$, with a noisy version of observation $\xzero_\mask$.
However, it is a heuristic approximation and can lead to inconsistency between inpainted region and observed region \citep{lugmayr2022repaint}. 
Additionally, the replacement method applies only to inpainting problems.
While recent work has extended the replacement method to linear inverse problems \citep[e.g.][]{kawar2022denoising},
the approach provides no accuracy guarantees. It is unclear how to extend these methods to arbitrary differentiable likelihoods.

\textbf{SMC samplers for diffusion models.}
Most closely related to the present work is SMC-Diff \citep{trippe2022diffusion},
which uses SMC to provide asymptotically accurate conditional samples for the inpainting problem.
However, this prior work 
(i) is limited to the inpainting case, and
(ii) provides asymptotic guarantees only under the assumption that the learned diffusion model exactly matches the forward noising process, which is rarely satisfied in practice.
Also, SMC-Diff does not leverage twisting functions.

In concurrent work, \citet{cardoso2023monte} propose MCGdiff, an alternative SMC algorithm that uses the framework of auxiliary particle filtering \citep{pitt1999filtering} to provide asymptotically exact conditional inference.  Compared to TDS, MCGdiff avoids computing gradients of the denoising network but applies only to linear inverse problems, with inpainting as a special case.
%In concurrent work, \citet{cardoso2023monte} propose an alternative SMC algorithm for conditional generation, MCGdiff. MCGdiff applies to linear inverse problems (with inpainting as a special case) and uses the framework of auxiliary particle filtering \citep{pitt1999filtering} to provide asymptotically exact inferences without relying on strong assumptions on the diffusion model.
%But compared to TDS, MCGdiff does not benefit from gradient guidance and (besides inpainting) does not apply to the conditional considered. While MCGdiff avoids computing gradients of the denoising network, TDS is not restricted to linear inverse problems and benefits from gradient guidance.



\section{Simulation study and conditional image generation}\label{sec:results}
\begin{figure*}[t]
    \tiny
    \includegraphics[scale=0.7]{sections/figures/toy/toy_A.png}
    \includegraphics[scale=0.7]{sections/figures/toy/toy_B.png}
    \includegraphics[scale=0.7]{sections/figures/toy/toy_C.png}
    \includegraphics[scale=0.6]{sections/figures/toy/toy_legend.png}
  \centering
    \caption{
   % Error on small-scale problems with ground truth.  Error of estimate of conditional mean (mean +- 2SEM) with 25 replicates. 
   Errors of conditional mean estimations with 2 SEM error bars averaged over 25 replicates. TDS applies to all three tasks and provides increasing accuracy with more particles.  
}
\label{fig:exp-toy}
\vspace{-1.5em}
\end{figure*}
We first test the dependence of the accuracy of TDS on the number of particles in synthetic settings with tractable exact conditionals in \Cref{subsec:exp-toy}.
\Cref{subsec:exp-mnist} compares TDS to alternative approaches on class-conditional image generation, and an image inpainting experiment is included in \Cref{supp:subsubsec:mnist-inpainting}. 
See \Cref{sec:supp_results} for all additional details.

Our evaluation includes:
(1) \textit{TDS}; (2) \textit{TDS-IS},  an importance sampler that uses TDS's proposal; %which can be viewed as TDS without resampling steps; 
(3) \textit{IS}, a naive importance sampler described in \Cref{subsec:SMC_framing};  and
(4) \textit{Gradient Guidance}. %, referred to the generalization of  prior works  \citep{chung2022diffusion,hovideo,song2023pseudoinverse}, and implemented as TDS with one particle (see \Cref{sec:related-works}).
For inpainting-type problems we further include: (5) \textit{Replacement method}; and (6) \textit{SMC-Diff}. %Additionally, We include a comparision to Classifier Guidance in high-dimensional image generation tasks in \Cref{supp:subsubec:inet-and-cifar}. 
\iffalse
\begin{enumerate}[label=(\roman*)] 
\item \emph{TDS}: twisted diffusion sampler; 
\item \emph{TDS-IS}: an importance sampler that uses twisted proposals and weight set to $\twist(y\mid \xzero)$. TDS-IS can be viewed as TDS without resampling steps;
\item \emph{Gradient guidance}: viewed as TDS with one particle. This approach coincides with the \emph{reconstruction guidance} \citep[e.g.][]{chung2022diffusion,hovideo,song2023pseudoinverse};
\item \emph{IS}: the naive importance sampler described in \Cref{sec:method}; 
\item \emph{Replacement}: a heuristic sampler for inpainting tasks that replaces $\xt_{\mask}$ with a noisy version of the observed $\xzero_{\mask}$ at each $t$ \citep{song2020score};  
\item \emph{SMC-Diff}: an SMC algorithm with replacement method as proposals for inpainting tasks \citep{trippe2022diffusion}.
TDS and SMC-Diff are implemented with systematic resampling.
\end{enumerate}
\fi 

Each SMC sampler forms an approximation to $\pmodel(\xzero \mid y)$ with $K$ weighted particles $(\sum_{k^\prime=1}^K w_{k^\prime})^{-1} \cdot \sum_{k=1}^K w_k \delta_{\xzero_k}$. 
Gradient guidance and replacement method are considered to form a similar particle-based approximation with $K$ independent samples viewed as $K$ particles with uniform weights. 

\paragraph{Compute Cost:} %
Compared to unconditional generation, TDS has compute cost that is 
(i) larger by a constant factor due to the need to backpropogate through the denoising network when computing the conditional score approximation and
(ii) linear in the number of particles.
As a result, the compute cost at inference cost is potentially large relative for accurate inference to be achieved.
%However 
%Extra (time and memory) cost from backprop, but memory cost can be reduced by use of gradient checkpointing \citep{chen2016training}.

By comparison, conditional training methods
provide fast inference by \emph{amortization} \citep{gershman2014amortized},
in which most computation is done ahead of time.
Hence they may be preferable for applications where sampling time computation is a primary concern.
On the other hand, TDS may be preferable when the likelihood criterion is readily available (e.g.\ through an existing classifier on clean data) but training an amortized model poses challenges;
for example, the labeled data, neural-network engineering expertise, and up-front compute resources required for amortization training can be prohibitive.

%

\subsection{Applicability and precision of TDS in two dimensional simulations}\label{subsec:exp-toy}
We explore two questions in this section:
(i) what sorts of conditioning information can be handled by TDS and other methods, and 
(ii) how does the precision of TDS depend on the number of particles?

To study these questions, we first consider an unconditional diffusion model $\pmodel(\xzero)$ approximation of a bivariate Gaussian.
For this choice, the marginals of the forward process are also Gaussian, and so we may define $\pmodel(\dt{x}{0:T})$ with an analytically tractable score function without neural network approximation.
Consequently, we can analyze the performance without the influence of score network approximation errors.
And the choice of a two-dimensional diffusion permits close approximation of exact conditional distributions by numerical integration that can then be used as ground truth.

We consider three test cases defining the conditional information: 
(1)  Smooth likelihood: $y$ is an observation of the Euclidean norm of $x$ with Laplace noise, with $\pmodel(y\mid \xzero)=\exp\{| \|\xzero \|_2 - y|\} / 2 .$  This likelihood is smooth almost everywhere.\footnote{This likelihood is smooth except at the point $x=(0, 0)$.}
    (2) Inpainting:  $y$ is an observation of  the first dimension of $\xzero$, with $\pmodel(y\mid \xzero, \mask=0) = \delta_{y}(\xzero_0)$.  
    (3) Inpainting with degrees-of-freedom: $y$ is a an observation of either the first or second dimension of $\xzero,$ with $\mathcal{M}=\{0, 1\}$ and $\pmodel(y \mid \xzero) = \frac{1}{2}[\delta_y(\xzero_0)+\delta_y(\xzero_1)].$. 
\iffalse
\begin{enumerate}
    \item {Smooth likelihood: $y$ is an observation of the Euclidean norm of $x$ with Laplace noise, with $\pmodel(y\mid \xzero)=\exp\{| \|\xzero \|_2 - y|\} / 2 .$  This likelihood is smooth almost everywhere.\footnote{This likelihood is smooth except at the point $x=(0, 0)$.} }
    \item{Inpainting:  $y$ is an observation of $\xzero$'s first dimension, with $\pmodel(y\mid \xzero, \mask=0) = \delta_{y}(\xzero_0).$}
    %\item{Inpainting:  $\pmodel(\xzero \mid y=1)$ where $y$ representing the first dimension of $\xzero$ is set to 1}
    %\item{Inpainting with degrees-of-freedom: $\pmodel(\xzero \mid y = 0 \textrm{ or } 1)$ where y -- first dimension of $\xzero$ -- can be either 0 or 1. }
    \item{Inpainting with degrees-of-freedom: $y$ is a an observation of either the first or second dimension of $\xzero,$ with $\mathcal{M}=\{0, 1\}$ and $\pmodel(y \mid \xzero) = \frac{1}{2}[\delta_y(\xzero_0)+\delta_y(\xzero_1)].$}
\end{enumerate}
\fi 
In all cases we fix $y=0$ and consider estimating $\E_{\pmodel}[\xzero\mid y].$


\Cref{fig:exp-toy} reports the estimation error for the mean of the  desired conditional distribution, i.e. $\|\sum w_k \xzero_k  - \E_q[\xzero \mid y ] \|_2$.
TDS provides a computational-statistical trade-off:
using more particles decreases mean square estimation error at the $O(1/K)$ parametric rate
(note the slopes of $-1$ in log-log scale) as expected from standard SMC theory \citep[Ch. 11]{chopin2020introduction}.
This convergence rate is shared by TDS, TDS-IS, and IS in the smooth likelihood case,
and by TDS, SMCDiff and in the inpainting case;
TDS-IS, IS and SMCDiff are applicable however only in these respective cases, 
whereas TDS applicable in all cases.
The only other method which applies to all three settings is Gradient Guidance, which exhibits significant estimation error and does not improve with many particles.




\subsection{Class-conditional image generation}\label{subsec:exp-mnist}
%\textbf{Main message}: Utility with a few particles

%We apply TDS to image inpainting and class-conditioned generation tasks on the MNIST dataset. For both tasks, we trained an unconditional diffusion model $\pmodel$  on 60k images with a modified architecture from a public codebase\footnote{\url{https://github.com/openai/guided-diffusion}}. See appendix for details.   

%\subsubsection{Class-conditioned generation}

%\begin{figure}
%  \centering
%  \begin{subfigure}[t]{0.9\linewidth}
%    \centering
%   \includegraphics[width=\textwidth]{sections/figures/mnist-class/classifier-guidance-sample-comparision-label7-P64-resampleP16.pdf}
%      \caption{Conditional samples of each method ran with 64 particles given class $y=7$.}
%      \label{subfig:mnist-class-sample}
%  \end{subfigure}\hfill 
%    \begin{subfigure}[t]{0.46\linewidth}
%    \centering
%    \includegraphics[width=\textwidth]{sections/figures/mnist-class/class-ess-legendTrue.pdf}
%      \caption{ESS trace avg. over 100 runs ($K=64$)}%for SMC-based methods averaged over 100 runs.}
%      \label{subfig:mnist-class-ess}
%  \end{subfigure}%\hfill 
%  \begin{subfigure}[t]{0.5\linewidth}
%    \centering
%    \includegraphics[width=\linewidth]{sections/figures/mnist-class/class-accuracy-legendTrue.pdf}
%      \caption{Classification accuracy avg. over 1k runs.}% Results are averaged over 1{,}000 runs.}
%      \label{subfig:mnist-class-accuracy}
%  \end{subfigure}
%  \caption{MNIST class-conditional generation. TDS  generates a diverse set of desired digits, and has higher particle efficiency compared to other SMC samplers.  Notably, TDS with $K=2$ particles outperform the Guidance baseline in terms of classification accuracy.}
%  \label{fig:mnist-class}
%\end{figure}




%Add a sentence to motivate this task. Neural network approximation. 
We next study the performance of TDS on diffusion models with neural network approximations to the score functions.
In particular, we study the class-conditional image generation task, which involves sampling an image from $\pmodel(\xzero \mid y) \propto \pmodel (\xzero) \plikelihood{y}{\xzero}$, where $\pmodel (\cdot)$ is a pretrained diffusion model on images $\xzero$, $y$ is a given image class, and $\plikelihood{y}{\cdot}$ is the classification likelihood. %\footnote{We trained an unconditional diffusion model  on 60{,}000 training images with 1{,}000 diffusion steps and the architecture proposed in  \citep{dhariwal2021diffusion}.The  classification likelihood $\plikelihood{y}{\cdot}$ is parameterized by a pretrained ResNet50 model taken from  \url{https://github.com/VSehwag/minimal-diffusion}.} 
To assess the faithfulness of generation, we evaluate  \emph{classification accuracy} on predictions of conditional samples  $\xzero$ given $y$, made by the same classifier that specifies the likelihood. 
In all experiments, we follow the standard practice of returning the denoising mean on the final sample \cite{ho2020denoising}.

On the MNIST dataset, we compare TDS to TDS-IS, Gradient Guidance, and IS. 
%The results are summarized in \Cref{fig:image-class-cond}. 
\Cref{subfig:mnist-class-sample}  compares the conditional samples of TDS and Gradient Guidance given class $y=7$.  
Samples from Gradient Guidance have noticeable artifacts, and most of them do not resemble the digit 7; by contrast, TDS produces authentic and correct digits.

\Cref{subfig:mnist-class-accuracy} presents an ablation study of the effect of \# of particles $K$ on MNIST. For all SMC samplers,  more particles improve the classification accuracy,  with $K=64$ leading to nearly perfect accuracy. The performance of Gradient Guidance is constant with respect to $K$ (in expectation). Notably, for fixed $K$, TDS and TDS-IS have comparable performance and outperform Gradient guidance and IS. 
%This observation suggests that one can use TDS-truncate %to avoid effective sample size drop and to encourage the generation diversity, without compromising the performance.  

We next apply TDS to higher dimension datasets.  \Cref{subfig:inet-tds-ind-samples} shows  samples from TDS $(K=16)$ using a pre-trained diffusion model  and a pretrained classifier on the ImageNet dataset ($256\times 256 \times 3$ dimensions). These samples are qualitatively good and capture the class label. \Cref{supp:subsubec:inet-and-cifar} provides more samples, comparision to Classifier Guidance \citep{dhariwal2021diffusion}, and results for the CIFAR-10 dataset. 

TDS can be extended by exponentiating twisting functions with a \emph{twist scale}. This extension is related to the existing literature of Classifier Guidance   that considers re-scaling the
gradient of the log classification probability. 
See \Cref{supp:subsubsec:mnist-class} for details and ablation study. 


\begin{figure*}[t!]
  \centering
  \begin{subfigure}[t]{0.38\linewidth}
    \centering
   \includegraphics[width=\textwidth]{sections/figures/mnist-class/two-cols-classifier-guidance-sample-comparision-label7-P64-resampleP16-showmethodnameTrue.pdf}
     % \caption{Approximate conditional samples of each method with 64 particles given class $y=7$.}
      \caption{MNIST: samples by TDS ($K=64$) and Gradient Guidance   given class `7'. TDS samples are randomly selected out of 64 particles in a single SMC run.}
      \label{subfig:mnist-class-sample}
  \end{subfigure}
\hfill
  \begin{subfigure}[t]{0.34\linewidth}
    \centering\includegraphics[width=\linewidth]{sections/figures/mnist-class/class-accuracy-legendTrue.pdf}
      \caption{MNIST: classification accuracy vs.\ $K$. Results are averaged over 1{,}000 runs with error bars denoting 2 standard errors.}
      \label{subfig:mnist-class-accuracy}
  \end{subfigure}
  \hfill
    \begin{subfigure}[t]{0.22\linewidth}
    \centering\includegraphics[width=\linewidth]{sections/figures/inet-and-cifar10/inet_TDS_independent_samples_N4.pdf}
      \caption{ImageNet: samples by TDS $(K=16)$ given class `brambling' from 4 independent SMC runs.}% Results are averaged over 1{,}000 runs.}
      \label{subfig:inet-tds-ind-samples}
  \end{subfigure}
\caption{Image class-conditional generation task.}
\vspace{-1.5em}
%  \caption{MNIST class-conditional generation. TDS  generates a diverse set of desired digits, and has higher particle efficiency compared to other SMC samplers.  Notably, TDS with $K=2$ particles outperform the Guidance baseline in terms of classification accuracy.}
  \label{fig:image-class-cond}
\end{figure*}

%Figure 3 shows more samples given randomly selected classes. We also reported results of the unconditional model from [1] that are evaluated on 50k samples with 250 sampling steps (Table 1). TDS and CG provide similar classification accuracy. TDS has similar FIDs compared to the unconditional model and better inception score. CG’s FID and inception score are better than TDS. We suspect this difference is attributed to the sample correlation (and hence less diversity) within particles in a single run of TDS(P=16).
%TDS can be applied to higher dimensional datasets, e.g. CIFAR10 and ImangeNet 256. 




%In addition, we include an MNIST inpainting experiment in \Cref{sec:supp_results}. We observe that TDS consistently outperforms TDS-IS, Guidance, IS, Replacement and SMC-Diff across $K$. We also encounter the ESS drop in the final 10 steps, and we likewise employ a truncation procedure there.% to alleviate the issue. 

\iffalse 

\subsubsection{Inpainting} 
\begin{figure}[t]
  \centering
  \begin{subfigure}[t]{0.28\linewidth}
    \centering
    \includegraphics[width=\textwidth]{sections/figures/mnist-inpainting/inpainting-quarter-ess.pdf}
      \caption{Effective sample size v.s. $t$}
  \end{subfigure}\hfill 
    \begin{subfigure}[t]{0.72\linewidth}
    \centering
    \includegraphics[width=\textwidth]{sections/figures/mnist-inpainting/inpainting-quarter-accuracy.pdf}
      \caption{Classification and Bayes Accuracy v.s. $K$}
  \end{subfigure}
  \caption{MNIST inpainting task}
  \label{fig:MNIST-inpaiting}
\end{figure}


The goal of inpainting is to sample from $\pmodel \left(\xzero \mid y=f(\xzero)\right)$ where $ f(\xzero) := \xzero_M$ denotes an indexing function with $M$ indices of observed pixels, and $y$ is the observed part of an image.  

We consider two accuracy criteria through the use of a separately pre-trained classifier as mentioned above. We know the true label $\dt{z}{0}$ for $\xzero$. 
\begin{itemize}
    \item Bayes accuracy $:=$
    %\begin{align}
     $   \mathds{1} \left\{  \left(\argmax_{z=1:C} \int \pmodelapprox (\xzero \mid y) p(z \mid \xzero)  d\xzero \right) = \dt{z}{0} \right\} $
    %\end{align}
    \item Weighted accuracy $:=$
    $  \int \pmodelapprox (\xzero \mid y) \mathds{1} \{ \argmax_{z=1:C} p(z \mid \xzero) = \truelabel \} d\xzero$
    
\end{itemize}

%\textbf{Bayes-optimal classification under partial information.}
%Consider the task assigning to an image of a digit $\xzero$ a classification $z$ of as one of the 10 digits, 1 through 10, when only a part of the image, $y,$ has been observed. For example, we consider $y$ to be the left half of an image and assume $z$ is represented as a one-hot vector. Provided that both (1) the distribution over images $\pmodel(\xzero),$ and (2) the true labeling mechanism $p(z\mid \xzero)$ are known, the Bayes optimal classifier (with respect to the zero-one loss) is
\begin{align}\label{eqn:bayes_opt_classifier}
\begin{split}
    \hat z^*(y) &= \argmax_{\ell} p(z_\ell=1 \mid y) \\
    &= \argmax_{\ell} \int \pmodel(\xzero \mid y) p(z_\ell=1 \mid \xzero )d\xzero,
\end{split}
\end{align}
where $p(z_\ell=1 \mid \xzero)$ is given by the optimal classifier on full images.

\Cref{eqn:bayes_opt_classifier} suggests that, if samples $\dt{x_m}{0}\sim\pmodel(\xzero \mid y)$ are available, we can use a Monte Carlo estimate of $\hat z^*(y)$ as
as 
$$
    \hat z^*(y) \approx M\inv \sum_{m=1}^M p(z \mid \dt{x_m}{0}),
$$
to approximate the optimal Bayes rule.

This observation provides a grounds to evaluate the accuracy of approximations to $p(\xzero \mid y)$ in a non-trivial setting;
approximation error in $p(\xzero \mid y)$ can only decrease the accuracy of the resulting classifier.
\textbf{Evaluation set-up and details:}
To examine the impact of number particles used on performance,
and the interaction between this knob and the quality of the score approximation (and there more twisting functions),
we consider classification of MNIST digits under partial information.
A diffusion model was fit to a training set of $50,000$ images.
Then on a test set of $10,000$ labeled test images we compared the performance of several inpainting algorithms on the following evaluation procedure, applied for on each test image:
\begin{enumerate}
    \item{Mask the right side of the image,}
    \item{Run the inpainting algorithm to sample completions of the image,}
    \item{Run the classifier of the completed images, and}
    \item{Make a prediction for the image as the label highest average predicted probability.}
\end{enumerate}
We compare the following inpainting algorithms on the basis of mean classification accuracy:
\begin{enumerate}
    \item{Twisted SMC (K=1 corresponds to reconstruction guidance)}
    \item{Naive SMC (importance sampling)}
    \item{SMCDiff}
    \item{Twisted SMC without resampling (SIS)}
    \item{The replacement method.}
    \item{Baseline Bayes' optimal using KDE?}
\end{enumerate}


We next asked if there was particle collapse, by looking the effective sample size per step.
This helped to pick between several choices of scaling on reconstruction guidance too.
\begin{itemize}
    \item{Twisted SMC with different variance schedules}
    \item{Twisted SMC with different \#'s of steps?}
    \item{SMCDiff}
\end{itemize}


\fi 

\section{Case study in computational protein design: the motif-scaffolding problem}\label{sec:exp-protein}

The biochemical functions of proteins are typically imparted by a small number of atoms, known as a \emph{motif}, that are stabilized by the overall protein structure, known as the \emph{scaffold} \citep{wang2022scaffolding}.
A central task in protein design is to identify stabilizing scaffolds in response to motifs expected to confer function.
We here describe an application of TDS to this task, and compare TDS to the state-of-the-art conditionally-trained model, RFdiffusion. See additional details in \Cref{sec:supp_results_protein}.
%We apply TDS to this task, and compare to the state-of-the-art conditionally-trained model, RFdiffusion.


Given a generative model supported on designable protein structures $\pmodel(\xzero)$,
suitable scaffolds may be constructed by solving a conditional generative modeling problem \citep{trippe2022diffusion}.
Complete structures are first segmented into a motif $\xzero_\mask$ and a scaffold $\xzero_\umask,$ i.e.\@ $\xzero = [\xzero_\mask, \xzero_{\umask}].$ 
Putative compatible scaffolds are then identified by (approximately) sampling from $\pmodel(\xzero_\umask \mid \xzero_\mask)
$ \citep{trippe2022diffusion}.

While the conditional generative modeling approach to motif-scaffolding has produced functional, experimentally validated structures for certain motifs \citep{watson2022broadly}, the general problem remains open.
Moreover, current methods for motif-scaffolding require one to specify the location of the motif within the primary sequence of the full scaffold;
this choice can require expert knowledge and trial and error.

We hypothesized that improved motif-scaffolding could be achieved through accurate conditional sampling.
To this end, we applied TDS to FrameDiff, a Riemannian diffusion model that parameterizes protein backbones as a collection of $N$ rigid bodies 
 (known as residues) in the manifold $SE(3)^N$ \citep{yim2023se}.\footnote{\url{https://github.com/jasonkyuyim/se3_diffusion}}
Each of the $N$ elements of $SE(3)$ consists of a rotation matrix and a translation that parameterize the locations of the backbone atoms of each residue.

\textbf{Likelihood, twisting, and degrees of freedom.}
The basis of our approach to the motif scaffolding problem is analogous to the inpainting case described in \Cref{subsec:conditioning-problems}.
We let $\plikelihood{y}{\xzero}=\delta_y(\xzero_\mask),$
where $y\in SE(3)^M$ describes the coordinates of backbone atoms of an $M=|\mathcal{M}|$ residue motif.
As such, to define twisting function we adopt the Riemannian TDS formulation described in  \Cref{sec:TDS_riemannian_main}.


To eliminate the requirement that the placement of the motif within the scaffold be pre-specified, we incorporate the motif placement as a degree of freedom.
We 
(i) treat the indices of the motif within the chain as a mask $\mask,$ 
(ii) let $\mathcal{M}$ be a set of possible masks of size equal to the length of the motif, and 
(iii) apply \Cref{eqn:mask_dof_twist} to average over these possible masks.

For some scaffolding problems, it is known that all motif residues must appear with contiguous indices.
In this case we choose $\mathcal{M}$ to be the set of all possible contiguous masks.
However, when motif residues are only known to appear in two or more possibly discontiguous blocks,
the number of possible placements can be too large and we choose $\mathcal{M}$ by
randomly sampling at most some maximum number (\texttt{\# Motif Locs.}) of masks.


%Prior motif-scaffolding approaches have required of pre-specification of the location of the motif within the final scaffold,
%translation and rotation of the motif \citep{wang2022scaffolding,trippe2022diffusion,watson2022broadly}.
%However, these nuisance degrees of freedom are ancillary to the biochemical function of a motif, which is determined by its internal geometry.
%TDS can eliminate this degree of freedom and condition on the motif appearing at \emph{any} combination of indices 
%by applying the strategy described in \Cref{subsec:conditioning-problems} (see \Cref{sec:supp_results_protein} for details).

We similarly eliminate the global translation and rotation of the motif as a degree of freedom.
The motivation is that restricting to one possible pose of the motif narrows the conditional distribution, thereby making inference  more challenging.
For translation, we use a likelihood that is invariant to the motif's center-of-mass and placement in the final scaffold by choosing $\plikelihood{y}{x} = \delta_{Py}(P x_\mask)$ where $P$ is a projection matrix that removes the center of mass of a vector \citep[see e.g.][section 3.3]{yim2023se}.
For rotation, we average the likelihood across some number (\texttt{\# Motif Rots.}) of possible rotations.



%\subsection{Sensitivity to number of particles, degrees of freedom and twist scale}\label{sec:sensitivity_5IUS}
\begin{figure}
  \centering
  \begin{subfigure}[c]{0.63\linewidth}
  \centering
    \includegraphics[width=1.0\textwidth]{sections/figures/protein/success_rate_5IUS}
    \vspace{-1.5em}
      \caption{TDS motif-scaffolding success rate (test case \texttt{5IUS}) 
      improves with more particles, and degrees of freedom, and twist-scale.}
      %on (Left) number of particles, (Middle Left) number of motif location degrees of freedom, (Middle Right) number of rotation degrees of freedom, and twist scale.  Using more particles and allowing more degrees of freedom increases success rates.}
%    \includegraphics[width=.9\textwidth]{sections/figures/protein/benchmark_success_rates}
  \label{fig:5IUS_results}
  \end{subfigure}
  \hfill 
%  \centering
    \begin{subfigure}[c]{0.32\linewidth}
    {
   %\tiny
   %\footnotesize
   \fontsize{7pt}{8.6pt}\selectfont
       \renewcommand{\arraystretch}{1.5} % increase row height by 50%
            \begin{tabular}{|lcc|}
    \hline
    Scaffold & TDS \& & RF \\ 
    size & FrameDiff & diffusion \\ %\hline
    <100 res. & 9 & 3 \\ %\hline
    $\ge$ 100 res. & 2 & 8 \\ 
    \hline
    Overall & 11 & 11 \\ \hline
    \end{tabular}
    }
  \caption{\# problems with higher success rate}
    \label{tab:motif_summary}
  \end{subfigure}
  \caption{Protein motif-scaffolding case study results}
    \label{fig:protein-results}
    \vspace{-1.5em}
  \end{figure}
  
\textbf{Ablation study.}
We first examine the impact of several parameters of TDS on
success rate in an in silico \emph{self-consistency} evaluation \citep{trippe2022diffusion}.
We begin with single problem (\texttt{5IUS}) in the 
benchmark set introduced by \citep{watson2022broadly},
before testing on the full set.
\Cref{fig:5IUS_results} (Left) shows that success rate increases monotonically with the number of particles.
Non-zero success rates in this setting required accounting for multiple motif locations; \Cref{fig:5IUS_results} (Left) uses $1{,}000$ possible motif locations and $100$ rotations.

We tested the impact of the degrees of freedom by evaluating the success rate of TDS (K=1) with increasing motif locations and 100 rotations (\Cref{fig:5IUS_results} Center Left), 
and increasing rotations and 1{,}000 locations (\Cref{fig:5IUS_results} Center Right).
The success rate was 0\% without accounting for either degree of freedom, and increased with larger numbers of locations and rotations.
We also explored including a heuristic twist scale as considered for image tasks (\Cref{subsec:exp-mnist});
in this case, the twist scale is a multiplicative factor on the logarithm of the twisting function.
 \Cref{fig:5IUS_results} (Right) shows larger twist scales gave higher success rates on this test case, where we use 8 particles, 1,000 possible motif locations and 100 rotations.
However, this trend is not monotonic for all problems (\Cref{fig:prot-twist-scale}).
%We use an effective sample size threshold of $K/2$ in all cases, that is, we trigger resampling steps only when the effective sample size falls below this level.


%\subsection{TDS on motif-scaffolding benchmark set}\label{sec:motif_benchmark}
\textbf{Evaluation on full benchmark.}
We next evaluate on TDS on a benchmark set of 24 motif-scaffolding problems \citep{watson2022broadly} and compare to the previous state of the art, RFdiffusion.
RFdiffusion operates on the same rigid body representation of protein backbones as FrameDiff.
TDS is run with K=8, twist scale=2, and 
$100$ rotations and $1{,}000$ motif location degrees of freedom ($100{,}000$ combinations total). 

Overall, TDS (applied to FrameDiff) and RFdiffusion have comparable performance 
(\Cref{tab:motif_summary});
each provides a success rate higher than the other in 11/24 cases; on two problems both methods have a  0\% success rate (full results in \Cref{fig:full_benchmark}).
This performance is obtained despite the fact that FrameDiff, unlike RFdiffusion, is not trained to perform motif scaffolding.
The division between problems on which each method performs well is primarily explained by total scaffold length, with TDS providing higher success rates on smaller scaffolds.

We suspect the shift in performance with scaffold length owes properties of the underlying diffusion models.
First, long backbones generated unconditionally by RFdiffusion are designable with higher frequency than those generated by FrameDiff \citep{yim2023se}.
Second, unlike RFdiffusion, FrameDiff can not condition on the fixed motif sequence.



\section{Discussion}\label{sec:discussion}
\ifnum \workshopsubmission=0
We propose TDS, a practical and asymptotically exact conditional sampling algorithm for diffusion models. 
We compare TDS to other approaches and demonstrate the effectiveness and flexibility of TDS on image class conditional generation and inpainting tasks. 
On protein motif-scaffolding problems with short scaffolds, TDS outperforms the current (conditionally trained) state of the art. 
\fi

A limitation of TDS is its requirement for  additional computes to simulate multiple particles.
While we observe improved performance with just two particles in some cases, the optimal number is problem dependent.
Moreover, the computational efficiency depends on how closely the twisting functions approximate exact conditionals,
which depends on the unconditional model and conditioning information.
Lastly, choosing twisting functions for generic constraints may be challenging.
Our future work will focus on addressing these limitations and improving the computational efficiency. 

%Future works include extensions 
%to large-scale image models and improvements on protein applications. In particular, the latter involves experimenting with a stronger unconditional protein model, and exploring the use of conditionally trained diffusion models as effective proposals for the SMC procedure. 
%Furthermore, there is a potential for exploring TDS's applications to compositional conditional tasks. 

\section*{Acknowledgements}
We thank Hai-Dang Dau, Arnaud Doucet, Joe Watson, and David Juergens for helpful discussion,
Jason Yim for discussions and code,
and David Baker for additional guidance.
%\newpage

\bibliographystyle{plainnat}

\bibliography{references}
\newpage




\newpage 
\appendix
\section*{Appendix}
%\setcounter{table}{0}
\addcontentsline{toc}{section}{Appendix}
\startcontents[appendixtoc]
\printcontents[appendixtoc]{l}{1}{\setcounter{tocdepth}{1}}

%\addcontentsline{toc}{section}{Appendix} % Add the appendix text to the document TOC
%\part{Appendix} % Start the appendix part
%\parttoc % Insert the appendix TOC





% Related Works 


% Method 
%% theorem, algorithm, modifications with DDPM structure etc

% Experiments 




% Change algorithm counter representation to uppercase letters
%\renewcommand{\thealgorithm}{\Alph{algorithm}}
\renewcommand{\thefigure}{\Alph{figure}}

%\section{Sequential Monte Carlo details}\label{sec:SMC_details}


% Original text: 
%However, for SMC to be practical the weights must be regular. Othewise when only a small fraction of particles receive nontrivial weights, the rest will be discarded and the procedure provides a poor approximation of the target distributions. See \Cref{sec:SMC_details} for details.  



%The number of particles needed for accurate estimation of each $\target_t$ is exponential in the KL divergence between successive targets \citep{chatterjee2018sample}: \begin{align}\label{eq:kl-intermediate-targets}\mathrm{KL}(\target_{t-1} \| \target_t)=\ln \E_{\target_{t-1}}[\weight_{t-1}\inv(\xt, \xtpo)] - \E_{\target_{t-1}}[\ln \weight_{t-1}\inv(\xt, \xtpo)].\end{align}
%This divergence will be large when a small fraction of particles receive a disproportionately large weight. Otherwise only a small fraction of the particles will receive nontrivial weights, and the rest will be discarded. In that case most particles are discarded, the procedure provides only small effective number of samples and therefore a poor approximation of the target distribution. 


Each step of SMC involves importance sampling of a target $\target_{t-1}$ with $\target_t$ as the proposal, and so the procedure inherits the operating characteristics of importance sampling. Notably this includes an exponential dependence of the number of samples needed for accurate estimation of $\target_{t-1}$ on the Kullback-Leibler divergence of each $\target_t$ from $\target_{t-1}$ \citep{chatterjee2018sample}. In the context of  \Cref{eq:final-target}, we see that $$\mathrm{KL}(\target_{t-1} \| \target_t)=\ln \E_{\target_{t-1}}[\weight_{t-1}\inv(\xt, \xtpo)] - \E_{\target_{t-1}}[\ln \weight_{t-1}\inv(\xt, \xtpo)].$$ This divergence will be large when a small fraction of particles receive a disproportionately large weight.
In that case, only a small fraction of the particles receive nontrivial weights, and the rest will be discarded. Consequently, the procedure provides only small effective number of samples and therefore a poor approximation of the target distribution. 




\iffalse 
Sequential Monte Carlo (SMC) is a set of techniques for inference in time-series models \citep{doucet2001sequential,naesseth2019elements,chopin2020introduction}.
%Feynman-Kac models provide a language for constructing a sequence of probability distributions concluding at a probability distribution of interest in a manner that is amenable to computational approximations by particle filtering. 
Any SMC algorithm may be described as an instance of a generic algorithm for inference in a Feynman-Kac (FK) model
%We review Feynman-Kac models (following \citet{chopin2020introduction}) and present the generic particle filtering algorithm 
\citep[Algorithm 10.1]{chopin2020introduction}.\footnote{Our presentation differs from that of \citet{chopin2020introduction} only in that it is stated in reverse time in order to mirror the time reversal in the diffusion model.}
%We then state a theorem on its asymptotic accuracy.
%Together, this theorem and algorithm make clear the requirements for a Feynman-Kac model to be used for asymptotically accurate conditional sampling.
This algorithm requires three ingredients.
First are an \emph{initial distribution} $M_T(\dt{x}{T})$ and \emph{transition kernels} $M_t(\dt{x}{t}\mid \dt{x}{t+1})$ for $t=T-1, \dots, 0$
that define proposal distributions for an ensemble of \emph{particles},  $\{ x^{(t)}_n\}_{n=1}^N$.
Second are positive real valued \emph{potential functions} $G_T(\dt{x}{T})$ and $G_t(\dt{x}{t}, \dt{x}{t+1})$ for $t=T-1, \dots, 0$ that define weights for the particles.
A \emph{sequence of Feynman-Kac models} is then defined as the sequence of distributions across $t,$
\begin{equation}\label{eqn:feynman_kac_sequence}
    \Q_t(\dt{x}{t:T})=L_t\inv 
        \left\{M_T(d\dt{x}{T}) \prod_{s=t}^{T-1} M_s(\dt{x}{s}\mid  \dt{x}{s+1})\right\}
        \left\{G_T(\dt{x}{T})\prod_{s=t}^{T-1} G_s(\dt{x}{s}, \dt{x}{s+1})\right\}
    ,
\end{equation}
where $L_t\inv$ is a normalizing constant.%,
\begin{equation}
L_t =  \int_{\dt{x}{0:t}} \left\{ G_T()\prod_{s=T-1}^t G_s(\dt{x}{s+1}, \dt{x}{s})\right\}
\left\{M_T(d\dt{x}{T}) \prod_{s=T-1}^t M_s(\dt{x}{s+1}, d \dt{x}{s})\right\}.
\end{equation}

The third ingredient is a resampling function (e.g.\ systematic resampling).
The generic particle filtering algorithm with $N$ particles is given by \cref{alg:FeymanKac_PF}.
The defining property that connects \Cref{alg:FeymanKac_PF} and the sequence of Feynman-Kac models is given by \Cref{theorem:FeynmanKac_PF}.
It provides that in the limit that many particles are used, 
the output of \Cref{alg:FeymanKac_PF} permits arbitrarily accurate estimates of expectations under each $\Q_t.$


\begin{algorithm}
\caption{Sequential Monte Carlo \citep[Algorithm 10.1]{chopin2020introduction}}\label{alg:FeymanKac_PF}
\begin{algorithmic}
\LineComment Input: $N (\#\mathrm{particles}),\{M_t \}_{t=0}^{T}$ and $ \{G_t \}_{t=0}^{T}$
\LineComment Operations involving index $n$ are performed for each particle $n = 1,\dots, N.$
\State $\dt{x_n}{T} \sim M_T(\dt{x}{T})$
\State $\dt{w_n}{T} \gets G_T(\dt{x}{T})$
\State $\dt{W_n}{T} \gets\dt{w_n}{T} / \sum_{m=1}^N \dt{w_m}{T}$
\State $\dt{A_{1:N}}{T} \sim \texttt{resample}(\dt{W_n}{T})$
\For{$t=T-1, \dots, 0$} 
    %\State $\dt{A_{1:N}}{t} \sim \texttt{resample}(\dt{W_n}{t+1})$ \Comment{User-defined resampling}
    \State $\dt{x_n}{t} \sim M_t(\dt{x_{\dt{A_n}{t+1}}}{t}\mid \dt{x}{t+1})$
    \State $\dt{w_n}{t} \gets G_t(\dt{x_{\dt{A_n}{t+1}}}{t+1}, \dt{x_n}{t})$
    \State $\dt{W_n}{t} \gets\dt{w_n}{t} / \sum_{m=1}^N \dt{w_m}{t}$
    \State $\dt{A_{1:N}}{t} \sim \texttt{resample}(\dt{W_n}{t})$
\EndFor 
\end{algorithmic}
\end{algorithm}
\begin{theorem}[Almost sure convergence of estimates {{\citep[Proposition 11.4]{chopin2020introduction}}}]\label{theorem:FeynmanKac_PF}
Assume the potential functions $G_t$ are bounded and multinomial resampling is used,
for every continuous and bounded function $\phi,$
$$
\sum_{n=1}^N \dt{W_n}{t}\phi(\dt{x_n}{t}) \overset{a.s.}{\rightarrow} \int_{\dt{x}{t}} \phi(\dt{x}{t}) \Q_t(\dt{x}{t})d\dt{x}{t},
$$
where $\overset{a.s.}{\rightarrow}$ denotes almost sure convergence in the limit that $N$ tends to infinity.
\end{theorem}
\fi 
%
\section{Naive SMC: importance sampling from unconditional diffusion models}
We here provide details of the naive procedure, showing that it is a case of self-normalized importance sampling.
As mentioned in \Cref{subsec:SMC_framing} importance sampling is a special case of of SMC with a single step (or equivalently several steps, but no resampling at intermediate steps).
It involves a target $\nu$,
and proposal $r$. 
$K$ samples $x_1,\dots, x_K$ are drawn i.i.d.\ from $r$,
and an estimate of $v$ is computed as 
$\hat \nu = \sum w_k x_k$ with each $w_k = \tilde w_k / \sum_k \tilde w_k$
for $\tilde w_k = \nu(x_k)/r(x_k).$
Note hat that the procedure in \Cref{subsec:SMC_framing} is a special case of importance sampling with proposal $r(\dt{x}{0:T})=\pmodel(\dt{x}{0:T})$ and target $\pmodel(\dt{x}{0:T} \mid y)$.

%As a step building up to a practical algorithm, we consider an instance that corresponds to importance sampling using (unconditional) diffusion model as the proposal. 
%As an initial proof-of-concept, we first present a simple instance that wraps the unconditional diffusion generation process into an SMC sampler, defined by the following proposals and weighting functions. 
%\begin{fkModel}[Naive importance sampling]\label{sampler:naive}
%\begin{align*}
%    M_T(\xT)&=p(\xT)\  \mathrm{and} \ 
%  M_t(\xt \mid \xtpo)=\pmodel (\xt \mid \xtpo )\  \mathrm{ and\  for }\ \ T>t\ge 0,
%%    &\mathrm{ and }\ \  M_0( \dt{x}{0}\mid \dt{x}{1}) = \pmodel(\dt{x}{0} \mid \dt{x}{1}, y)
%\end{align*}
%and 
%\begin{align*}
%    &G_T(\xT)=G_t(\xt, \xtpo)=1 \ \ \mathrm{for }\ \ T>t>1,\ \mathrm{and }\ G_0(\xzero, \dt{x}{1})=p(y\mid \xzero).
%\end{align*}
%\end{fkModel}

%This sampler relies on the tractable unconditional sampling mechanism of $\pmodel$ in every transition kernel, i.e. not conditioning on measurement $y$. 
%And the weighting function is uninformative until the final step, where each particle of $\xzero$ is evaluated under the measurement model. 
%And the potential function used to weight the particles is equal to one at all but the final step.
%The final proposal $M_0(\xzero \mid \dt{x}{1})$ is analytically tractable to sample from for many types of complex conditioning information.
%otably by Bayes' rule we can see that $M_0(\xzero \mid \dt{x}{1})= \pmodel(\xzero \mid \dt{x}{1})\pmodel(y \mid \xzero) / \int \pmodel(\dt{\tilde x}{0} \mid \dt{x}{1})\pmodel(y \mid \dt{x}{1})d\dt{\tilde x}{0},$
%here the first two terms are computable and the denominator an integral taken w.r.t. a Gaussian and so will often be easily and accurately approximated by e.g.\ Gaussian quadrature or importance sampling.\footnote{
%   \BLT{Clean this bit up... for figure out to to weave into assumptions of propositions and leave as a remark below defending the assumption.}}

\section{\methodname{} additional details}\label{sec:supp_method}
In this section we provide additional details on TDS.
\Cref{sec:VP_VE_TDS} describes its generalization to alternative common formulations of diffusion models.
\Cref{sec:TDS_hyperparams} describes the choices of proposal variance and resampling strategy.
\Cref{sec:twisted_intermediates_derivation} provides a derivation of \Cref{eq:twisted-intermediates}.
\Cref{sec:supp_method_inpaint_alt_final_step} describes modifications to the final step of TDS in inpainting and inpainting with degrees of freedom applications.
And \Cref{sec:theorem_statement_and_proof} provides a full statement and proof of \Cref{prop:twisted_SMC}.


\subsection{Diffusion models with non-constant noise variance and variance-preserving scaling}\label{sec:VP_VE_TDS}
TDS as developed in \Cref{sec:method} assumed noise with constant variance $\sigma^2$ added at each timestep.
This choice corresponds to a particular discretization of a ``variance exploding'' (VE) diffusion model \citep{song2020score}.
Here we describe generalizations of TDS to non-constant variance schedules,
and to variance preserving (VP) diffusion models \citep{ho2020denoising}.
We adopt the VP formulation in our experiments;
though it introduces additional notation, it is known to work well in practice \citep{song2020score}.

\paragraph{TDS algorithm for variance exploding diffusion models.}
VE diffusion models define the forward process $q$ by 
\begin{align}\label{eq:VE-forward-q}
    q(\dt{x}{1:T} \mid \xzero) := \prod_{t=1}^T q(\xt \mid \xprev), \qquad q(\xt \mid \xprev) := \mathcal{N}(\xt; \xprev, \sigma_t^2). 
\end{align}
where $\sigma_t^2$ is an increasing sequence of variances such that $q(\xT) \approx \mathcal{N}(0,\bar \sigma_T^2)$, with  $\bar \sigma_t^2 :=\sum_{t^\prime=1}^t \sigma_{t^\prime}^2$ for $t=1,\cdots, T$.  And so one can set $p(\xT) =  \mathcal{N}(0,\bar \sigma_T^2)$ to match $q(\xT)$.
Notably, \Cref{eq:VE-forward-q} implies the conditional $q(\xt \mid \xzero) = \mathcal{N}(\xt; \xzero, \bar\sigma_t^2)$. 

The reverse diffusion process $\pmodel$ is parameterized as  
\begin{align}
    \pmodel(\dt{x}{0:T}) :=p(\xT) \prod_{t=T}^{1}\pmodel(\xprev \mid \xt), \qquad 
    \pmodel(\xprev \mid \xt ) := \mathcal{N}(\xprev; \xt + \sigma_t^2 \scoremodel (\xt, t), \sigma_t^2)
\end{align}
where the score network $\scoremodel$ is modeled through a denoiser network $\denoiser$ by $\scoremodel(\xt, t) := (\denoiser (\xt, t; \theta) - \xt )/\bar \sigma_t^2$.
Note that the constant schedule is a special case where $\sigma_t^2=\sigma^2$, and $\bar \sigma_t^2=t\sigma^2$ for all $t$. 

The TDS algorithm for general VE models is described in \Cref{alg:TDS}, where $t\sigma^2$ in \Cref{alg_line:conditional_score} is replaced by $\bar \sigma_t^2$ and $\sigma^2$ in \Cref{alg_line:proposal} is replaced by $\sigma_t^2$. 



\paragraph{Extension to variance preserving diffusion models.}  Another widely used diffusion framework is variance preserving (VP) diffusion models \citep{ho2020denoising}. VP models define the forward process $q$ by 
\begin{align}\label{eq:VP-forward-q}
q(\dt{x}{1:T} \mid \xzero):=\prod_{t=1}^T q(\xt \mid \dt{x}{t-1}), \qquad 
    q(\xt \mid \dt{x}{t-1}) := \mathcal{N}(\xt; \sqrt{1-\sigma_t^2}\dt{x}{t-1}, \sigma_t^2)
\end{align}
where $\sigma_t^2$ is  a sequence of increasing variances chosen such that $q(\xT) \approx \mathcal{N}(0,1)$, and so one can set $p(\xT) = \mathcal{N}(0,1)$. Define $\alpha_t := 1-\sigma_t^2$,  $\bar \alpha_t := \prod_{t^\prime=1}^t\alpha_{t^\prime}$, and $\bar \sigma_t^2:=1-\bar \alpha_t$. 
Then the marginal conditional of \cref{eq:VE-forward-q} is $q(\xt \mid \xzero) = \mathcal{N}(\xt; \sqrt{\bar \alpha_t} \xzero, \bar \sigma_t^2)$. The reverse diffusion process $\pmodel$ is parameterized as  
\begin{align}
    \pmodel(\dt{x}{0:T}) :=p(\xT) \prod_{t=T}^{1}\pmodel(\xprev \mid \xt), \quad 
    \pmodel(\xprev \mid \xt ) := \mathcal{N}(\xprev; \frac{1}{\sqrt{\alpha_t}}\xt + \frac{\sigma_t^2}{\sqrt{\alpha_t}} \scoremodel (\xt, t), \sigma_t^2)
\end{align}
where $\scoremodel$ is now defined through the denoiser $\denoiser$ by $\scoremodel(\xt, t) := (\sqrt{\bar \alpha_t} \denoiser (\xt, t; \theta) - \xt )/\bar \sigma_t^2$.


TDS in \Cref{alg:TDS} extends to VP models as well,
where the conditional score approximation in \Cref{alg_line:conditional_score} 
is changed to  
$\tilde s_k \leftarrow (\sqrt{\bar \alpha_{t+1} }\denoiser(\xtpo_k) - \xtpo_k)/\bar \sigma_{t+1}^2 + \nabla_{\xtpo_k} \log \tilde p_k^{t+1}$, 
and the proposal in \Cref{alg_line:proposal} is changed to 
$\xt_k \sim \pTwist (\cdot \mid \xtpo_k, y) :=\mathcal{N}\left(\frac{1}{\sqrt{\alpha_{t+1}}}\xtpo_k + \frac{\sigma_{t+1}^2}{\sqrt{\alpha_{t+1}}} \tilde s_k, \tilde \sigma^2 \right).$


\subsection{TDS parameters}\label{sec:TDS_hyperparams}
\paragraph{Proposal variance.}
The proposal distribution in  \Cref{alg_line:proposal} of \Cref{alg:TDS} is associated with a variance parameter $\tilde \sigma^2$.
In general, this parameter can be dependent on the time step, i.e.\ replacing $\tilde \sigma^2$ 
 by some $\tilde \sigma_{t+1}^2$ in \Cref{alg_line:proposal}. 
Unless otherwise specified, we set $\tilde \sigma_{t+1}^2:=\mathrm{Var}_{\pmodel}[\dt{x}{t} \mid \xtpo]$ the variance of the unconditional diffusion model (typically learned along with the score network during training). 

\paragraph{Resampling strategy.} The mulinomial resampling strategy in \Cref{alg_line:resampling} of \Cref{alg:TDS} can be replaced by other strategies, see \cite[Chapter 2]{naesseth2019elements} for an overview.
In our experiments, we use the systematic resampling strategy.

In addition, one can consider setting an effective sample size (ESS) threshold (between 0 and 1), and only when the ESS is smaller than this threshold, the resampling step is triggered. 
ESS thresholds for resampling are commonly used to improve efficiency of SMC algorithms \citep[see e.g.][Chapter 2.2.2]{naesseth2019elements},
but for simplicity we use TDS with resampling at every step unless otherwise specified.


\subsection{Derivation of \Cref{eq:twisted-intermediates}}\label{sec:twisted_intermediates_derivation}
%: intermediate twisted targets approach the final target}\label{sec:intermediate_TDS_derivation}
\Cref{eq:twisted-intermediates} illustrated how the extended intermediate targets $\target_t(\dt{x}{0:T})$ provide approximations to the final target $\pmodel(\dt{x}{0:T} \mid y)$ that become increasingly accurate as $t$ approaches 0.
We obtain \Cref{eq:twisted-intermediates} by substituting the proposal distributions in \Cref{eq:diffusion-twisted-proposals} and weights in \Cref{eq:diffusion-twisted-weigthing-functions} into \Cref{eq:final-target} and simplifying.
\begin{align}
   \nu_t(\dt{x}{0:T}) &\propto
   \left[p(\xT)\prod_{t^\prime=0}^{T-1} \pTwist(\dt{x}{t^\prime} \mid \dt{x}{t^\prime+1}, y) \right]
   \left[
  \tilde p(y\mid \xT) \prod_{t^\prime=t}^{T-1} 
    \frac{\pmodel (\dt{x}{t^\prime} \mid \dt{x}{t^\prime +1}) \pTwist(y\mid \dt{x}{t^\prime})}{\pTwist(y \mid \dt{x}{t^\prime+1}) \pTwist(\dt{x}{t^\prime} \mid \dt{x}{t^\prime+1}, y)}
   \right]\\
   & \textrm{Rearrange \ and\ cancel\ }\pmodel \mathrm{\ and \ } \pTwist \textrm{\ terms.}\\ 
   &= \left[ p(\xT) \prod_{t^\prime=t}^{T-1} \pmodel(\dt{x}{t^\prime} \mid  \dt{x}{t^\prime+1}) \right] \left[\tilde p_\theta(y\mid \dt{x}{t})\prod_{t^\prime=0}^{t-1} \pTwist(\dt{x}{t^\prime} \mid \dt{x}{t^\prime +1}, y)\right]\\
   & \textrm{Group\ }\pmodel \textrm{\ terms by chain rule of probability.}\\ 
   &=  \pmodel(\dt{x}{t:T}) \pTwist (y\mid \xt) \prod_{t^\prime=0}^{t-1} \pTwist (\dt{x}{t^\prime} \mid \dt{x}{t^\prime +1}, y)\\
   & \textrm{Apply\ Bayes'\ rule and note that } \pmodel (y\mid \dt{x}{t:T}) = \pmodel(y\mid \xt). \\
   &\propto \pmodel(\dt{x}{t:T}\mid y) \frac{\tilde p_\theta(y\mid \xt)}{p(y \mid \xt)} \prod_{t^\prime=0}^{t-1} \pTwist (\dt{x}{t^\prime} \mid \dt{x}{t^\prime +1}, y)\\
%   & \textrm{Multiply\ by\  }\prod_{t^\prime=0}^{t-1} \pmodel(\dt{x}{t^\prime}\mid \dt{x}{t^\prime + 1}, y)/ \pmodel(\dt{x}{t^\prime}\mid \dt{x}{t^\prime + 1}, y)=1.\\
 & \textrm{Note that }  \pmodel(\dt{x}{0:T} \mid y) = \pmodel (\dt{x}{t:T} \mid y)   \prod_{t=0}^{t'-1} \pmodel(\dt{x}{t^\prime}\mid \dt{x}{t^\prime + 1}, y). \\
   &=  \pmodel(\dt{x}{0:T}\mid y) \left[\frac{\tilde p_\theta(y\mid \xt)}{p(y \mid \xt)} \prod_{t^\prime=0}^{t-1} \frac{ \pTwist (\dt{x}{t^\prime} \mid \dt{x}{t^\prime +1}, y)}{
   \pmodel(\dt{x}{t^\prime}\mid \dt{x}{t^\prime + 1}, y)}\right].
\end{align}
The final line is the desired expression in \Cref{eq:twisted-intermediates}.
\subsection{Inpainting and degrees of freedom final steps for asymptotically exact target }\label{sec:supp_method_inpaint_alt_final_step}
The final ($t=0$) twisting function for inpainting and inpainting with degrees of freedom described in \Cref{subsec:conditioning-problems} do not satisfy the assumption of \Cref{thm:twisted_SMC_complete} that $\pTwist (y\mid \xzero)=\plikelihood{y}{\xzero}.$
This choice introduces error in the final target of TDS relative to the exact conditional $\pmodel(\dt{x}{0:T} \mid y).$

For inpainting, to maintain asymptotic exactness one may instead choose the final proposal and weights as 
\ifnum\workshopsubmission=0
\begin{align}\label{eqn:inpaint_final_proposal_and_weight}
    \pTwist(\xzero \mid \dt{x}{1}, y; \mask) :=  \delta_{y}(\xzero_\mask)\pmodel(\xzero_\umask \mid \dt{x}{1}) \quad \mathrm{and} \quad \weight_0(\xzero, \dt{x}{1}) = 1.
\end{align}
\else
\begin{align}\label{eqn:inpaint_final_proposal_and_weight}
    \pTwist(\xzero \mid \dt{x}{1}, y; \mask) :=  \delta_{y}(\xzero_\mask)\pmodel(\xzero_\umask \mid \dt{x}{1})
\end{align}
and $\weight_0(\xzero, \dt{x}{1}) = 1.$
\fi
One can verify the resulting final target is $\pmodel(\xzero \mid \xzero_\mask =y)$ according to \cref{eq:final-target}.  

Similarly, for inpainting with degrees of freedom, one may define the final proposal and weight as
\ifnum\workshopsubmission=0
\begin{equation}
    \pTwist(\xzero \mid \dt{x}{1}, y, \mathcal{M}) :=  \sum_{\mask \in \mathcal{M}} \frac{\pmodel(\xzero_\mask \mid \dt{x}{1}, \mask)}{\sum_{\mask^\prime \in \mathcal{M}} \pmodel(\xzero_{\mask^\prime} \mid \dt{x}{1}, \mask^\prime)
    }
    \delta_{y}(\xzero_\mask)\pmodel(\xzero_\umask \mid \dt{x}{1}) 
    \quad \mathrm{and} \quad \weight_0(\xzero, \dt{x}{1}) = 1.
\end{equation}
\else
$\weight_0(\xzero, \dt{x}{1}) = 1$ and 
\begin{align*}
\begin{split}
    &\pTwist(\xzero \mid \dt{x}{1}, y; \mathcal{M}) :=  \sum_{\mask \in \mathcal{M}} \pmodel(\xzero_\mask \mid \dt{x}{1};\mask)\delta_{y}(\xzero_\mask)\pmodel(\xzero_\umask \mid \dt{x}{1}).
\end{split}
\end{align*}
\fi

\subsection{Asymptotic accuracy of TDS -- additional details and full theorem statement}\label{sec:theorem_statement_and_proof}
In this section we (i) characterize sufficient conditions on the model and twisting functions under which TDS provides arbitrarily accurate estimates as the number of particles is increased and 
(ii) discuss when these conditions will hold in practice for the twisting function $\pTwist(y\mid \xt)$ introduced in \Cref{sec:method}.

We begin with a theorem providing sufficient conditions for asymptotic accuracy of TDS.
\begin{theorem}\label{thm:twisted_SMC_complete}
Let $\pmodel(\dt{x}{0:T})$ be a diffusion generative model
(defined by \cref{eqn:diffusion_marginal,eq:unconditional-diffusion-transition})
with 
$$
\pmodel(\xt \mid \xtpo) = \mathcal{N}\left(\xt \mid \xtpo + \sigma_{t+1}^2 \scoremodel(\xtpo), \sigma_{t+1}^2 I
\right),$$
with variances $\sigma_1^2, \dots, \sigma_T^2.$ 
Let $\pTwist(y \mid \xt)$ be twisting functions,
and 
$$
\proposal_t(\xt \mid \xtpo) = 
\mathcal{N}\left(
\xt \mid \xtpo + \sigma_{t+1}^2 [\scoremodel(\xtpo) + \nabla_{\xtpo}\log \pTwist(y\mid \xtpo)], 
\tilde \sigma_{t+1}^2 \right)
$$
be proposals distributions for $t=0,\dots, T-1,$
and let $\mathbb{P}_{K}=\sum_{k=1}^K \dt{w_k}{0}\delta_{\dt{x_k}{0}}$ for weighted particles $\{(\dt{x_k}{0}, \dt{w_k}{0})\}_{k=1}^K$ returned by \Cref{alg:TDS} with $K$ particles.
Assume
\begin{enumerate}[label=(\alph*)]
    \item{the final twisting function is the likelihood, $\pTwist(y\mid \xzero)=\plikelihood{y}{\xzero},$} \label{item:final_twist}
    \item{the first twisting function $\pTwist(y \mid \xT),$ and the ratios of subsequent twisting functions $\pTwist(y \mid \xt) / \pTwist(y \mid \xtpo) $ are positive and bounded,}\label{item:twist_assumption}
    %\item{each $\nabla_{\xt} \log \pTwist(y\mid \xt)$ with $t>0$ is continuous and has bounded gradients, and}\label{item:twist_gradient}
    \item{each $\log \pTwist(y\mid \xt)$ with $t>0$ is continuous and has bounded gradients in $\xt$, and}\label{item:twist_gradient}
    \item{the proposal variances are larger than the model variances, i.e.\ for each $t,\ \tilde \sigma_t^2 >\sigma_t^2.$} \label{item:proposal_var}
\end{enumerate}
Then $\mathbb{P}_{K}$ converges setwise to $\pmodel(\xzero\mid y)$ with probability one,
that is 
for every set $A,$ 
$\lim_{K\rightarrow\infty} \mathbb{P}_{K}(A) = \int_A \pmodel(\xzero\mid y) d\xzero.$
\end{theorem}

The assumptions of \Cref{thm:twisted_SMC_complete} are readily satisfied in common applications.
\begin{itemize}
\item {Assumption \ref{item:final_twist} may be satisfied by construction by choosing $\pTwist(y\mid \xzero)=\plikelihood{y}{\xzero}$
by defining $\denoiser(\xzero, t=0)=\xzero.$
}
\item{Assumption \ref{item:twist_assumption} is satisfied if (i) $\plikelihood{y}{x}$ is smooth in $x$ and everywhere positive and (ii) $\pTwist(y\mid \xt) = \plikelihood{y}{\denoiser(\xt)}$ where $\denoiser(\xt)$ takes values in some compact domain.
An alternative sufficient condition for Assumption \ref{item:twist_assumption} is for $\plikelihood{y}{x}$ to be positive and bounded away from zero; 
this latter condition will hold when, for example, $\plikelihood{y}{x}$ is a classifier fit with regularization.}
\item{Assumption \ref{item:twist_gradient} is the strongest assumption.
It will be satisfied, for example, if (i) $\pTwist(y\mid \xt) = \plikelihood{y}{\denoiser(\xt)},$ and (ii) $\plikelihood{y}{x}$ and $\denoiser(\xt)$ are smooth in $x$ and $\xt,$ with uniformly bounded gradients.
While smoothness of $\denoiser(\cdot, t)$ can be encouraged by the use of skip-connections and regularization,
$\denoiser(\cdot, t)$ may present sharp transitions,
particularly for $t$ close to zero.}
\item{Assumption \ref{item:proposal_var}, that the proposal variances satisfy $\tilde \sigma_{t}^2 > \sigma_t^2$ is likely not needed for the result to hold.
However, this assumption permits usage of existing SMC theoretical results in the proof;
in practice, our experiments use $\tilde \sigma_t^2 = \sigma_t^2,$
but alternatively the assumption could be met by inflating each $\tilde \sigma_t^2$ by some arbitrarily small $\delta$ without markedly impacting the behavior of the sampler.}
\end{itemize}

%\paragraph{\BLT{Cut this paragraph for NeurIPS appendix submission} Inpainting and inpaiting with degrees of freedom.}

%\begin{cor}
%The inpainting and DOF cases are consistent for their conditionals.
%\end{cor}
%\begin{proof}
%We can see that early truncation at step $t=1$ forms a special case.
%By reindexing $t$ we find we are consistent for $\pmodel(\dt{x}{1} \mid y).$
%Then those final samplers sample $\xzero \sim \pmodel(\xzero \mid \dt{x}{1}, y)$ in the last step.
%So by the law of total probability and chain rule of probability they give exact samples in the limit as desired.
%\end{proof}
%
%Why is bounded assumption reasonable in the inpainting case?
%It's because each $\sigma_t^2>0$ and so with 
%$\pTwist(y\mid \xt; \mask) = \mathcal{N}(y; \denoiser(\xt)_\mask, \bar \sigma_t^2)=p(y\mid \dt{x}{1})$
%with $\bar \sigma_t^2 = \sum_{t^\prime=1}^t\sigma_{t^\prime}^2,$
%the twisting variance are a decreasing sequence.
%And since the denoiser is also bounded and each $\bar \sigma_t^2>0$, there is some $0 <c <C< \infty$ such that 
%$c<\pTwist(y\mid \xt) <C$ and so the ratio is bounded.

\paragraph{Proof of \Cref{thm:twisted_SMC_complete}:}
\Cref{thm:twisted_SMC_complete} characterizes a set of conditions under which SMC algorithms converge.
We restate this result below in our own notation.
\begin{theorem}[\citet{chopin2020introduction} -- Proposition 11.4]\label{thm:chopin_theorem}
Let $\{(\dt{x_k}{0}, \dt{w_k}{0})\}_{k=1}^K$ be the particles and weights returned at the last iteration of a sequential Monte Carlo algorithm with $K$ particles using multinomial resampling.
If each weighting function $\weight_t(\xt, \xtpo)$ is positive and bounded, 
then for every bounded, $\target_0$-measurable function $\phi$ of $\xt$
$$
\lim_{K\rightarrow \infty} \sum_{k=1}^K \dt{w_k}{0} \phi(\dt{x_k}{0})
=
\int \phi(\xzero)\target_0(\xzero)d\xzero.
$$
with probability one.
\end{theorem}


An immediate consequence of \Cref{thm:chopin_theorem} is the setwise convergence of the discrete measures, $\hat P_K = \sum_{k=1}^K \dt{w_k}{0} \delta_{\dt{x_k}{0}}.$
This can be seen by taking for each $\phi(x) = \I[x \in A]$ for any $\target_0$-measurable set $A$.
The theorem applies both in the Euclidean setting, where each $\dt{x_k}{t}\in \R^D,$ as well as the Riemannian setting.

We now proceed to prove \Cref{thm:twisted_SMC_complete}.
\begin{proof}
To prove the theorem we show
(i) the $\xzero$ marginal final target $\target_0$ is $\pmodel(\xzero \mid y)$ and then
(ii) $\mathbb{P}_K$ converges setwise to $\target_0.$

We first show (i) by manipulating $\target_0$ in \Cref{eq:final-target} to obtain $\pmodel(\dt{x}{0:T} \mid y).$
From \Cref{eq:final-target} we first have
\begin{align}
\begin{split}
   \target_0(\dt{x}{0:T}) &= \frac{1}{\mathcal{L}_0} 
        \left[ \proposal(\dt{x}{T})  \prod_{t=0}^{T-1} \proposal_{t}(\dt{x}{t}\mid  \dt{x}{t+1})  \right] 
        \left[ \weight_T(\dt{x}{T})
        \prod_{t=0}^{T-1} \weight_{t^\prime}(\dt{x}{t}, \dt{x}{t+1}) \right]   \\
    &\mathrm{Substitute\ in\ weights\ from\ \cref{eq:diffusion-twisted-weigthing-functions}.} \\
    &= \frac{1}{\mathcal{L}_0} 
        \left[ p(\dt{x}{T}) \prod_{t=0}^{T-1} r_t (\xt \mid \xtpo)
        \right] 
        \left[ \pTwist(y\mid \dt{x}{T})
        \prod_{t=0}^{T-1} 
        \frac{\pmodel(\xt \mid \xtpo)\pTwist(y\mid \xt)}{\pTwist(y \mid \xtpo)r_t(\xt \mid \xtpo)}
        \right] \\
    &\mathrm{Rearrange\ }\pmodel\mathrm{\ and \ } r_t\ \mathrm{\ terms, \ and\ cancel\ out\ }r_t \mathrm{\ terms.} \\
    &= \frac{1}{\mathcal{L}_0} 
        \left[ p(\dt{x}{T}) \prod_{t=0}^{T-1} \pmodel(\xt \mid \xtpo)
        \right] 
        \left[ \pTwist(y\mid \dt{x}{T})
        \prod_{t =0}^{T-1} 
        \frac{\cancel{r_t(\xt \mid \xtpo)}\pTwist(y\mid \xt)}{\pTwist(y \mid \xtpo)\cancel{r_t(\xt \mid \xtpo)}}
        \right] \\
    &\mathrm{Collapse\ }\pmodel\mathrm{\ terms.}\\ %,\ rearrange\ and\ cancel\ }\pTwist \mathrm{\ terms.} \\
    &= \frac{1}{\mathcal{L}_0} 
        \pmodel(\dt{x}{0:T})
        \left[ 
        \prod_{t=0}^{T-1} 
        \frac{\pTwist(y\mid \xt)}{\pTwist(y \mid \xtpo)}\right]\pTwist(y \mid \dt{x}{T})\\
    &\mathrm{Cancel \ out \ } \pTwist \mathrm{\ terms.}\\
    &= \frac{1}{\mathcal{L}_0} 
        \pmodel(\dt{x}{0:T}) \pTwist (y \mid \xzero) \\ 
    &\mathrm{Recognize\ }\pTwist(y\mid \dt{x}{0})=\plikelihood{y}{\xzero}\mathrm{\ by \ Assm. \ref{item:final_twist}\ \ and\  apply\  Bayes'\  rule\  with\  }\mathcal{L}_0=\pmodel(y). \\
    &= \pmodel(\dt{x}{0:T}\mid y).
\end{split}
\end{align}
The final line reveals that once we marginalize out $\dt{x}{1:T}$
we obtain $\target_0(\xzero)=p(\xzero \mid y)$ as desired.

We next show that $\mathbb{P}_K$ converges to $\target_0$ with probability one by applying \Cref{thm:chopin_theorem}.
To apply \Cref{thm:chopin_theorem} it is sufficient to show that the weights at each step are upper bounded, as they are defined through (ratios of) probabilities and hence are positive. 
Since there are a finite number of steps $T,$
it is enough to show that each $w_t$ is bounded.
The inital weight is the initial twisting function, $w_T(\xT)=\pTwist(y\mid \xT),$ which is bounded by Assumption \ref{item:twist_assumption}.
So we proceed to intermediate weights.

To show that the weighting functions at subsequent steps are bounded,
we decompose the log-weighting functions as 
$$
\log \weight_t(\xt, \xtpo) = \log \frac{\pTwist(y\mid \xt)}{\pTwist(y\mid \xtpo)} + \log \frac{\pmodel(\xt\mid \xtpo)}{r_t(\xt \mid \xtpo)} ,
$$
and show independently that $\log \pTwist(y\mid \xt)/\pTwist(y\mid \xtpo)$ and
$\log \pmodel(\xt\mid \xtpo)/r_t(\xt \mid \xtpo)$ are bounded.
The first term $\log \pTwist(y\mid \xt)/\pTwist(y\mid \xtpo)$ is again bounded by Assumption \ref{item:twist_assumption},
and we proceed to the second.

That $\log \pmodel(\xt\mid \xtpo)/r_t(\xt \mid \xtpo)$ is bounded follows from Assumptions \ref{item:twist_gradient} and \ref{item:proposal_var}.
First write
$$
%\pmodel(\xt \mid \xtpo) = \mathcal{N}\left(\xt \mid \xtpo + \sigma_{t+1}^2 \scoremodel(\xtpo), \sigma_{t+1}^2 I
\pmodel(\xt \mid \xtpo) = \mathcal{N}\left(\xt \mid \hat \mu, \sigma_{t+1}^2 I
\right)$$
with $\hat \mu=\xtpo + \sigma_{t+1}^2 \scoremodel(\xtpo),$
and 
$$
\proposal_t(\xt \mid \xtpo) = 
\mathcal{N}\left(
\xt \mid  \hat \mu_\psi, \tilde \sigma_{t+1}^2 I \right),
$$
for $\hat \mu_\psi = \hat \mu + \sigma_{t+1}^2 \nabla_{\xtpo}\log \pTwist(y\mid \xtpo)$.
The log-ratio then simplifies as 
\begin{align}
\log \frac{\pmodel(\xt\mid \xtpo)}{r_t(\xt \mid \xtpo)} &= 
\log \frac{| 2\pi \sigma_{t+1}^2 I|^{-1/2} \exp\{-(2\sigma_{t+1}^2)^{-1}\|\hat \mu - \xt \|^2 \}}
{ | 2\pi \tilde \sigma_{t+1}^2 I|^{-1/2} \exp\{-(2\tilde \sigma_{t+1}^2)^{-1}\|\hat \mu_\psi- \xt \|^2 \}} \\
&\textrm{Rearrange \ and\  let  \ } C=\log | 2\pi \sigma_{t+1}^2 I|^{-1/2} /  | 2\pi \tilde \sigma_{t+1}^2 I|^{-1/2}\\
&= \frac{-1}{2}\left[ 
\sigma_{t+1}^{-2}\|\hat \mu - \xt \|^2    -
\tilde \sigma_{t+1}^{-2}\|\hat \mu_\psi - \xt \|^2 \right]  +C \\
&\textrm{Expand and rearrange \ } \|\hat \mu_\psi - \xt \|^2 = \| \hat \mu  - \xt \|^2 + 2 \langle \hat\mu_\psi - \hat \mu, \hat \mu - \xt  \rangle + \|\hat \mu_\psi - \hat \mu\|^2 \\
&= \frac{-1}{2}\left[ ( \sigma_{t+1}^{-2} - \tilde \sigma_{t+1}^{-2})\|\hat \mu - \xt \|^2  
- 2 \tilde \sigma_{t+1}^{-2}\langle \hat\mu_\psi - \hat \mu, \hat \mu - \xt  \rangle - \tilde \sigma_{t+1}^2 \|\hat \mu_\psi - \hat \mu\|^2 
\right]+ C\\
&\textrm{Let \ } C^\prime {=} C - \frac{1}{2}\tilde \sigma_{t+1}^2 \|\hat \mu_\psi - \hat \mu\|^2 \textrm{\ and rearrange.  Note that $\|\hat \mu_\psi {-} \hat \mu\|^2{<}\infty$ by Assm. \ref{item:twist_gradient}.}\\
&= \frac{-1}{2}(\sigma_{t+1}^{-2} - \tilde \sigma_{t+1}^{-2})\|\hat \mu - \xt \|^2 +
\tilde \sigma_{t+1}^{-2}\langle \hat\mu_\psi - \hat \mu, \hat \mu - \xt  \rangle + C^\prime\\
&\textrm{Apply Cauchy-Schwarz}\\
&\le  \frac{-1}{2} (\sigma_{t+1}^{-2} - \tilde \sigma_{t+1}^{-2})\|\hat \mu - \xt \|^2
+ \tilde \sigma_{t+1}^{-2}\| \hat\mu_\psi - \hat \mu\|\cdot \| \hat \mu - \xt  \| + C^\prime\\
&\textrm{Upper-bounding using that $\mathrm{max}_{x}\frac{-a}{2}x^2 + bx =\frac{b^2}{2a}$ for $a = \sigma_{t+1}^{-2} - \tilde \sigma_{t+1}^{-2}>0$ by Assm. (d).}\\
&\le  \frac{1}{2} \frac{(\tilde \sigma_{t+1}^{-2}\| \hat\mu_\psi - \hat \mu\|)^2}{\sigma_{t+1}^{-2} - \tilde \sigma_{t+1}^{-2}} + C^\prime \\
&= \frac{\tilde \sigma_{t+1}^{-4}}{2( \sigma_{t+1}^{-2} - \tilde \sigma_{t+1}^{-2})}\| \hat\mu_\psi - \hat \mu\|^2  + C^\prime \\
& \textrm{Note that } \hat \mu_\psi = \hat \mu + \sigma_{t+1}^2 \nabla_{\xtpo}\log \pTwist(y\mid \xtpo). \\ 
&= \frac{\tilde \sigma_{t+1}^{-4}}{2( \sigma_{t+1}^{-2} - \tilde \sigma_{t+1}^{-2})}\sigma_{t+1}^4 \| \nabla_{\xtpo} \log \pTwist(y \mid \xtpo) \|^2  + C^\prime \\
&\le C^{\prime\prime}.
\end{align}
The final line follows from Assumption \ref{item:twist_gradient}, that the gradients of the twisting functions are bounded.
The above derivation therefore provides that each $\weight_t$ is bounded, concluding the proof.
\end{proof}

\section{Riemannian \methodname{}}\label{sec:riemannian}

This section provides additional details on the extension of TDS to Riemannian diffusion models introduced in \Cref{sec:method}.
We first introduce the tangent normal distribution.
We then provide with background on Riemannian diffusion models, which we parameterize with the tangent normal.
Then we describe the extension of TDS to these models.
Finally we show how \Cref{alg:TDS} modifies to this setting.

\paragraph{The tangent normal distribution.}
Just as the Gaussian is the parametric family underlying Euclidean diffusion models in \Cref{eq:unconditional-diffusion-transition}, 
the tangent normal (see e.g.\ \citep{chirikjian2014gaussian,deriemannian})
underlies generation in Riemannian diffusion models so we review it here.

We take the tangent normal to be the distribution is implied by a two step procedure.
Given a variable $x$ in the manifold, the first step is to sample a variable $\bar y$ in $\mathcal{T}_{x},$ the tangent space at $x;$
if $\{h_1, \dots, h_D\}$ is an orthonormal basis of $\mathcal{T}_{x}$
one may generate
\begin{align}
    \bar y = \mu + \sum_{d=1}^D \sigma \epsilon_d \cdot h_d,
\end{align}
with $\epsilon_d\overset{i.i.d.}{\sim}\mathcal{N}(0, 1),$
and $\sigma^2>0$ a variance parameter. 
The second step is to project $\bar y$ back onto the manifold to obtain
$y = \exp_{x}\{ \bar y\}$
where $\exp_{x}\{ \cdot \}$ denotes the exponential map at $x.$ The resulting distribution of $y$ is denoted as $\mathcal{TN}_{x}(\mu, \sigma^2)$. 
By construction, $\mathcal{T}_x$ is a Euclidean space with its origin at $x,$
so when $\|\mu\|_2=0, \mathcal{TN}_{x}(\mu, \sigma^2)$ is centered on $x.$
And since the geometry of a Riemannian manifold is locally Euclidean, when $\sigma^2$ is small
the exponential map is close to the identity map and the tangent normal is essentially a narrow Gaussian distribution in the manifold at $x.$
Finally, we use $\mathcal{TN}_{x}(y;\mu, \sigma^2)$ to denote the density of the tangent normal evaluated at $y$.

Because this procedure involves a change of variables from $\bar y$ to $y$ (via the exponential map),
to compute the tangent normal density one computes 
\begin{equation}\label{eq:TN_density}
    \mathcal{TN}_{x}(y;\mu, \sigma^2) = \mathcal{N}(\exp^{-1}_x\{y\}; \mu,  \sigma^2) \left| \frac{\partial}{\partial y} \exp_x^{-1}\{y\}\right|
\end{equation}
where $\exp_x^{-1}\{y\}$ is the inverse of the exponential map (from the manifold into $\mathcal{T}_x$), 
and $\left| \frac{\partial}{\partial y} \exp_x^{-1}\{y\}\right|$ is the determinant of the Jacobian of the exponential map; when the manifold lives in a higher dimensional subset of $R,$ we take 
$\left| \frac{\partial}{\partial y} \exp_x^{-1}\{y\}\right|$ to be the product of the positive singular values of the Jacobian.

\paragraph{Riemannian diffusion models and the tangent normal distribution.}
Riemannian diffusion models proceeds through a geodesic random walk \citep{deriemannian}.
At each step $t,$
one first samples a variable $\dt{\bar x}{t}$ in $\mathcal{T}_{\xtpo};$
if $\{h_1, \dots, h_D\}$ is an orthonormal basis of $\mathcal{T}_{\xtpo}$
one may generate
\begin{align}
    \dt{\bar x}{t}  =\sigma_{t+1}^2 s_\theta(\dt{x}{t+1})+ \sum_{d=1}^D \sigma_{t+1} \epsilon_d \cdot h_d,
\end{align}
with $\epsilon_d\overset{i.i.d.}{\sim}\mathcal{N}(0, 1).$
One then projects $\dt{\bar x}{t}$ back onto the manifold to obtain
$\xt = \exp_{\xtpo}\{ \dt{\bar x}{t}\}$
where $\exp_{x}\{ \cdot \}$ denotes the exponential map at $x.$

This is equivalent to sampling $\xt$ from a tangent normal distribution as 
\begin{equation}\label{eqn:reverse_transition_unconditional_riemannian}
    \dt{x}{t} \sim p(\dt{ x}{t} \mid \xtpo) = \mathcal{TN}_{\xtpo}\left(\dt{ x}{t};\sigma_{t+1}^2 s_\theta(\dt{x}{t+1}),
\sigma_{t+1}^2 \right).
\end{equation}

\paragraph{TDS for Riemannian diffusion models.}
To extend TDS, appropriate analogues of the twisted proposals and weights are all that is needed.
For this extension we require that the diffusion model is also associated with a manifold-valued denoising estimate $\denoiser$ as will be the case when, for example, $\scoremodel(\xt, t) := \nabla_{\xt} \log q(\xt \mid \xzero=\denoiser)$ for $\denoiser = \denoiser(\xt, t).$
In contrast to the Euclidean case, a relationship between a denoising estimate and a computationally tractable score approximation may not always exist for arbitrary Riemannian manifolds;
however for Lie groups when the the forward diffusion is the Brownian motion, tractable score approximations do exist \citep[Proposition 3.2]{yim2023se}.

For the case of positive and differentiable $\plikelihood{y}{\xzero},$ we again choose twisting functions $\twist(y\mid \xt):= \plikelihood{y}{\denoiser(\xt)}.$

Next are the inpainting and inpainting with degrees of freedom cases.
Here, assume that $\xzero$ lives on a multidimensional manifold (e.g.\ $SE(3)^N$)
and the unmasked observation $y=\xzero_\mask$ with $\mask \subset \{1, \dots, N\}$ on a lower-dimensional sub-manifold (e.g.\ $SE(3)^{|\mask|},$ with $|\mask|<N$).
In this case, twisting functions are constructed exactly as in \Cref{subsec:conditioning-problems}, 
except with the normal density in \Cref{eqn:inpaint_twist} replaced with a Tangent normal as
\begin{align}
\twist(y\mid \xt) {=} \mathcal{TN}_{\denoiser(\xt)_\mask}( y ; 0, \bar \sigma_t^2).
\end{align}

For all cases, we propose the twisted proposal as
\begin{equation}\label{eqn:likelihood_approx_riemannian_supp}
    \pTwist(\dt{ x}{t} \mid \xtpo, y) {=} \mathcal{TN}_{\xtpo}\left(\dt{ x}{t};\sigma_{t+1}^2 s_\theta(\dt{x}{t+1}, y), \tilde\sigma_{t+1}^2 \right)
\end{equation}
%where $\bar y := \exp_{\denoiser(\dt{x}{t})_\umask}^{-1}(y)\in \mathcal{T}_{\denoiser(\dt{x}{t})_\umask}$ 
%is $\dt{x}{0}$ mapped into $\mathcal{T}_{\denoiser(\dt{x}{t})}$ by the inverse of the exponential map at $\denoiser(\dt{x}{t})_\umask,$
where as in the Euclidean case  $s_\theta(\dt{x}{t}, y)=s_\theta(\dt{x}{t}) + \nabla_{\dt{x}{t}} \log \tilde p(y \mid \dt{x}{t})$.


Weights at intermediate steps are computed as in the Euclidean case (\Cref{eq:diffusion-twisted-weigthing-functions}): 
%However, 
%since the tangent normal includes a projection step back to the manifold, its density involves not only usual Gaussian density but also the Jacobian of the exponential map to account for the change of variables.
%For example, in the case that the ambient dimension of the manifold and the tangent space are the same dimension,
%$$
%\mathcal{TN}_{x}(x^\prime;\mu, \sigma^2)= \mathcal{N}(\exp^{-1}_x\{x^\prime\}; \mu, \sigma^2)\left|\frac{\partial}{\partial x} \exp_x^{-1}\{x^\prime\}\right|,
%$$
%where $\exp^{-1}_x\{x^\prime\}$ is the inverse of the exponential map at $x$ from the manifold into $\mathcal{T}_x,$
%and 
%$\left|\frac{\partial}{\partial x} \exp_x^{-1}\{x^\prime\}\right|$ is the determinant of its Jacobian.
%The weights are then computed as
\begin{align}\label{eq:diffusion-twisted-weigthing-functions-riemannian}
    \weight_t (\xt , \xtpo ):&= \frac{\pmodel (\xt \mid \xtpo) \twist(y \mid \xt)}{\twist(y \mid \xtpo) \pTwist(\xt \mid \xtpo, y)} \\
    &= \frac{\mathcal{TN}_{\xtpo}(\xt; \sigma_{t+1}^2\scoremodel(\xtpo), \sigma_{t+1}^2) \twist(y \mid \xt)
    }{\twist(y \mid \xtpo) \mathcal{TN}_{\xtpo}(\xt; \sigma_{t+1}^2\scoremodel(\xtpo, y), \tilde\sigma_{t+1}^2)}.
\end{align}
While the proposal and target contribute identical Jacobian determinant terms that cancel out, they remain in the twisting functions.


\paragraph{Adapting the TDS algorithm to the Riemannian setting.}
To translate the TDS algorithm to the Riemannian setting we require only two changes.

The first is on \Cref{alg:TDS} 
\Cref{alg_line:conditional_score}.
Here we assume that the unconditional score is computed through the transition density function: 
\begin{align}
\scoremodel(\xtpo) := \nabla_{\xtpo} \log q_{t+1{|}0}(\xtpo\mid \denoiser) \quad \mathrm{\ for\ } \denoiser = \denoiser(\xtpo).
\end{align}
Note that the gradient above ignores the dependence of $\denoiser$ on $\xtpo$. 

%\BLT{Is it clear that the gradient above ignores the dependence of $\denoiser$ on $\xt$?}
The conditional score approximation in \Cref{alg_line:conditional_score} is then replaced with 
$$
\tilde s_k \leftarrow \pTwist (\cdot \mid \xtpo_k, y):= \nabla_{\xtpo_k} \log q_{t+1{|} 0}(\xtpo_k\mid \denoiser) + \nabla_{\xtpo_k} \log \tilde p_k^{t+1} \quad \mathrm{ \ for\  } \denoiser = \denoiser(\xtpo_k).
$$
Notably, $\tilde s_k$ is a $\mathcal{T}_{\xtpo_k}$-valued Riemannian gradient.

The second change is to make the proposal on \Cref{alg:TDS} \Cref{alg_line:proposal} a tangent normal, as defined in \cref{eqn:likelihood_approx_riemannian_supp}.


\section{Additional Related Work}\label{sec:additional_related_work}
% Conditional training
%\paragraph{Conditionally trained models.}
%Much prior work relies on conditional training or fine-tuning for conditional generation. For example, conditional training for text-conditional image generation can involve augmenting the diffusion model to accommodate learned text embeddings as additional inputs ($y$), assembling a data set of paired image and text pairs, and performing denoising training with $y$ provided (\citep{ramesh2022hierarchical}). This strategy amounts to amortization of the conditional generation problem, and has been used for inpainting problems as well with random masks \citep{saharia2022palette,watson2022broadly}. A limitation of these approaches is that they demand the time and expertise needed to (1) construct a satisfactory set of conditional training examples and (2) design and train a neural network that can accomplish the conditional inference task.

\textbf{Training with conditioning information.}
Several approaches involve training a neural network to directly sample from a conditional diffusion models.
These approaches include
(i) conditional training with embeddings of conditioning information, e.g.\ for denoising images \citep{saharia2022image}, and text-to-image generation \citep{ramesh2022hierarchical},
(ii) conditioning training with a subset of the state space, e.g.\ for protein design \citep{watson2022broadly}, and 
image inpainting \citep{saharia2022palette},
(iii) classifier-guidance \citep{dhariwal2021diffusion,song2020score}, an approach for generating samples of a desired class, which requires training a time-dependent classifier to approximate, for each $t,\ \pmodel (y\mid \xt)$, and training such a time-dependent classifier may be inconvenient and costly), and % or even impossible if, for example, one does not have access to labeled data. }
(iv) classifier-free guidance \citep{ho2022classifier}, an approach that builds on classifier-guidance without training a classifier, and instead trains a diffusion model with class information as  additional input.


\textbf{Langevin and Metropolis-Hastings steps.}
Some prior work has explored using Markov chain Monte Carlo steps in sampling schemes to better approximate conditional distributions.
For example unadjusted Langevin dynamics \citep{hovideo,song2020score} or Hamiltonian Monte Carlo \citep{du2023reduce} in principle permit asymptotically exact sampling in the limit of many steps.
However these strategies are only guaranteed to target the conditional distributions of the joint model under the (unrealistic) assumption that the unconditional model exactly matches th forward process, and in practice adding such steps can worsen the resulting samples, presumably as a result of this approximation error \citep{karras2022elucidating}.
By contrast, TDS does not require the assumption that the learned diffusion model exactly matches the forward noising process.
\section{Empirical results additional details}\label{sec:supp_results}


\subsection{Synthetic diffusion models on two dimensional problems}
\begin{figure}[t]
    \tiny
    \includegraphics[scale=0.7]{sections/figures/toy/toy_MOG_A.png}
    \includegraphics[scale=0.7]{sections/figures/toy/toy_MOG_B.png}
    \includegraphics[scale=0.7]{sections/figures/toy/toy_MOG_C.png}
    \includegraphics[scale=0.6]{sections/figures/toy/toy_legend.png}
  \centering
    \caption{
   % Error on small-scale problems with ground truth.  Error of estimate of conditional mean (mean +- 2SEM) with 25 replicates. 
   Errors of conditional mean estimations with 2 SEM error bars averaged over 25 replicates on mixture of Gaussians unconditional target.
   TDS applies to all three tasks and provides increasing accuracy with more particles.  
}
\label{fig:exp-toy-MOG}
\end{figure}

\paragraph{Forward process.}
Our forward process is variance preserving (as described in \Cref{sec:supp_method}) with $T=100$ steps and a quadratic variance schedule.
We set $\sigma^2_t=  \sigma^2_\mathrm{min} + (\frac{t}{T})^2 \sigma^2_{\mathrm{max}}$ with $\sigma^2_\mathrm{min}=10^{-5}$
and $\sigma^2_\mathrm{max}=10^{-1}.$


\paragraph{Unconditional target distributions and likelihood.}
We evaluate the different methods with two different unconditional target distributions:
\begin{enumerate}
    \item A \textbf{bivariate Gaussian} with mean at $(\frac{1}{2},\frac{1}{2})$ and covariance $0.9$ and  
    \item {A \textbf{Gaussian mixture} with three components with mixing proportions $[0.3, 0.5, 0.2],$ means $[(1.54, -0.29), (-2.18, 0.57), (-1.09, -1.40)]$, and $0.2$ standard deviations.}
\end{enumerate}
We evaluate on conditional distributions defined by the three likelihoods described in \Cref{subsec:exp-toy}.

\Cref{fig:exp-toy-MOG} provides results analogous to those in \Cref{fig:exp-toy} but with the mixture of Gaussians unconditional target.
In this second example we evaluate a with a variation of the inpainting degrees of freedom case wherein we consider $y=1$ and $\mathcal{M}=\{[1], [2]\},$ so that $\plikelihood{y}{\xzero}=\delta_y(\xzero_1) + \delta_y(\xzero_2).$


\subsection{Image conditional generation experiments}
We perform extensive ablation studies on class-conditional generatoin task and inpainting task using the small-scale MNIST dataset ($28\times 28 \times 1$ dimensions) in \Cref{supp:subsubsec:mnist-class} and \Cref{supp:subsubsec:mnist-inpainting}. And we show TDS's applicability to higher-dimensional datasets, namely CIFAR10 ($32\times \times 3$ dimensions) and ImageNet256 ($256\times 256 \times 3$ dimensions) in \Cref{supp:subsubec:inet-and-cifar}.


\subsubsection{Class-conditional generation on MNIST}
\label{supp:subsubsec:mnist-class}

\paragraph{Set-up.} For MNIST, we set up a diffusion model using the variance preserving framework. The model architecture is based on the guided diffusion codebase\footnote{\url{https://github.com/openai/guided-diffusion}} with the following specifications: number of channels = 64, attention resolutions = "28,14,7", number of residual blocks = 3, 
learn sigma (i.e. to learn the variance of $\pmodel(\xprev \mid \xt)$) = True,  resblock updown = True, dropout = 0.1, variance schedule = "linear". We trained the model  for 60k epochs with a batch size of 128 and a learning rate of $10^{-4}$ on 60k MNIST training images. The model uses $T=1,000$ for training and $T=100$ for sampling.  

The classifier used for class-conditional generation and evaluation is a pretrained ResNet50 model.\footnote{Downloaded from \url{https://github.com/VSehwag/minimal-diffusion}}
This classifier is trained on the same set of MNIST training images used by diffusion model training. 


In addition we include a
variation called \textit{TDS-truncate} that truncates the TDS procedure at $t = 10$ and returns the prediction $\denoiser (\dt{x}{10})$. 

\paragraph{Sample plots.} 
To supplement the sample plot conditioned on class $7$ in \Cref{subfig:mnist-class-sample}, we present samples conditioned on each of the remaining 9 classes and from other methods. We observe that samples from Gradient Guidance have noticeable artifacts, whereas the other 4 methods produce authentic and
correct digits. However, most samples from IS or TDS-IS are identical due to the collapse of importance weights. By contrast, samples from
TDS and TDS-truncate have greater diversity, with the latter exhibiting slightly more variations
%These results exhibit the same characteristics of relative diversity and quality discussed in \Cref{subsec:exp-mnist}.

\begin{figure}
  \centering
  \begin{subfigure}[t]{0.75\linewidth}
    \centering
   \includegraphics[width=\textwidth]{sections/figures/mnist-class/classifier-guidance-sample-comparision-label0-P64-resampleP16-showmethodnameFalse.pdf}     
\end{subfigure}\hfill 
  \begin{subfigure}[t]{0.75\linewidth}
    \centering
   \includegraphics[width=\textwidth]{sections/figures/mnist-class/classifier-guidance-sample-comparision-label1-P64-resampleP16-showmethodnameFalse.pdf}     
\end{subfigure}\hfill 
  \begin{subfigure}[t]{0.75\linewidth}
    \centering
   \includegraphics[width=\textwidth]{sections/figures/mnist-class/classifier-guidance-sample-comparision-label2-P64-resampleP16-showmethodnameFalse.pdf}     
\end{subfigure}\hfill 
  \begin{subfigure}[t]{0.75\linewidth}
    \centering
   \includegraphics[width=\textwidth]{sections/figures/mnist-class/classifier-guidance-sample-comparision-label3-P64-resampleP16-showmethodnameFalse.pdf}     
\end{subfigure}\hfill 
  \begin{subfigure}[t]{0.75\linewidth}
    \centering
   \includegraphics[width=\textwidth]{sections/figures/mnist-class/classifier-guidance-sample-comparision-label4-P64-resampleP16-showmethodnameFalse.pdf}     
\end{subfigure}\hfill 
  \begin{subfigure}[t]{0.75\linewidth}
    \centering
   \includegraphics[width=\textwidth]{sections/figures/mnist-class/classifier-guidance-sample-comparision-label5-P64-resampleP16-showmethodnameFalse.pdf}     
\end{subfigure}\hfill 
  \begin{subfigure}[t]{0.75\linewidth}
    \centering
   \includegraphics[width=\textwidth]{sections/figures/mnist-class/classifier-guidance-sample-comparision-label6-P64-resampleP16-showmethodnameFalse.pdf}     
\end{subfigure}\hfill 
  \begin{subfigure}[t]{0.75\linewidth}
    \centering
   \includegraphics[width=\textwidth]{sections/figures/mnist-class/classifier-guidance-sample-comparision-label8-P64-resampleP16-showmethodnameFalse.pdf}     
\end{subfigure}\hfill 
  \begin{subfigure}[t]{0.75\linewidth}
    \centering
   \includegraphics[width=\textwidth]{sections/figures/mnist-class/classifier-guidance-sample-comparision-label9-P64-resampleP16-showmethodnameTrue.pdf}     
\end{subfigure}\hfill 
\caption{MNIST class-conditional generation. 16 randomly chosen conditional samples from 64 particles in a single run, given class $y$. From top to bottom, $y=0,1,2,3,4,5,6,8,9$. }
\end{figure}

\begin{figure}[t]
\begin{subfigure}[t]{\linewidth}
    \centering
    \includegraphics[width=0.9\linewidth]{sections/figures/mnist-class/sampler-comparison.pdf}
    \caption{Classification accuracy (measured by the neural network classifier) v.s. number of particles $K$, under different twist scales. Results are averaged over 1{,}000 random runs with error bands indicating 2 standard errors. Across the panels we see that for all the methods, the larger the twist scale (i.e.\ the darker the line color), the higher the classification accuracy. For TDS and TDS-IS, this improvement is more significant for smaller value of $K$,}
    \label{subfig:cg-twist-scale-classifier}
    \end{subfigure}
\begin{subfigure}[t]{\linewidth}
    \centering
    \includegraphics[width=0.6\linewidth]{sections/figures/mnist-class/sampler-comparison-human-rated.pdf}
    \caption{Human-rated classification accuracy v.s. twist scale, for TDS ($K= 64, 2$), TDS-IS ($K= 64, 2$) and Gradient Guidance. Results are averaged over 640 randomly chosen samples with error bands indicating 2 standard errors. For TDS ($K=64$), increasing the twist scale generally decreases the human-rated accuracy.  For remaining methods, a moderate increase in twist scale improves the human-rated accuracy; however, excessively large twist scale can hurt the accuracy.}
    \label{subfig:cg-twist-scale-human}    
\end{subfigure}
    \caption{MNIST class-conditional generation: classification accuracy under different twist scales, computed by a neural network classifier (top panel) and a human (bottom panel).}
    \label{fig:cg-twist-scale}
\end{figure}


\begin{figure}[tp]
\begin{subfigure}[t]{0.95\linewidth}
    \centering
    \includegraphics{sections/figures/mnist-class/twist-scale-comparision-tsmc-label6-P64.pdf}
    \caption{TDS ($K=64$). \# of good-quality samples counted by human is 62, 61, 63, 60, 62, from left to right.}
    % number of artifacts 2, 3, 1, 4, 2
    \label{subfig:twist-scale-sample-TDS-K64}
\end{subfigure}
\begin{subfigure}[t]{0.95\linewidth}
    \centering
    \includegraphics{sections/figures/mnist-class/twist-scale-comparision-tsmc-label6-P2.pdf}
    \caption{TDS ($K=2$). \# of good-quality samples counted by human 49, 51, 51,  46, 45, from left to right.}
    % number of artifacts 15, 13, 13, 18, 19
    \label{subfig:twist-scale-sample-TDS-K2}
\end{subfigure}
\begin{subfigure}[t]{0.95\linewidth}
    \centering
    \includegraphics{sections/figures/mnist-class/twist-scale-comparision-tsis-label6-P64.pdf}
    \caption{TDS-IS ($K=64$). \# of good-quality samples counted by a human is 62, 61, 64, 64, 61, from left to right.}
    % number of artifacts 2, 3, 0, 0, 3
    \label{subfig:twist-scale-sample-TDS-IS-K64}
\end{subfigure}
\begin{subfigure}[t]{0.95\linewidth}
    \centering
    \includegraphics{sections/figures/mnist-class/twist-scale-comparision-tsis-label6-P2.pdf}
    \caption{TDS-IS ($K=2$). \# of good-quality samples counted by a human is 49, 52, 53, 52, 48, from left to right.}
    % number of artifacts   15, 12, 11, 12, 16 
    \label{subfig:twist-scale-sample-TDS-IS-K2}
\end{subfigure}
\begin{subfigure}[t]{0.95\linewidth}
    \centering
    \includegraphics{sections/figures/mnist-class/twist-scale-comparision-guidance-label6-P2.pdf}
    \caption{Gradient Guidance. \# of good-quality samples counted by a human is 31, 38, 41, 39, 34, from left to right.}
    % number of artifacts 33,26,23,25,30
    \label{subfig:twist-scale-sample-guidance}
\end{subfigure}

\caption{MNIST class-conditional generation: random samples selected from 64 random runs conditioned on class 6 under different twist scales. Top to bottom: TDS ($K=64$), TDS ($K=2$), TDS-IS ($K=64$), TDS-IS ($K=2$) and Gradient Guidance. In general, moderately large twist scales improve the sample quality. However, overly large twist scale (e.g. 10) would distort the digit shape with more artifacts, though retaining useful features that may allow a neural network classifier to identity the class.}
\label{fig:cg-twist-scale-samples}
\end{figure}


\paragraph{Ablation study on twist scales.} We consider exponentiating and re-normalizing twisting functions by a \emph{twist scale} $\gamma$, i.e. setting new twisting functions to $\pmodelapprox (y \mid \xt; \gamma) \propto \plikelihood{y}{\denoiser(\xt)}^\gamma$. In particular, when $t=0$, we set $\plikelihood{y}{ \xzero;\gamma}   \propto \plikelihood{y }{\xzero}^\gamma$. This modification suggests that the targeted conditional distribution is now  
\begin{align}
    \pmodel (\xzero \mid y; \gamma) \propto \pmodel( \xzero ) \plikelihood{y}{\xzero}^\gamma . 
\end{align}
By setting $\gamma > 1$, the classification likelihood becomes sharper, which is potentially helpful for twisting the samples towards a specific class. 
The TDS algorithm (and likewise TDS-IS and Gradient Guidance) still apply with this new definition of twisting functions.
The use of twist scale is similar to the \emph{guidance scale} introduced in the literature of Classifier Guidance, which is used to multiply the gradient of the log classification probability \citep{dhariwal2021diffusion}. 


In \Cref{fig:cg-twist-scale}, we examine the effect of varying twist scales on classification accuracy of TDS, TDS-IS and Gradient Guidance. We consider two ways to evaluate the accuracy. First, \emph{classification accuracy} computed by a \emph{neural network classifier}, where the evaluation setup is the same as in \Cref{subsec:exp-mnist}. Second, the \emph{human-rated} classification accuracy, where a human (one of the authors) checks if a generated digit has the right class and does not have artifacts. 
Since human evaluation is expensive, we only evaluate TDS, TDS-IS (both with $K= 64, 2$) and Gradient Guidance. In each  run, we randomly sample one particle out of $K$ particles according to the associated weights. We conduct 64 runs for each class label, leading to a total of 640 samples for human evaluation.


\Cref{subfig:cg-twist-scale-classifier} depicts the classification accuracy measured by a neural network classifier. We observe that in general larger twist scale improves the classification accuracy. For TDS and TDS-IS, the improvement is more significant for smaller number of particles $K$ used. 

\Cref{subfig:cg-twist-scale-human} depicts the human-rated accuracy. In this case, we find that larger twist scale is not necessarily better. A moderately large twist scale ($\gamma=2,3$) generally helps increasing the accuracy, while an excessively large twist scale ($\gamma =10$) decreases the accuracy. An exception is TDS with $K=64$ particles, where any $\gamma >1$ leads to worse accuracy compared to the case of $\gamma=1$.  Study on twist scale aside, we find that using more particles $K$ help improving human-rated accuracy (recall that Gradient Guidance is a special case of TDS with $K=1$): given the same twist scale, TDS with $K=1,2$ or 64 has increasing accuracy.  In addition, both TDS and TDS-IS with $K=64$ and $\gamma=1$ have almost perfect accuracy. The effect of $K$ on human-rated accuracy is consistent with previous findings with neural network classifier evaluation in \Cref{subsec:exp-mnist}. 

We note that there is a discrepancy on the effects of twist scales between neural network evaluation and human evaluation. We suspect that when overly large twist scale is used, the generated samples may fall out of the data manifold; however, they may still retain features recognizable to a neural network classifier, thereby leading to a low human-rated accuracy but a high classifier-rated accuracy. To validate this hypothesis, 
we present samples conditioned on class 6 in \Cref{fig:cg-twist-scale-samples}. For example, in \Cref{subfig:twist-scale-sample-guidance},  Gradient Guidance with $\gamma=1$ has 31 good-quality samples out of 64, and the rest of the samples often resamble the shape of other digits, e.g. 3,4,8; and Gradient Guidance with $\gamma=10$ has 34 good-quality samples, but most of the remaining samples resemble 6 with many artifacts. 

\paragraph{Effective sample size.} Effective sample size (ESS) is a common
metric used to diagnose the performance of SMC samplers, which is defined as 
 $ (\sum_{k=1}^K \dt{w}{t}_k)^2 / (\sum_{k=1}^K (\dt{w}{t}_k)^2 )$ for $K$ weighted particles  $\{\xt_k, \dt{w}{t}_k\}_{k=1}^K$. 
Note that ESS is always bounded between $0$ and $K$.

\Cref{fig:mnist-class-ess} depicts the ESS trace comparison of TDS, TDS-IS, and IS. TDS has a general upward trend of ESS approaching $K = 64$. Though in the final few steps ESS of TDS drops by a half, it is still higher than that of TDS-IS and IS that deteriorates to around 1 and 6 respectively. In practice, one can use the TDS-truncate variant introduced above to avoid the ESS drop in the final few steps. 

\begin{figure}
\centering
 %   \begin{subfigure}[t]{0.46\linewidth}
  %  \centering
    \includegraphics[width=0.4\textwidth]{sections/figures/mnist-class/class-ess-legendTrue.pdf}
      \caption{ESS trace avg. over 100 runs ($K=64$)}%for SMC-based methods averaged over 100 runs.}
      \label{fig:mnist-class-ess}
  %\end{subfigure}%\hfill 
  %\caption{Caption}
  %  \label{fig:class-cond-ess-all}
\end{figure}


\subsubsection{Image inpainting on MNIST}
\label{supp:subsubsec:mnist-inpainting}
The inpainting task is to sample images from $\pmodel(\xzero \mid \xzero_\mask=y)$ given observed part $\xzero_\mask = y$. Here we segment $\xzero=[\xzero_\mask, \xzero_\umask]$ into observed dimensions $\mask$ and unobserved dimensions $\umask$, 
as described in  \Cref{subsec:conditioning-problems}. 

In this experiment, we consider two types of observed dimensions: (1) $\mask$ = ``half'', where the left half of an image is observed, and (2) $\mask$ = ``quarter'', where the upper left quarter of an image is observed. 

We run TDS, TDS-IS, Gradient Guidance,  SMC-Diff, and Replacement method to inpaint 10{,}000 validation images. We also include TDS-truncate that truncates the TDS procedure at $t = 10$ and returns $\denoiser(\dt{x}{10})$. 
 
%Similar to the class-conditional generation task, we can use a twist scale to exponentiate the twisting functions. In the inpainting case, it is equivalent to changing the variance of the twisting function   in \cref{eqn:inpaint_twist} by $\gamma_t$. Denote the new variance by $\tilde \gamma_t$. By default TDS sets $\tilde \gamma_t := (\tilde \sigma_t^2 \tau^2)/(\tilde \sigma_t^2 + \tau^2)$, where $\tilde \sigma_t^2:=(1-\bar \alpha_t)/\bar \alpha_t$ and $\tau^2=0.12$ is the population variance of $\xzero$ estimated from training data. 


\paragraph{Twisting function variance schedule.} We use a flexible variance schedule $\{\hat \sigma_t^2\}$ to extend the twisting functions in \cref{eqn:inpaint_twist} to  
\begin{align}
 \pTwist (y \mid \xt, \mask) &\coloneqq \mathcal{N}(y \mid \denoiser(\xt)_\mask, \hat \sigma_t^2). 
\end{align}
Ideally, $\hat \sigma_t^2$ should match $\Var_{\pmodel} [y \mid \xt, \mask]$. In TDS, we choose $\hat \sigma_t^2 = \frac{ (\bar \sigma_t^2 / \bar \alpha_t ) \cdot \tau^2}{ \bar \sigma_t^2 / \bar \alpha_t + \tau^2}$ where $\tau^2 = 0.12$  is the estimated sample variance of the training data (averaged over pixels). This choice is inspired from \citet[Appendix A.3]{song2023pseudoinverse}: if the data distribution is $q(\xzero) = \mathcal{N}(0, \tau^2)$, and $q(\xt \mid \xzero) = \mathcal{N}(\xt \mid \sqrt{\bar \alpha_t} \xzero, \sigma_t^2)$, then by Bayes rule $\Var_q[\xzero \mid \xt] =  \frac{ (\bar \sigma_t^2 / \bar \alpha_t ) \cdot \tau^2}{ (\bar \sigma_t^2 / \bar \alpha_t) + \tau^2}$. 

Changing $\hat \sigma_t^2$ has a similar effect as using a twist scale to exponentiate the twisting functions in \Cref{supp:subsubsec:mnist-class}. In particular, for a Gaussian distribution, exponentiating the distribution corresponds to re-scaling the variance: $\mathcal{N}(x\mid \mu, (\sigma/\sqrt{\gamma})^2)  \propto \mathcal{N}(x\mid \mu, \sigma^2)^{\gamma}$.




%By default TDS sets $\tilde \gamma_t := (\tilde \sigma_t^2 \tau^2)/(\tilde \sigma_t^2 + \tau^2)$, where $\tilde \sigma_t^2:=(1-\bar \alpha_t)/\bar \alpha_t$ and $\tau^2=0.12$ is the population variance of $\xzero$ estimated from training data. 




\paragraph{Metrics.} We consider 3 metrics. (1) We use the effective sample size (ESS) to compare the particle efficiency among different SMC samplers (namely TDS, TDS-IS, and SMC-Diff).

(2) In addition, we ground the performance of 
a sampler in the downstream task of classifying a partially observed image $\xzero_\mask = y$: we use a classifier to predict the class $\hat z$ of an inpainted image $\xzero$, and compute the accuracy against the true class $z^*$ of the unmasked image ($z^*$ is provided in MNIST dataset). This prediction is made by the same classifier used in \Cref{subsec:exp-mnist}. 

Consider $K$ weighted particles $\{ \xzero_k; \dt{w}{0}_k\}_{k=1}^K$ drawn from a sampler conditioned on $\xzero_\mask = y$ and assume the weights are normalized. 
%Notice that $ \sum_{k=1}^K \dt{w}{0}_k\delta_{\xzero_k} (\xzero) \approx \pmodel(\xzero_0\mid \xzero_\mask = y)$. This observation motivates us to define a metric, 
We define the \emph{Bayes accuracy} (BA) as 
\begin{align}\label{eq:BA}
    \mathds{1} \left\{ \hat z(y)= z^* \right\}, \quad  
   \mathrm{ with } \quad \hat z (y): = \argmax_{z=1:10} \sum_{k=1}^K \dt{w}{0}_k p(z;\xzero_k),  
\end{align}
where $\hat z(y)$ is viewed as an approximation to the Bayes optimal classifier $\hat z^* (y)$ given by 
\begin{align}\label{eq:bayes_opt_classifier}
\begin{split}
    \hat z^*(y) & := \argmax_{z=1:10} \pmodel(z \mid \xzero_\mask = y) = \argmax_{z=1:10} \int \pmodel(\xzero \mid \xzero_\mask = y) p(z;  \xzero )d\xzero.
\end{split}
\end{align}
(In \cref{eq:bayes_opt_classifier} we assume the classifier $\plikelihood{\cdot}{\xzero}$ is the optimal classifier on full images.)

(3) We also consider \emph{classification accuracy} (CA) defined as the following 
\begin{align}\label{eq:CA}
    \sum_{k=1}^K \dt{w}{0}_k \mathds{1} \{ \hat z(\xzero_k) = z^*\}, \quad \mathrm{ with } \quad \hat z(\xzero_k) := \argmax_{z=1:10} \plikelihood{z}{\xzero_k}. 
\end{align}

 BA and CA evaluate different aspects of a sampler. BA is focused on the optimal prediction among multiple particles, whereas CA is focused on the their weighted average prediction. 

%Note that Guidance and Replacement with $K$ independent samples can be viewed as $K$ particles with uniform weights. Their CA is independent of $K$ in expectation which we use $K=64$ to estimate, and their BA is usually dependent on $K$. 

\begin{figure}[t]
\begin{subfigure}[t]{0.27\linewidth}
    \centering
    \includegraphics[width=1.0\linewidth]{sections/figures/mnist-inpainting/inpainting-half-ess-includeISFalse-includet0False.pdf}
    \caption{$\mask$ = ``half''. ESS trace.}
    \label{subfig:mnist-inpainting-ess-half}
\end{subfigure}\hfill
\begin{subfigure}[t]{0.7\linewidth}
    \centering
    \includegraphics[width=1.0\linewidth]{sections/figures/mnist-inpainting/inpainting-half-accuracy-includeISFalse.pdf}
    \caption{$\mask$ = ``half''. Class. accuracy (left), and Bayes accuracy (right) v.s. $K$.}
    \label{subfig:mnist-inpainting-accuracy-half} 
    \end{subfigure}\hfill 
\begin{subfigure}[t]{0.27\linewidth}
    \centering
    \includegraphics[width=1.0\linewidth]{sections/figures/mnist-inpainting/inpainting-quarter-ess-includeISFalse-includet0False.pdf}
     \caption{$\mask$ = ``quarter''. ESS trace.}
     \label{subfig:mnist-inpainting-ess-quarter}
\end{subfigure}\hfill
\begin{subfigure}[t]{0.7\linewidth}
    \centering
    \includegraphics[width=1.0\linewidth]{sections/figures/mnist-inpainting/inpainting-quarter-accuracy-includeISFalse.pdf}
    \caption{$\mask$ = ``quarter''. Class. accuracy (left), and Bayes accuracy (right) v.s.  $K$.}
    \label{subfig:mnist-inpainting-accuracy-quarter} 
    \end{subfigure}
\caption{MNIST image inpainting. Results for observed dimension $\mask$ = ``half'' are shown in the top panel, and $\mask$ = ``quarter'' in the bottom panel.  
(i) Left column: ESS traces are averaged over 100 different images inpainted by TDS, TDS-IS, and SMC-Diff, all with $K=64$ particles. 
 TDS's ESS is generally increasing untill the final 20 steps where it drops to around 1, suggesting significant particle collapse in the end of generation trajectory. TDS-IS's ESS is always around 1. SMC-Diff has higher particle efficiency compared to TDS. 
(ii) Right column: Classification accuracy and Bayes accuracy are averaged over 10{,}000 images.
In general increasing the number of particles $K$ would improve the performance of all samplers. TDS and TDS-truncate have the highest accuracy among all given the same $K$. 
(iii) Finally, comparison of top and bottom panels shows that in a harder inpainting problem where $\mask$ = ``quarter'', TDS's has a higher ESS but lower CA and BA. }
\label{fig:mnist-inpaiting-ess-and-class-accuracy}
\end{figure}

\paragraph{Comparison results of different methods.}
\Cref{fig:mnist-inpaiting-ess-and-class-accuracy} depicts the ESS trace, BA and CA for different samplers. 
The overall observations are similar to the observations in the class-conditional generation task in \Cref{subsec:exp-mnist}, except that SMC-Diff and Replacement methods are not available there.  

Replacement has the lowest CA and BA across all settings. Comparing TDS to SMC-Diff, we find that SMC-Diff's ESS is consistently greater than TDS; however, SMC-Diff is outperformed by TDS in terms of both CA and BA. 

We also note that despite Gradient Guidance's CA is lower,  its BA is comparable to TDS. This result is due to that as long as Gradient Guidance generates a few good-quality samples out of $K$ particles, the optimal prediction can be accurate, thereby resulting in a high BA.  


\begin{figure}[t]
\begin{subfigure}[t]{\linewidth}
        \centering
    \includegraphics[width=0.9\linewidth]{sections/figures/mnist-inpainting/comparison_inpainting_half_pred_accuracy.pdf}
    \caption{$\mask$ = ``half''. Classification accuracy under different twist scale schemes.}
    \label{subfig:mnist-inpainting-twist-scale-mask-half}
\end{subfigure}
\begin{subfigure}[t]{\linewidth}
        \centering
    \includegraphics[width=0.9\linewidth]{sections/figures/mnist-inpainting/comparison_inpainting_quarter_pred_accuracy.pdf}
    \caption{$\mask$ = ``quarter''. Classification accuracy under different twist scale schemes.}
    \label{subfig:mnist-inpainting-twist-scale-mask-quarter}
\end{subfigure}
\caption{MNIST image inpainting: Ablation study on twist scales. Top: observed dimensions $\mask$ = ``half''. Bottom: $\mask$ = ``quarter''. In most case, the performance of our choice of twist scale is similar to that of $\Pi$GDM, and is better compared to DPS.}
\label{fig:mnist-inpainting-twist-scale}
\end{figure}

\paragraph{Ablation study on twisting function variance schedule.}
We compare the following three options: 
\begin{enumerate}
    \item DPS: $\hat \sigma_t:= 2\| \denoiser(\xt)_{\mask} - y\|_2 \sigma_t^2$, adapted from the DPS method \citep[Algorithm 1]{chung2022diffusion}, 
    \item $\Pi$GDM: $\hat \sigma_t:=\sigma_t^2  / \sqrt{\bar \alpha_t}$, adapted from  the $\Pi$GDM method  \citep[Algorithm 1]{song2023pseudoinverse},  
    \item TDS: $\hat \sigma_t^2 = \frac{ (\bar \sigma_t^2 / \bar \alpha_t ) \cdot \tau^2}{ (\bar \sigma_t^2 / \bar \alpha_t) + \tau^2}$ where $\tau^2 = 0.12$. 
\end{enumerate}

\Cref{fig:mnist-inpainting-twist-scale} shows the classification accuracy of TDS, TDS-IS and Gradient Guidance with different twist scale schemes.  We find that our choice  has similar performance to that of $\Pi$GDM, and outperforms DPS in most cases. Exceptions are when $\mask$ = ``quarter'' and for large $K$, TDS with twist scale choice of $\Pi$GDM or DPS has higher CA, as is shown in the left panel in  \Cref{subfig:mnist-inpainting-twist-scale-mask-quarter}.


\subsubsection{Class-conditional generation on ImageNet and CIFAR-10}
\label{supp:subsubec:inet-and-cifar}


\paragraph{ImageNet.} We compare TDS (with \# of particles $K=1,16$) to Classifier Guidance (CG) using the same unconditional model from \url{https://github.com/openai/guided-diffusion/tree/main}. CG uses a classifier trained on noisy inputs taken from this repository as well. For TDS, we use the same classifier evaluated at timestep = 0 to mimic a standard trained classifier. We generate 16 images for each of the 1000 class labels, using 100 sampling steps. TDS uses a twist scale of 10, and CG uses a guidance scale of 10. Notably, given a fixed class, TDS ($K=16$) generates correlated samples in a single SMC run, and TDS ($K=1$) and CG generate 16 independent samples. 


In \Cref{fig:inet1}, we observe that TDS can faithfully capture the class and have comparable image quality to CG’s , although with less diversity than CG and TDS ($K=1$). 

\Cref{fig:inet-comparison} shows more samples given randomly selected classes. We also reported results of the unconditional model from the original paper \citep{dhariwal2021diffusion} that are evaluated on 50k samples with 250 sampling steps in \Cref{tab:inet-comparison}. TDS and CG provide similar classification accuracy. TDS has similar FIDs compared to the unconditional model and better inception score. CG’s FID and inception score are better than TDS. We suspect this difference is attributed to the sample correlation (and hence less diversity) within particles in a single run of TDS ($K=16$).


\begin{figure}
    \centering
    
    \begin{subfigure}{0.3\textwidth}
        \includegraphics[width=\linewidth]{\detokenize{sections/figures/inet-and-cifar10/inet-samples/tds-samples-10-brambling__Fringilla_montifringilla.png}}
        \caption{TDS (P=16)}
    \end{subfigure}
    \begin{subfigure}{0.3\textwidth}
        \includegraphics[width=\linewidth]{\detokenize{sections/figures/inet-and-cifar10/inet-samples/tdsg-samples-10-brambling__Fringilla_montifringilla.png}}
        \caption{TDS (P=1)}
    \end{subfigure}
    \begin{subfigure}{0.3\textwidth}
        \includegraphics[width=\linewidth]{\detokenize{sections/figures/inet-and-cifar10/inet-samples/gd-samples-10-brambling__Fringilla_montifringilla.png}}
        \caption{Classifier Guidance}
    \end{subfigure}

    \caption{256$\times$256 ImageNet. Samples from TDS and Classifier Guidance given class `brambling'.}
    \label{fig:inet1}
\end{figure}

\begin{figure}
    \centering
    \includegraphics[scale=0.1]{sections/figures/inet-and-cifar10/inet-samples/samples-989-212-883-869-15-558-276-71-544-613-690-443-731-130-763.png}
    \caption{256$\times$256 ImageNet. 15 random class samples. Top to bottom: TDS(P=16), TDS(P=1), Classifier Guidance.}
    \label{fig:inet-comparison}
\end{figure}

\begin{table}[!t]
\centering 
\begin{tabular}{|l|l|l|l|}
\hline
Method               & Classification Accuracy $\uparrow$ & FID $\downarrow$   & Inception Score $\uparrow$    \\ \hline
TDS ($K=16$)           & 99.90\%                 & 26.65 & 64.03  \\ \hline
TDS ($K=1$)            &  99.63\%                & 26.05 & 44.45       \\ \hline
Classifier Guidance  & 99.17\%                 & 14.03 & 100.33 \\ \hline
Unconditional model             & n/a                     & 26.21 & 39.70  \\ \hline
%Unocnd. w. Guidance* & n/a                     & 12.00 & 95.41  \\ \hline
%Cond.*               & n/a                     & 10.94 & 100.98 \\ \hline
\end{tabular}
\label{tab:inet-comparison}
\caption{ImageNet. Comparison of sample quality. The top three methods are evaluated on 16k samples generated with 100 steps. The unconditional model performance is provided in \citep{dhariwal2021diffusion}, which is evaluated on 50k samples generated with 250 steps.}
\end{table}



\paragraph{CIFAR-10.} We ran TDS ($K=16$) with the twist scale = 1 and 100 sampling steps using diffusion model from \url{https://github.com/openai/improved-diffusion} and classifier from \url{https://github.com/VSehwag/minimal-diffusion/tree/main}. TDS generates faithful and diverse images, see \Cref{fig:cifar}. However, we found TDS can occasionally generate incorrect samples for the class ‘truck’ (the last panel in \Cref{fig:cifar}).



\begin{figure}[!t]
    \centering
    
    \begin{subfigure}{0.15\textwidth}
        \includegraphics[width=\linewidth]{\detokenize{sections/figures/inet-and-cifar10/cifar10-samples/sample_label0_iter260.png}}
        \caption{airplane}
    \end{subfigure}
    \begin{subfigure}{0.15\textwidth}
        \includegraphics[width=\linewidth]{\detokenize{sections/figures/inet-and-cifar10/cifar10-samples/sample_label4_iter324.png}}
        \caption{deer}
    \end{subfigure}
    \begin{subfigure}{0.15\textwidth}
        \includegraphics[width=\linewidth]{\detokenize{sections/figures/inet-and-cifar10/cifar10-samples/sample_label7_iter327.png}}
        \caption{horse}
    \end{subfigure}
    \begin{subfigure}{0.15\textwidth}
        \includegraphics[width=\linewidth]{\detokenize{sections/figures/inet-and-cifar10/cifar10-samples/sample_label8_iter448.png}}
        \caption{ship}
    \end{subfigure}
    \begin{subfigure}{0.15\textwidth}
        \includegraphics[width=\linewidth]{\detokenize{sections/figures/inet-and-cifar10/cifar10-samples/sample_label9_iter219.png}}
        \caption{truck}
    \end{subfigure}
    \begin{subfigure}{0.15\textwidth}
        \includegraphics[width=\linewidth]{\detokenize{sections/figures/inet-and-cifar10/cifar10-samples/sample_label9_iter439.png}}
        %\label{subfig:truck-failure-mode}
        \caption{failure mode}
    \end{subfigure}

    \caption{CIFAR10. Samples from TDS with $K=16$  particles given different class labels (within a single SMC run). In most cases TDS generate authentic and diverse images. The last panel presents a failure mode for the class `truck'.}
    \label{fig:cifar}
\end{figure}

Finally, we present the ESS trace plot of the three image datasets in \Cref{fig:ess-all}. We see there is a sudden ESS drop in the final stage for MINST, sometimes for CIFAR10, but usually not on ImageNet. Mechanically, the drop in ESS implies a large discrepancy between the final and intermediate target distributions. We suspect such discrepancies might arise from irregularities in the denoising network near $t=0$. In practice, one can consider early truncating the TDS procedure to prevent such ESS drops and promote particle diversity, as is studied in \Cref{supp:subsubsec:mnist-class,supp:subsubsec:mnist-inpainting}. 



\begin{figure}[!t]
        \centering
        \includegraphics[width=0.9\linewidth]{sections/figures/inet-and-cifar10/ESS_mnist_cifar_inet.pdf}
        \caption{2 example effective sample size traces of TDS ($K=16$, 100 sampling steps) on MNIST, CIFAR-10 and ImageNet models, respectively.}
        \label{fig:ess-all}
\end{figure}




\section{Motif-scaffolding application details}\label{sec:supp_results_protein}


\textbf{Unconditional model of protein backbones.}
We here use the FrameDiff, a diffusion generative model described by \citet{yim2023se}.
FrameDiff parameterizes protein $N$-residue protein backbones as a collection of rigid bodies defined by rotations and translations as $\xzero=[(R_1, z_1), \dots, (R_N, z_N)]\in SE(3)^N$.
$SE(3)$ is the special Euclidean group in three dimensions (a Riemannian manifold). 
Each $R_n\in SO(3)$ (the special orthogonal group in three dimensions) is a rotation matrix and each $z_n\in \R^3$ in a translation.
Together, $R_n$ and $z_n$ describe how one obtains the coordinates of the three backbone atoms \texttt{C}, $\texttt{C}_\alpha$, and \texttt{N} for each residue by translating and rotating the coordinates of an \emph{idealized} residue with $C_\alpha$ carbon and the origin. 
The conditioning information is then a motif $y=\xzero_\mask\in SE(3)^{|\mask|}$ for some $\mask\subset \{1,\dots, N\}.$
FrameDiff is a continuous time diffusion model and includes the number of steps as a hyperparmeter; we use 200 steps in all experiments.
We refer the reader to \citep{yim2023se} for details on the neural network architecture,
and details of the forward and reverse diffusion processes.

\paragraph{Twisting functions for motif scaffolding.}
To define a twisting function, we use a tangent normal approximation to $\pmodel(y\mid \xt)$ as introduced in \cref{eqn:tangent_normal}.
In this case the tangent normal factorizes across each residue and across the rotational and translational components.
The translational components are represented in $\R^3,$ and so are treated as in the previous Euclidean cases, using the variance preserving extension described in \Cref{sec:supp_method}.
For the rotational components, represented in the special orthogonal group $SO(3),$ the tangent normal may be computed as described in \Cref{sec:riemannian}, using the known exponential map on $SO(3)$.
In particular, we use as the twisting function
\begin{align}
\twist (y\mid \xt, \mask) = \prod_{m=1}^{|\mask |} \mathcal{N}(y^T_m; \denoiser(\xt)^T_{\mask_m}, 1-\bar \alpha_t)\mathcal{TN}_{\denoiser(\xt)^R_{\mask_m}}(y_m^R; 0, \bar \sigma_t^2),
\end{align}
where $y^T_m, \denoiser(\xt)^T_{\mask_m} \in \R^3$ represent the translations associated with the $m^{th}$ residue of the motif and its prediction from $\xt,$ and
$y^R_m, \denoiser(\xt)^R_{\mask_m} \in SO(3)$ represent the analogous quantities for the rotational component.
Next, $1-\bar \alpha_t=\Var_q[\xt \mid \xzero]$ and  $\bar \sigma_t^2$ is the time of the Brownian motion associated with the forward process at step $t$ \citep{yim2023se}.
For further simplicity, our implementation further approximates log density of the tangent normal on rotations using the squared Frobenius norm of the difference between the rotations associated with the motif and the denoising prediction (as in \citep{watson2022broadly}), which becomes exact in the limit that $\bar \sigma_t$ approaches 0 but avoids computation of the inverse exponential map and its Jacobian.





%Though $\nabla_{\xt}\log \twist(y; \xt)$ is a $\mathcal{T}_{\xt}$-valued Riemannian gradient, it is nevertheless computable by automatic differentiation.


%\BLT{Could postpone for arxiv version.} Significantly higher success-rates are obtained when using ESMfold rather than AF2, and when not constraining the sequence (see appendix).


\textbf{Motif rotation and translation degrees of freedom.}  We similarly seek to eliminate the rotation and translation of the motif as a degree of freedom. 
We again represent our ambivalence about the rotation of the motif with randomness,
augment our joint model to include the rotation with a uniform prior, 
and write
\begin{align}
\plikelihood{y}{\xzero}=\int p(R)\plikelihood{y}{\xzero, R}\quad  \mathrm{for}\quad \plikelihood{y}{\xzero, R}=\delta_{Ry}(\xzero)
\end{align}
and with $p(R)$ the uniform density (Haar measure) on SO(3), and $Ry$ represents rotating the rigid bodies described by $y$ by $R.$ 
For computational tractability, we approximate this integral with a Monte Carlo approximation defined by subsampling a finite number of rotations, $\mathcal{R} = \{R_k\}$ (with $|\mathcal{R}|=$\texttt{\# Motif Rots.} in total).
Altogether we have
$$
\twist(y\mid \xt) = |\mathcal{R}|\inv \cdot |\mathcal{M}|\inv \sum_{R \in \mathcal{R}} \sum_{\mask \in \mathcal{M}} \twist(Ry\mid \xt, \mask).
$$

The subsampling above introduces the number of motif locations and rotations subsampled as a hyperparameter.
Notably the conditional training approach does not immediately address these degrees of freedom,
and so prior work randomly sampled a single value for these variables \citep{trippe2022diffusion,wang2022scaffolding,watson2022broadly}.

\paragraph{Evaluation details.}
We use \emph{self-consistency} evaluation approach for generated backbones \citep{trippe2022diffusion} that 
(i) uses fixed backbone sequence design (inverse folding) to generate a putative amino acid sequence to encode the backbone, 
(ii) forward folds sequences to obtain backbone structure predictions, and
(iii) judges the quality of initial backbones by their agreement (or self-consistency) with predicted structures.
We inherit the specifics of our evaluation and success criteria set-up following \citep{watson2022broadly}, including
using ProteinMPNN \citep{dauparas2022robust} for step (i) and AlphaFold \citep{jumper2021highly} on a single sequence (no multiple sequence alignment) for (ii).

In this evaluation we use ProteinMPNN \citep{dauparas2022robust} with default settings to generate 8 sequences for each sampled backbone.
Positions not indicated as resdesignable in \citep[Methods Table 9]{watson2022broadly} are held fixed.
We use AlphaFold \citep{jumper2021highly} for forward folding.
We define a ``success'' as a generated backbone for which at least one of the 8 sequences has backbone atom with both 
scRMSD < 1 {\AA} on the motif and 
scRMSD < 2 {\AA} on the full backbone.

We benchmarked TDS on 24/25 problems in the benchmark set introduced \citep[Methods Table 9]{watson2022broadly}.
A 25th problem (6VW1) is excluded because it involves multiple chains, which cannot be represented by FrameDiff. FrameDiff requires specifying a total length of scaffolds.
In all replicates, we fixed the scaffold the median of the \texttt{Total Length} range specified by \citet[Methods Table 9]{watson2022broadly}.
For example, 116 becomes 125 and 62-83 becomes 75.


As discussed in the main text, in our handling of the motif-placement degree of freedom we restrict the number of possible placements considered for discontiguous motifs.
To select these placements we
(i) restrict to masks which place the motif indices in a pre-specified order that we do not permute
and do not separate residues that appear contiguously in source PDB file,
and 
(ii) when there are still too many possible placements 
sub-sample randomly to obtain the set of masks $\mathcal{M}$ of at most some maximum size (\texttt{\# Motif Locs.}).
We do not enforce the spacing between segments specified in the ``contig'' description described by \citep{watson2022broadly}.


\paragraph{Resampling.}
In all motif-scaffolding experiments, we use systematic resampling.
An effective sample size threshold is used $K/2$ in all cases, that is, we trigger resampling steps only when the effective sample size falls below this level.


\subsection{Additional Results}

\begin{figure}[t]
  \centering
    \includegraphics[width=1.0\textwidth]{sections/figures/protein/benchmark_success_rates}
      \caption{Results on full motif-scaffolding benchmark set.  The y-axis is the fraction of successes across at least $200$ replicates.  Error bars are $\pm 2$ standard errors of the mean.  Problems with scaffolds shorter than 100 residues are bolded.  TDS (K=8) outperforms the state-of-the-art (RFdiffusion) on most problems with short scaffolds.}
  \label{fig:full_benchmark}
\end{figure}

\paragraph{Impact of twist scale on additional motifs.}
\Cref{fig:5IUS_results} showed monotonically increasing success rates with the twist scale.
However, this trend does not hold for every problem.
\Cref{{fig:prot-twist-scale}} demonstrates this by comparing the success rates with different twist-scales on five additional benchmark problems.

\begin{figure}[t]
    \tiny
    \includegraphics[scale=1.0]{sections/figures/protein/twist_scale_figure.png}
  \centering
    \caption{Increasing the twist scale has a different impact across motif-scaffolding benchmark problems.}
\label{fig:prot-twist-scale}
\end{figure}

\paragraph{Effective sample size varies by problem.}
\Cref{fig:prot-ESS} shows two example effective sample size traces over the course of sample generation.
For \texttt{6EXZ-med} resampling was triggered 38 times (with 14 in the final 25 steps),
and for \texttt{5UIS} resampling was triggered 63 times (with 13 in the final 25 steps).
The traces are representative of the larger benchmark.

\begin{figure}[t]
    \tiny
    \includegraphics[scale=1.0]{sections/figures/protein/6EXZ_med_38_resamples.png}
    \includegraphics[scale=1.0]{sections/figures/protein/5UIS_63_resample.png}
  \centering
    \caption{
    Effective sample size traces for two motif-scaffolding examples (Left) \texttt{6EXZ-med} and (Right) \texttt{5IUS}.
    In both case $K=8$ and a resampling threshold of $0.5K$ is used.  Dashed red lines indicate resampling times. 
}
\label{fig:prot-ESS}
\end{figure}


%\paragraph{Impact of not constraining sequence, using ESMFold instead of AF2}
%\BLT{This can be left for the arxiv and skipped for NeurIPS appendix.}


\paragraph{Application of TDS to RFdiffusion:}
We also tried applying TDS to RFdiffusion \cite{watson2022broadly}.
RFdiffusion is a diffusion model that uses the same backbone structure representation as FrameDiff.
However, unlike FrameDiff, RFdiffusion trained with a mixture of conditional examples, in which a segment of the backbone is presented as input as a desired motif,
and unconditional training examples in which the only a noised backbone is provided.
Though in our benchmark evaluation provided the motif information as explicit conditioning inputs, 
we reasoned that TDS should apply to RFdiffusion as well (with conditioning information not provided as input).
However, we were unable to compute numerically stable gradients (with respect to either the rotational or translational components of the backbone representation); 
this lead to twisted proposal distributions that were similarly unstable and trajectories that frequently diverged even with one particle.
We suspect this instability owes to RFdiffusion's structure prediction pretraining and limited fine-tuning, which may allow it to achieve good performance without having fit a smooth score-approximation.






\end{document}
